% Options for packages loaded elsewhere
\PassOptionsToPackage{unicode}{hyperref}
\PassOptionsToPackage{hyphens}{url}
%
\documentclass[
  11pt,
  a4paper,
]{scrbook}

\usepackage{amsmath,amssymb}
\usepackage{iftex}
\ifPDFTeX
  \usepackage[T1]{fontenc}
  \usepackage[utf8]{inputenc}
  \usepackage{textcomp} % provide euro and other symbols
\else % if luatex or xetex
  \usepackage{unicode-math}
  \defaultfontfeatures{Scale=MatchLowercase}
  \defaultfontfeatures[\rmfamily]{Ligatures=TeX,Scale=1}
\fi
\usepackage{lmodern}
\ifPDFTeX\else  
    % xetex/luatex font selection
\fi
% Use upquote if available, for straight quotes in verbatim environments
\IfFileExists{upquote.sty}{\usepackage{upquote}}{}
\IfFileExists{microtype.sty}{% use microtype if available
  \usepackage[]{microtype}
  \UseMicrotypeSet[protrusion]{basicmath} % disable protrusion for tt fonts
}{}
\makeatletter
\@ifundefined{KOMAClassName}{% if non-KOMA class
  \IfFileExists{parskip.sty}{%
    \usepackage{parskip}
  }{% else
    \setlength{\parindent}{0pt}
    \setlength{\parskip}{6pt plus 2pt minus 1pt}}
}{% if KOMA class
  \KOMAoptions{parskip=half}}
\makeatother
\usepackage{xcolor}
\usepackage[top=25mm,bottom=25mm,left=25mm,right=25mm]{geometry}
\setlength{\emergencystretch}{3em} % prevent overfull lines
\setcounter{secnumdepth}{5}
% Make \paragraph and \subparagraph free-standing
\makeatletter
\ifx\paragraph\undefined\else
  \let\oldparagraph\paragraph
  \renewcommand{\paragraph}{
    \@ifstar
      \xxxParagraphStar
      \xxxParagraphNoStar
  }
  \newcommand{\xxxParagraphStar}[1]{\oldparagraph*{#1}\mbox{}}
  \newcommand{\xxxParagraphNoStar}[1]{\oldparagraph{#1}\mbox{}}
\fi
\ifx\subparagraph\undefined\else
  \let\oldsubparagraph\subparagraph
  \renewcommand{\subparagraph}{
    \@ifstar
      \xxxSubParagraphStar
      \xxxSubParagraphNoStar
  }
  \newcommand{\xxxSubParagraphStar}[1]{\oldsubparagraph*{#1}\mbox{}}
  \newcommand{\xxxSubParagraphNoStar}[1]{\oldsubparagraph{#1}\mbox{}}
\fi
\makeatother

\usepackage{color}
\usepackage{fancyvrb}
\newcommand{\VerbBar}{|}
\newcommand{\VERB}{\Verb[commandchars=\\\{\}]}
\DefineVerbatimEnvironment{Highlighting}{Verbatim}{commandchars=\\\{\}}
% Add ',fontsize=\small' for more characters per line
\usepackage{framed}
\definecolor{shadecolor}{RGB}{241,243,245}
\newenvironment{Shaded}{\begin{snugshade}}{\end{snugshade}}
\newcommand{\AlertTok}[1]{\textcolor[rgb]{0.68,0.00,0.00}{#1}}
\newcommand{\AnnotationTok}[1]{\textcolor[rgb]{0.37,0.37,0.37}{#1}}
\newcommand{\AttributeTok}[1]{\textcolor[rgb]{0.40,0.45,0.13}{#1}}
\newcommand{\BaseNTok}[1]{\textcolor[rgb]{0.68,0.00,0.00}{#1}}
\newcommand{\BuiltInTok}[1]{\textcolor[rgb]{0.00,0.23,0.31}{#1}}
\newcommand{\CharTok}[1]{\textcolor[rgb]{0.13,0.47,0.30}{#1}}
\newcommand{\CommentTok}[1]{\textcolor[rgb]{0.37,0.37,0.37}{#1}}
\newcommand{\CommentVarTok}[1]{\textcolor[rgb]{0.37,0.37,0.37}{\textit{#1}}}
\newcommand{\ConstantTok}[1]{\textcolor[rgb]{0.56,0.35,0.01}{#1}}
\newcommand{\ControlFlowTok}[1]{\textcolor[rgb]{0.00,0.23,0.31}{\textbf{#1}}}
\newcommand{\DataTypeTok}[1]{\textcolor[rgb]{0.68,0.00,0.00}{#1}}
\newcommand{\DecValTok}[1]{\textcolor[rgb]{0.68,0.00,0.00}{#1}}
\newcommand{\DocumentationTok}[1]{\textcolor[rgb]{0.37,0.37,0.37}{\textit{#1}}}
\newcommand{\ErrorTok}[1]{\textcolor[rgb]{0.68,0.00,0.00}{#1}}
\newcommand{\ExtensionTok}[1]{\textcolor[rgb]{0.00,0.23,0.31}{#1}}
\newcommand{\FloatTok}[1]{\textcolor[rgb]{0.68,0.00,0.00}{#1}}
\newcommand{\FunctionTok}[1]{\textcolor[rgb]{0.28,0.35,0.67}{#1}}
\newcommand{\ImportTok}[1]{\textcolor[rgb]{0.00,0.46,0.62}{#1}}
\newcommand{\InformationTok}[1]{\textcolor[rgb]{0.37,0.37,0.37}{#1}}
\newcommand{\KeywordTok}[1]{\textcolor[rgb]{0.00,0.23,0.31}{\textbf{#1}}}
\newcommand{\NormalTok}[1]{\textcolor[rgb]{0.00,0.23,0.31}{#1}}
\newcommand{\OperatorTok}[1]{\textcolor[rgb]{0.37,0.37,0.37}{#1}}
\newcommand{\OtherTok}[1]{\textcolor[rgb]{0.00,0.23,0.31}{#1}}
\newcommand{\PreprocessorTok}[1]{\textcolor[rgb]{0.68,0.00,0.00}{#1}}
\newcommand{\RegionMarkerTok}[1]{\textcolor[rgb]{0.00,0.23,0.31}{#1}}
\newcommand{\SpecialCharTok}[1]{\textcolor[rgb]{0.37,0.37,0.37}{#1}}
\newcommand{\SpecialStringTok}[1]{\textcolor[rgb]{0.13,0.47,0.30}{#1}}
\newcommand{\StringTok}[1]{\textcolor[rgb]{0.13,0.47,0.30}{#1}}
\newcommand{\VariableTok}[1]{\textcolor[rgb]{0.07,0.07,0.07}{#1}}
\newcommand{\VerbatimStringTok}[1]{\textcolor[rgb]{0.13,0.47,0.30}{#1}}
\newcommand{\WarningTok}[1]{\textcolor[rgb]{0.37,0.37,0.37}{\textit{#1}}}

\providecommand{\tightlist}{%
  \setlength{\itemsep}{0pt}\setlength{\parskip}{0pt}}\usepackage{longtable,booktabs,array}
\usepackage{calc} % for calculating minipage widths
% Correct order of tables after \paragraph or \subparagraph
\usepackage{etoolbox}
\makeatletter
\patchcmd\longtable{\par}{\if@noskipsec\mbox{}\fi\par}{}{}
\makeatother
% Allow footnotes in longtable head/foot
\IfFileExists{footnotehyper.sty}{\usepackage{footnotehyper}}{\usepackage{footnote}}
\makesavenoteenv{longtable}
\usepackage{graphicx}
\makeatletter
\newsavebox\pandoc@box
\newcommand*\pandocbounded[1]{% scales image to fit in text height/width
  \sbox\pandoc@box{#1}%
  \Gscale@div\@tempa{\textheight}{\dimexpr\ht\pandoc@box+\dp\pandoc@box\relax}%
  \Gscale@div\@tempb{\linewidth}{\wd\pandoc@box}%
  \ifdim\@tempb\p@<\@tempa\p@\let\@tempa\@tempb\fi% select the smaller of both
  \ifdim\@tempa\p@<\p@\scalebox{\@tempa}{\usebox\pandoc@box}%
  \else\usebox{\pandoc@box}%
  \fi%
}
% Set default figure placement to htbp
\def\fps@figure{htbp}
\makeatother

\usepackage{amsmath,amssymb,amsfonts}
\usepackage{bm}
\usepackage{algorithm2e}
\usepackage{booktabs}
\makeatletter
\@ifpackageloaded{tcolorbox}{}{\usepackage[skins,breakable]{tcolorbox}}
\@ifpackageloaded{fontawesome5}{}{\usepackage{fontawesome5}}
\definecolor{quarto-callout-color}{HTML}{909090}
\definecolor{quarto-callout-note-color}{HTML}{0758E5}
\definecolor{quarto-callout-important-color}{HTML}{CC1914}
\definecolor{quarto-callout-warning-color}{HTML}{EB9113}
\definecolor{quarto-callout-tip-color}{HTML}{00A047}
\definecolor{quarto-callout-caution-color}{HTML}{FC5300}
\definecolor{quarto-callout-color-frame}{HTML}{acacac}
\definecolor{quarto-callout-note-color-frame}{HTML}{4582ec}
\definecolor{quarto-callout-important-color-frame}{HTML}{d9534f}
\definecolor{quarto-callout-warning-color-frame}{HTML}{f0ad4e}
\definecolor{quarto-callout-tip-color-frame}{HTML}{02b875}
\definecolor{quarto-callout-caution-color-frame}{HTML}{fd7e14}
\makeatother
\makeatletter
\@ifpackageloaded{bookmark}{}{\usepackage{bookmark}}
\makeatother
\makeatletter
\@ifpackageloaded{caption}{}{\usepackage{caption}}
\AtBeginDocument{%
\ifdefined\contentsname
  \renewcommand*\contentsname{Table of contents}
\else
  \newcommand\contentsname{Table of contents}
\fi
\ifdefined\listfigurename
  \renewcommand*\listfigurename{List of Figures}
\else
  \newcommand\listfigurename{List of Figures}
\fi
\ifdefined\listtablename
  \renewcommand*\listtablename{List of Tables}
\else
  \newcommand\listtablename{List of Tables}
\fi
\ifdefined\figurename
  \renewcommand*\figurename{Figure}
\else
  \newcommand\figurename{Figure}
\fi
\ifdefined\tablename
  \renewcommand*\tablename{Table}
\else
  \newcommand\tablename{Table}
\fi
}
\@ifpackageloaded{float}{}{\usepackage{float}}
\floatstyle{ruled}
\@ifundefined{c@chapter}{\newfloat{codelisting}{h}{lop}}{\newfloat{codelisting}{h}{lop}[chapter]}
\floatname{codelisting}{Listing}
\newcommand*\listoflistings{\listof{codelisting}{List of Listings}}
\makeatother
\makeatletter
\makeatother
\makeatletter
\@ifpackageloaded{caption}{}{\usepackage{caption}}
\@ifpackageloaded{subcaption}{}{\usepackage{subcaption}}
\makeatother

\usepackage{bookmark}

\IfFileExists{xurl.sty}{\usepackage{xurl}}{} % add URL line breaks if available
\urlstyle{same} % disable monospaced font for URLs
\hypersetup{
  pdftitle={Introduction to the Finite Element Method},
  pdfauthor={Shmuel Osovski},
  hidelinks,
  pdfcreator={LaTeX via pandoc}}


\title{Introduction to the Finite Element Method}
\author{Shmuel Osovski}
\date{2025-01-01}

\begin{document}
\frontmatter
\maketitle

\renewcommand*\contentsname{Table of contents}
{
\setcounter{tocdepth}{1}
\tableofcontents
}

\mainmatter
\bookmarksetup{startatroot}

\chapter{Preface}\label{preface}

These lecture notes accompany the course \emph{Introduction to the
Finite Element Method}. They cover the mathematical foundations, 1D and
multi-dimensional finite element formulations, implementation details,
and error analysis.

The notes are designed so that students can follow along during lectures
without needing to copy derivations --- instead, focus on understanding
the concepts and annotating where needed.

\section{How to Use These Notes}\label{how-to-use-these-notes}

The chapters follow the lecture sequence. Each chapter contains:

\begin{itemize}
\tightlist
\item
  Mathematical derivations presented step by step
\item
  Python code demonstrating the concepts computationally (click to
  expand code blocks)
\item
  Key definitions and results highlighted in callout boxes
\end{itemize}

\section{Prerequisites}\label{prerequisites}

Students should be familiar with:

\begin{itemize}
\tightlist
\item
  Linear algebra (matrix operations, eigenvalues)
\item
  Calculus (integration by parts, partial derivatives)
\item
  Basic Python programming
\item
  Introductory continuum mechanics
\end{itemize}

\part{Part I: Foundations}

\chapter{Introduction and 1D Setup}\label{introduction-and-1d-setup}

\section{Setup}\label{setup}

\section{The Strong Form}\label{the-strong-form}

\begin{Shaded}
\begin{Highlighting}[]
\InformationTok{\textasciigrave{}\textasciigrave{}\textasciigrave{}\{python\}}
\InformationTok{\#| code{-}fold: true}
\InformationTok{fem.display\_strong\_form()}
\InformationTok{\textasciigrave{}\textasciigrave{}\textasciigrave{}}
\end{Highlighting}
\end{Shaded}

\begin{verbatim}
<IPython.core.display.HTML object>
\end{verbatim}

$\displaystyle A{\left(x \right)} \frac{d^{2}}{d x^{2}} y{\left(x \right)} + B{\left(x \right)} \frac{d}{d x} y{\left(x \right)} + C{\left(x \right)} y{\left(x \right)} = F{\left(x \right)}$

\begin{verbatim}
<IPython.core.display.HTML object>
\end{verbatim}

$\displaystyle y{\left(0 \right)} = 0$

$\displaystyle y{\left(L \right)} = 0$

\begin{verbatim}
<IPython.core.display.HTML object>
\end{verbatim}

\section{Deriving the Weak Form}\label{deriving-the-weak-form}

\begin{Shaded}
\begin{Highlighting}[]
\InformationTok{\textasciigrave{}\textasciigrave{}\textasciigrave{}\{python\}}
\InformationTok{\#| code{-}fold: true}
\InformationTok{weak\_eq = fem.derive\_weak\_form()}
\InformationTok{\textasciigrave{}\textasciigrave{}\textasciigrave{}}
\end{Highlighting}
\end{Shaded}

\begin{verbatim}
<IPython.core.display.HTML object>
\end{verbatim}

\begin{verbatim}
<IPython.core.display.HTML object>
\end{verbatim}

\begin{verbatim}
<IPython.core.display.HTML object>
\end{verbatim}

\begin{verbatim}
<IPython.core.display.HTML object>
\end{verbatim}

\begin{verbatim}
<IPython.core.display.HTML object>
\end{verbatim}

\begin{verbatim}
<IPython.core.display.HTML object>
\end{verbatim}

\begin{verbatim}
<IPython.core.display.HTML object>
\end{verbatim}

\begin{verbatim}
<IPython.core.display.HTML object>
\end{verbatim}

$\displaystyle \left(A{\left(x \right)} \frac{d^{2}}{d x^{2}} y{\left(x \right)} + B{\left(x \right)} \frac{d}{d x} y{\left(x \right)} + C{\left(x \right)} y{\left(x \right)}\right) v{\left(x \right)} = F{\left(x \right)} v{\left(x \right)}$

\begin{verbatim}
<IPython.core.display.HTML object>
\end{verbatim}

\begin{verbatim}
<IPython.core.display.HTML object>
\end{verbatim}

\begin{verbatim}
<IPython.core.display.HTML object>
\end{verbatim}

\begin{verbatim}
<IPython.core.display.HTML object>
\end{verbatim}

\begin{verbatim}
<IPython.core.display.HTML object>
\end{verbatim}

\begin{verbatim}
<IPython.core.display.HTML object>
\end{verbatim}

\begin{verbatim}
<IPython.core.display.HTML object>
\end{verbatim}

$\displaystyle \int\limits_{0}^{L} \left(A{\left(x \right)} \frac{d^{2}}{d x^{2}} y{\left(x \right)} + B{\left(x \right)} \frac{d}{d x} y{\left(x \right)} + C{\left(x \right)} y{\left(x \right)}\right) v{\left(x \right)}\, dx = \int\limits_{0}^{L} F{\left(x \right)} v{\left(x \right)}\, dx$

\begin{verbatim}
<IPython.core.display.HTML object>
\end{verbatim}

\begin{verbatim}
<IPython.core.display.HTML object>
\end{verbatim}

\begin{verbatim}
<IPython.core.display.HTML object>
\end{verbatim}

\begin{verbatim}
<IPython.core.display.HTML object>
\end{verbatim}

\begin{verbatim}
<IPython.core.display.HTML object>
\end{verbatim}

\begin{verbatim}
<IPython.core.display.HTML object>
\end{verbatim}

\begin{verbatim}
<IPython.core.display.HTML object>
\end{verbatim}

\begin{verbatim}
<IPython.core.display.HTML object>
\end{verbatim}

$\displaystyle \int\limits_{0}^{L} A{\left(x \right)} v{\left(x \right)} \frac{d^{2}}{d x^{2}} y{\left(x \right)}\, dx + \int\limits_{0}^{L} B{\left(x \right)} v{\left(x \right)} \frac{d}{d x} y{\left(x \right)}\, dx + \int\limits_{0}^{L} C{\left(x \right)} v{\left(x \right)} y{\left(x \right)}\, dx = \int\limits_{0}^{L} F{\left(x \right)} v{\left(x \right)}\, dx$

\begin{verbatim}
<IPython.core.display.HTML object>
\end{verbatim}

\begin{verbatim}
<IPython.core.display.HTML object>
\end{verbatim}

\begin{verbatim}
<IPython.core.display.HTML object>
\end{verbatim}

\begin{verbatim}
<IPython.core.display.HTML object>
\end{verbatim}

\begin{verbatim}
<IPython.core.display.HTML object>
\end{verbatim}

\begin{verbatim}
<IPython.core.display.HTML object>
\end{verbatim}

\begin{verbatim}
<IPython.core.display.HTML object>
\end{verbatim}

\begin{verbatim}
<IPython.core.display.HTML object>
\end{verbatim}

\begin{verbatim}
<IPython.core.display.HTML object>
\end{verbatim}

\begin{verbatim}
<IPython.core.display.HTML object>
\end{verbatim}

\begin{verbatim}
<IPython.core.display.HTML object>
\end{verbatim}

\begin{verbatim}
<IPython.core.display.HTML object>
\end{verbatim}

\begin{verbatim}
<IPython.core.display.HTML object>
\end{verbatim}

\begin{verbatim}
<IPython.core.display.HTML object>
\end{verbatim}

\begin{verbatim}
<IPython.core.display.HTML object>
\end{verbatim}

\begin{verbatim}
<IPython.core.display.HTML object>
\end{verbatim}

\begin{verbatim}
<IPython.core.display.HTML object>
\end{verbatim}

\begin{verbatim}
<IPython.core.display.HTML object>
\end{verbatim}

\begin{verbatim}
<IPython.core.display.HTML object>
\end{verbatim}

\begin{verbatim}
<IPython.core.display.HTML object>
\end{verbatim}

\begin{verbatim}
<IPython.core.display.HTML object>
\end{verbatim}

\begin{verbatim}
<IPython.core.display.HTML object>
\end{verbatim}

\begin{verbatim}
<IPython.core.display.HTML object>
\end{verbatim}

$\displaystyle \int\limits_{0}^{L} A{\left(x \right)} v{\left(x \right)} \frac{d^{2}}{d x^{2}} y{\left(x \right)}\, dx = - A{\left(0 \right)} v{\left(0 \right)} \left. \frac{d}{d x} y{\left(x \right)} \right|_{\substack{ x=0 }} + A{\left(L \right)} v{\left(L \right)} \frac{d}{d L} y{\left(L \right)} - \int\limits_{0}^{L} \left(A{\left(x \right)} \frac{d}{d x} v{\left(x \right)} + v{\left(x \right)} \frac{d}{d x} A{\left(x \right)}\right) \frac{d}{d x} y{\left(x \right)}\, dx$

\begin{verbatim}
<IPython.core.display.HTML object>
\end{verbatim}

\begin{verbatim}
<IPython.core.display.HTML object>
\end{verbatim}

\begin{verbatim}
<IPython.core.display.HTML object>
\end{verbatim}

\begin{verbatim}
<IPython.core.display.HTML object>
\end{verbatim}

\begin{verbatim}
<IPython.core.display.HTML object>
\end{verbatim}

\begin{verbatim}
<IPython.core.display.HTML object>
\end{verbatim}

\begin{verbatim}
<IPython.core.display.HTML object>
\end{verbatim}

\begin{verbatim}
<IPython.core.display.HTML object>
\end{verbatim}

\begin{verbatim}
<IPython.core.display.HTML object>
\end{verbatim}

\begin{verbatim}
<IPython.core.display.HTML object>
\end{verbatim}

\begin{verbatim}
<IPython.core.display.HTML object>
\end{verbatim}

\begin{verbatim}
<IPython.core.display.HTML object>
\end{verbatim}

\begin{verbatim}
<IPython.core.display.HTML object>
\end{verbatim}

\begin{verbatim}
<IPython.core.display.HTML object>
\end{verbatim}

\begin{verbatim}
<IPython.core.display.HTML object>
\end{verbatim}

\begin{verbatim}
<IPython.core.display.HTML object>
\end{verbatim}

\begin{verbatim}
<IPython.core.display.HTML object>
\end{verbatim}

\begin{verbatim}
<IPython.core.display.HTML object>
\end{verbatim}

$\displaystyle - \int\limits_{0}^{L} \left(A{\left(x \right)} \frac{d}{d x} v{\left(x \right)} + v{\left(x \right)} \frac{d}{d x} A{\left(x \right)}\right) \frac{d}{d x} y{\left(x \right)}\, dx + \int\limits_{0}^{L} B{\left(x \right)} v{\left(x \right)} \frac{d}{d x} y{\left(x \right)}\, dx + \int\limits_{0}^{L} C{\left(x \right)} v{\left(x \right)} y{\left(x \right)}\, dx = \int\limits_{0}^{L} F{\left(x \right)} v{\left(x \right)}\, dx$

\section{Shape Functions for
Discretization}\label{shape-functions-for-discretization}

\begin{Shaded}
\begin{Highlighting}[]
\InformationTok{\textasciigrave{}\textasciigrave{}\textasciigrave{}\{python\}}
\InformationTok{\#| code{-}fold: true}
\InformationTok{N = 5}
\InformationTok{phi, a, b, y\_approx, v\_approx = fem.define\_shape\_functions(N)}
\InformationTok{\textasciigrave{}\textasciigrave{}\textasciigrave{}}
\end{Highlighting}
\end{Shaded}

\begin{verbatim}
<IPython.core.display.HTML object>
\end{verbatim}

\begin{verbatim}
<IPython.core.display.HTML object>
\end{verbatim}

\begin{verbatim}
<IPython.core.display.HTML object>
\end{verbatim}

\begin{verbatim}
<IPython.core.display.HTML object>
\end{verbatim}

\begin{verbatim}
<IPython.core.display.HTML object>
\end{verbatim}

\begin{verbatim}
<IPython.core.display.HTML object>
\end{verbatim}

\begin{verbatim}
<IPython.core.display.HTML object>
\end{verbatim}

\begin{verbatim}
<IPython.core.display.HTML object>
\end{verbatim}

\begin{verbatim}
<IPython.core.display.HTML object>
\end{verbatim}

\begin{verbatim}
<IPython.core.display.HTML object>
\end{verbatim}

$\displaystyle \left[ \phi_{1}{\left(x \right)}, \  \phi_{2}{\left(x \right)}, \  \phi_{3}{\left(x \right)}, \  \phi_{4}{\left(x \right)}, \  \phi_{5}{\left(x \right)}\right]$

\begin{verbatim}
<IPython.core.display.HTML object>
\end{verbatim}

\begin{verbatim}
<IPython.core.display.HTML object>
\end{verbatim}

\begin{verbatim}
<IPython.core.display.HTML object>
\end{verbatim}

\begin{verbatim}
<IPython.core.display.HTML object>
\end{verbatim}

\begin{verbatim}
<IPython.core.display.HTML object>
\end{verbatim}

\begin{verbatim}
<IPython.core.display.HTML object>
\end{verbatim}

\begin{verbatim}
<IPython.core.display.HTML object>
\end{verbatim}

\begin{verbatim}
<IPython.core.display.HTML object>
\end{verbatim}

\begin{verbatim}
<IPython.core.display.HTML object>
\end{verbatim}

$\displaystyle y{\left(x \right)} = a_{1} \phi_{1}{\left(x \right)} + a_{2} \phi_{2}{\left(x \right)} + a_{3} \phi_{3}{\left(x \right)} + a_{4} \phi_{4}{\left(x \right)} + a_{5} \phi_{5}{\left(x \right)}$

$\displaystyle v{\left(x \right)} = b_{1} \phi_{1}{\left(x \right)} + b_{2} \phi_{2}{\left(x \right)} + b_{3} \phi_{3}{\left(x \right)} + b_{4} \phi_{4}{\left(x \right)} + b_{5} \phi_{5}{\left(x \right)}$

\begin{verbatim}
<IPython.core.display.HTML object>
\end{verbatim}

\begin{verbatim}
<IPython.core.display.HTML object>
\end{verbatim}

\begin{verbatim}
<IPython.core.display.HTML object>
\end{verbatim}

\begin{verbatim}
<IPython.core.display.HTML object>
\end{verbatim}

\begin{verbatim}
<IPython.core.display.HTML object>
\end{verbatim}

\begin{verbatim}
<IPython.core.display.HTML object>
\end{verbatim}

\begin{verbatim}
<IPython.core.display.HTML object>
\end{verbatim}

\begin{verbatim}
<IPython.core.display.HTML object>
\end{verbatim}

\begin{verbatim}
<IPython.core.display.HTML object>
\end{verbatim}

\section{Assembling the FEM System}\label{assembling-the-fem-system}

\begin{Shaded}
\begin{Highlighting}[]
\InformationTok{\textasciigrave{}\textasciigrave{}\textasciigrave{}\{python\}}
\InformationTok{\#| code{-}fold: true}
\InformationTok{K, F\_vec = fem.assemble\_fem\_system(weak\_eq, phi, a, b, y\_approx, v\_approx, N)}
\InformationTok{\textasciigrave{}\textasciigrave{}\textasciigrave{}}
\end{Highlighting}
\end{Shaded}

\begin{verbatim}
<IPython.core.display.HTML object>
\end{verbatim}

\begin{verbatim}
<IPython.core.display.HTML object>
\end{verbatim}

\begin{verbatim}
<IPython.core.display.HTML object>
\end{verbatim}

\begin{verbatim}
<IPython.core.display.HTML object>
\end{verbatim}

\begin{verbatim}
<IPython.core.display.HTML object>
\end{verbatim}

\begin{verbatim}
<IPython.core.display.HTML object>
\end{verbatim}

\begin{verbatim}
<IPython.core.display.HTML object>
\end{verbatim}

\begin{verbatim}
<IPython.core.display.HTML object>
\end{verbatim}

\begin{verbatim}
<IPython.core.display.HTML object>
\end{verbatim}

\begin{verbatim}
<IPython.core.display.HTML object>
\end{verbatim}

\begin{verbatim}
<IPython.core.display.HTML object>
\end{verbatim}

\begin{verbatim}
<IPython.core.display.HTML object>
\end{verbatim}

\begin{verbatim}
<IPython.core.display.HTML object>
\end{verbatim}

\begin{verbatim}
<IPython.core.display.HTML object>
\end{verbatim}

\begin{verbatim}
<IPython.core.display.HTML object>
\end{verbatim}

\begin{verbatim}
<IPython.core.display.HTML object>
\end{verbatim}

\begin{verbatim}
<IPython.core.display.HTML object>
\end{verbatim}

\begin{verbatim}
<IPython.core.display.HTML object>
\end{verbatim}

\begin{verbatim}
<IPython.core.display.HTML object>
\end{verbatim}

\begin{verbatim}
<IPython.core.display.HTML object>
\end{verbatim}

\begin{verbatim}
<IPython.core.display.HTML object>
\end{verbatim}

\begin{verbatim}
<IPython.core.display.HTML object>
\end{verbatim}

\begin{verbatim}
<IPython.core.display.HTML object>
\end{verbatim}

\begin{verbatim}
<IPython.core.display.HTML object>
\end{verbatim}

$\displaystyle K_{ij} = \int_0^L \frac{d\phi_i}{dx} \frac{d\phi_j}{dx} dx + \int_0^L \phi_i(x) \frac{d\phi_j}{dx} dx + \int_0^L \phi_i(x) \phi_j(x) dx$

\begin{verbatim}
<IPython.core.display.HTML object>
\end{verbatim}

$\displaystyle F_i = \int_0^L \phi_i(x) F(x) dx$

\begin{verbatim}
<IPython.core.display.HTML object>
\end{verbatim}

\begin{verbatim}
<IPython.core.display.HTML object>
\end{verbatim}

\begin{verbatim}
<IPython.core.display.HTML object>
\end{verbatim}

\begin{verbatim}
<IPython.core.display.HTML object>
\end{verbatim}

\section{Putting It All Together}\label{putting-it-all-together}

\begin{Shaded}
\begin{Highlighting}[]
\InformationTok{\textasciigrave{}\textasciigrave{}\textasciigrave{}\{python\}}
\InformationTok{\#| code{-}fold: true}
\InformationTok{fem.finite\_element\_method\_1d\_demo()}
\InformationTok{\textasciigrave{}\textasciigrave{}\textasciigrave{}}
\end{Highlighting}
\end{Shaded}

\begin{verbatim}
<IPython.core.display.HTML object>
\end{verbatim}

\begin{verbatim}
<IPython.core.display.HTML object>
\end{verbatim}

\begin{verbatim}
<IPython.core.display.HTML object>
\end{verbatim}

\begin{verbatim}
<IPython.core.display.HTML object>
\end{verbatim}

\begin{verbatim}
<IPython.core.display.HTML object>
\end{verbatim}

$\displaystyle A{\left(x \right)} \frac{d^{2}}{d x^{2}} y{\left(x \right)} + B{\left(x \right)} \frac{d}{d x} y{\left(x \right)} + C{\left(x \right)} y{\left(x \right)} = F{\left(x \right)}$

\begin{verbatim}
<IPython.core.display.HTML object>
\end{verbatim}

$\displaystyle y{\left(0 \right)} = 0$

$\displaystyle y{\left(L \right)} = 0$

\begin{verbatim}
<IPython.core.display.HTML object>
\end{verbatim}

\begin{verbatim}
<IPython.core.display.HTML object>
\end{verbatim}

\begin{verbatim}
<IPython.core.display.HTML object>
\end{verbatim}

\begin{verbatim}
<IPython.core.display.HTML object>
\end{verbatim}

\begin{verbatim}
<IPython.core.display.HTML object>
\end{verbatim}

\begin{verbatim}
<IPython.core.display.HTML object>
\end{verbatim}

\begin{verbatim}
<IPython.core.display.HTML object>
\end{verbatim}

\begin{verbatim}
<IPython.core.display.HTML object>
\end{verbatim}

\begin{verbatim}
<IPython.core.display.HTML object>
\end{verbatim}

\begin{verbatim}
<IPython.core.display.HTML object>
\end{verbatim}

\begin{verbatim}
<IPython.core.display.HTML object>
\end{verbatim}

\begin{verbatim}
<IPython.core.display.HTML object>
\end{verbatim}

$\displaystyle \left(A{\left(x \right)} \frac{d^{2}}{d x^{2}} y{\left(x \right)} + B{\left(x \right)} \frac{d}{d x} y{\left(x \right)} + C{\left(x \right)} y{\left(x \right)}\right) v{\left(x \right)} = F{\left(x \right)} v{\left(x \right)}$

\begin{verbatim}
<IPython.core.display.HTML object>
\end{verbatim}

\begin{verbatim}
<IPython.core.display.HTML object>
\end{verbatim}

\begin{verbatim}
<IPython.core.display.HTML object>
\end{verbatim}

\begin{verbatim}
<IPython.core.display.HTML object>
\end{verbatim}

\begin{verbatim}
<IPython.core.display.HTML object>
\end{verbatim}

\begin{verbatim}
<IPython.core.display.HTML object>
\end{verbatim}

\begin{verbatim}
<IPython.core.display.HTML object>
\end{verbatim}

$\displaystyle \int\limits_{0}^{L} \left(A{\left(x \right)} \frac{d^{2}}{d x^{2}} y{\left(x \right)} + B{\left(x \right)} \frac{d}{d x} y{\left(x \right)} + C{\left(x \right)} y{\left(x \right)}\right) v{\left(x \right)}\, dx = \int\limits_{0}^{L} F{\left(x \right)} v{\left(x \right)}\, dx$

\begin{verbatim}
<IPython.core.display.HTML object>
\end{verbatim}

\begin{verbatim}
<IPython.core.display.HTML object>
\end{verbatim}

\begin{verbatim}
<IPython.core.display.HTML object>
\end{verbatim}

\begin{verbatim}
<IPython.core.display.HTML object>
\end{verbatim}

\begin{verbatim}
<IPython.core.display.HTML object>
\end{verbatim}

\begin{verbatim}
<IPython.core.display.HTML object>
\end{verbatim}

\begin{verbatim}
<IPython.core.display.HTML object>
\end{verbatim}

\begin{verbatim}
<IPython.core.display.HTML object>
\end{verbatim}

$\displaystyle \int\limits_{0}^{L} A{\left(x \right)} v{\left(x \right)} \frac{d^{2}}{d x^{2}} y{\left(x \right)}\, dx + \int\limits_{0}^{L} B{\left(x \right)} v{\left(x \right)} \frac{d}{d x} y{\left(x \right)}\, dx + \int\limits_{0}^{L} C{\left(x \right)} v{\left(x \right)} y{\left(x \right)}\, dx = \int\limits_{0}^{L} F{\left(x \right)} v{\left(x \right)}\, dx$

\begin{verbatim}
<IPython.core.display.HTML object>
\end{verbatim}

\begin{verbatim}
<IPython.core.display.HTML object>
\end{verbatim}

\begin{verbatim}
<IPython.core.display.HTML object>
\end{verbatim}

\begin{verbatim}
<IPython.core.display.HTML object>
\end{verbatim}

\begin{verbatim}
<IPython.core.display.HTML object>
\end{verbatim}

\begin{verbatim}
<IPython.core.display.HTML object>
\end{verbatim}

\begin{verbatim}
<IPython.core.display.HTML object>
\end{verbatim}

\begin{verbatim}
<IPython.core.display.HTML object>
\end{verbatim}

\begin{verbatim}
<IPython.core.display.HTML object>
\end{verbatim}

\begin{verbatim}
<IPython.core.display.HTML object>
\end{verbatim}

\begin{verbatim}
<IPython.core.display.HTML object>
\end{verbatim}

\begin{verbatim}
<IPython.core.display.HTML object>
\end{verbatim}

\begin{verbatim}
<IPython.core.display.HTML object>
\end{verbatim}

\begin{verbatim}
<IPython.core.display.HTML object>
\end{verbatim}

\begin{verbatim}
<IPython.core.display.HTML object>
\end{verbatim}

\begin{verbatim}
<IPython.core.display.HTML object>
\end{verbatim}

\begin{verbatim}
<IPython.core.display.HTML object>
\end{verbatim}

\begin{verbatim}
<IPython.core.display.HTML object>
\end{verbatim}

\begin{verbatim}
<IPython.core.display.HTML object>
\end{verbatim}

\begin{verbatim}
<IPython.core.display.HTML object>
\end{verbatim}

\begin{verbatim}
<IPython.core.display.HTML object>
\end{verbatim}

\begin{verbatim}
<IPython.core.display.HTML object>
\end{verbatim}

\begin{verbatim}
<IPython.core.display.HTML object>
\end{verbatim}

$\displaystyle \int\limits_{0}^{L} A{\left(x \right)} v{\left(x \right)} \frac{d^{2}}{d x^{2}} y{\left(x \right)}\, dx = - A{\left(0 \right)} v{\left(0 \right)} \left. \frac{d}{d x} y{\left(x \right)} \right|_{\substack{ x=0 }} + A{\left(L \right)} v{\left(L \right)} \frac{d}{d L} y{\left(L \right)} - \int\limits_{0}^{L} \left(A{\left(x \right)} \frac{d}{d x} v{\left(x \right)} + v{\left(x \right)} \frac{d}{d x} A{\left(x \right)}\right) \frac{d}{d x} y{\left(x \right)}\, dx$

\begin{verbatim}
<IPython.core.display.HTML object>
\end{verbatim}

\begin{verbatim}
<IPython.core.display.HTML object>
\end{verbatim}

\begin{verbatim}
<IPython.core.display.HTML object>
\end{verbatim}

\begin{verbatim}
<IPython.core.display.HTML object>
\end{verbatim}

\begin{verbatim}
<IPython.core.display.HTML object>
\end{verbatim}

\begin{verbatim}
<IPython.core.display.HTML object>
\end{verbatim}

\begin{verbatim}
<IPython.core.display.HTML object>
\end{verbatim}

\begin{verbatim}
<IPython.core.display.HTML object>
\end{verbatim}

\begin{verbatim}
<IPython.core.display.HTML object>
\end{verbatim}

\begin{verbatim}
<IPython.core.display.HTML object>
\end{verbatim}

\begin{verbatim}
<IPython.core.display.HTML object>
\end{verbatim}

\begin{verbatim}
<IPython.core.display.HTML object>
\end{verbatim}

\begin{verbatim}
<IPython.core.display.HTML object>
\end{verbatim}

\begin{verbatim}
<IPython.core.display.HTML object>
\end{verbatim}

\begin{verbatim}
<IPython.core.display.HTML object>
\end{verbatim}

\begin{verbatim}
<IPython.core.display.HTML object>
\end{verbatim}

\begin{verbatim}
<IPython.core.display.HTML object>
\end{verbatim}

\begin{verbatim}
<IPython.core.display.HTML object>
\end{verbatim}

$\displaystyle - \int\limits_{0}^{L} \left(A{\left(x \right)} \frac{d}{d x} v{\left(x \right)} + v{\left(x \right)} \frac{d}{d x} A{\left(x \right)}\right) \frac{d}{d x} y{\left(x \right)}\, dx + \int\limits_{0}^{L} B{\left(x \right)} v{\left(x \right)} \frac{d}{d x} y{\left(x \right)}\, dx + \int\limits_{0}^{L} C{\left(x \right)} v{\left(x \right)} y{\left(x \right)}\, dx = \int\limits_{0}^{L} F{\left(x \right)} v{\left(x \right)}\, dx$

\begin{verbatim}
<IPython.core.display.HTML object>
\end{verbatim}

\begin{verbatim}
<IPython.core.display.HTML object>
\end{verbatim}

\begin{verbatim}
<IPython.core.display.HTML object>
\end{verbatim}

\begin{verbatim}
<IPython.core.display.HTML object>
\end{verbatim}

\begin{verbatim}
<IPython.core.display.HTML object>
\end{verbatim}

\begin{verbatim}
<IPython.core.display.HTML object>
\end{verbatim}

\begin{verbatim}
<IPython.core.display.HTML object>
\end{verbatim}

\begin{verbatim}
<IPython.core.display.HTML object>
\end{verbatim}

\begin{verbatim}
<IPython.core.display.HTML object>
\end{verbatim}

\begin{verbatim}
<IPython.core.display.HTML object>
\end{verbatim}

\begin{verbatim}
<IPython.core.display.HTML object>
\end{verbatim}

\begin{verbatim}
<IPython.core.display.HTML object>
\end{verbatim}

\begin{verbatim}
<IPython.core.display.HTML object>
\end{verbatim}

$\displaystyle \left[ \phi_{1}{\left(x \right)}, \  \phi_{2}{\left(x \right)}, \  \phi_{3}{\left(x \right)}, \  \phi_{4}{\left(x \right)}, \  \phi_{5}{\left(x \right)}\right]$

\begin{verbatim}
<IPython.core.display.HTML object>
\end{verbatim}

\begin{verbatim}
<IPython.core.display.HTML object>
\end{verbatim}

\begin{verbatim}
<IPython.core.display.HTML object>
\end{verbatim}

\begin{verbatim}
<IPython.core.display.HTML object>
\end{verbatim}

\begin{verbatim}
<IPython.core.display.HTML object>
\end{verbatim}

\begin{verbatim}
<IPython.core.display.HTML object>
\end{verbatim}

\begin{verbatim}
<IPython.core.display.HTML object>
\end{verbatim}

\begin{verbatim}
<IPython.core.display.HTML object>
\end{verbatim}

\begin{verbatim}
<IPython.core.display.HTML object>
\end{verbatim}

$\displaystyle y{\left(x \right)} = a_{1} \phi_{1}{\left(x \right)} + a_{2} \phi_{2}{\left(x \right)} + a_{3} \phi_{3}{\left(x \right)} + a_{4} \phi_{4}{\left(x \right)} + a_{5} \phi_{5}{\left(x \right)}$

$\displaystyle v{\left(x \right)} = b_{1} \phi_{1}{\left(x \right)} + b_{2} \phi_{2}{\left(x \right)} + b_{3} \phi_{3}{\left(x \right)} + b_{4} \phi_{4}{\left(x \right)} + b_{5} \phi_{5}{\left(x \right)}$

\begin{verbatim}
<IPython.core.display.HTML object>
\end{verbatim}

\begin{verbatim}
<IPython.core.display.HTML object>
\end{verbatim}

\begin{verbatim}
<IPython.core.display.HTML object>
\end{verbatim}

\begin{verbatim}
<IPython.core.display.HTML object>
\end{verbatim}

\begin{verbatim}
<IPython.core.display.HTML object>
\end{verbatim}

\begin{verbatim}
<IPython.core.display.HTML object>
\end{verbatim}

\begin{verbatim}
<IPython.core.display.HTML object>
\end{verbatim}

\begin{verbatim}
<IPython.core.display.HTML object>
\end{verbatim}

\begin{verbatim}
<IPython.core.display.HTML object>
\end{verbatim}

\begin{verbatim}
<IPython.core.display.HTML object>
\end{verbatim}

\begin{verbatim}
<IPython.core.display.HTML object>
\end{verbatim}

\begin{verbatim}
<IPython.core.display.HTML object>
\end{verbatim}

\begin{verbatim}
<IPython.core.display.HTML object>
\end{verbatim}

\begin{verbatim}
<IPython.core.display.HTML object>
\end{verbatim}

\begin{verbatim}
<IPython.core.display.HTML object>
\end{verbatim}

\begin{verbatim}
<IPython.core.display.HTML object>
\end{verbatim}

\begin{verbatim}
<IPython.core.display.HTML object>
\end{verbatim}

\begin{verbatim}
<IPython.core.display.HTML object>
\end{verbatim}

\begin{verbatim}
<IPython.core.display.HTML object>
\end{verbatim}

\begin{verbatim}
<IPython.core.display.HTML object>
\end{verbatim}

\begin{verbatim}
<IPython.core.display.HTML object>
\end{verbatim}

\begin{verbatim}
<IPython.core.display.HTML object>
\end{verbatim}

\begin{verbatim}
<IPython.core.display.HTML object>
\end{verbatim}

\begin{verbatim}
<IPython.core.display.HTML object>
\end{verbatim}

\begin{verbatim}
<IPython.core.display.HTML object>
\end{verbatim}

\begin{verbatim}
<IPython.core.display.HTML object>
\end{verbatim}

\begin{verbatim}
<IPython.core.display.HTML object>
\end{verbatim}

\begin{verbatim}
<IPython.core.display.HTML object>
\end{verbatim}

\begin{verbatim}
<IPython.core.display.HTML object>
\end{verbatim}

\begin{verbatim}
<IPython.core.display.HTML object>
\end{verbatim}

\begin{verbatim}
<IPython.core.display.HTML object>
\end{verbatim}

\begin{verbatim}
<IPython.core.display.HTML object>
\end{verbatim}

\begin{verbatim}
<IPython.core.display.HTML object>
\end{verbatim}

\begin{verbatim}
<IPython.core.display.HTML object>
\end{verbatim}

\begin{verbatim}
<IPython.core.display.HTML object>
\end{verbatim}

\begin{verbatim}
<IPython.core.display.HTML object>
\end{verbatim}

$\displaystyle K_{ij} = \int_0^L \frac{d\phi_i}{dx} \frac{d\phi_j}{dx} dx + \int_0^L \phi_i(x) \frac{d\phi_j}{dx} dx + \int_0^L \phi_i(x) \phi_j(x) dx$

\begin{verbatim}
<IPython.core.display.HTML object>
\end{verbatim}

$\displaystyle F_i = \int_0^L \phi_i(x) F(x) dx$

\begin{verbatim}
<IPython.core.display.HTML object>
\end{verbatim}

\begin{verbatim}
<IPython.core.display.HTML object>
\end{verbatim}

\begin{verbatim}
<IPython.core.display.HTML object>
\end{verbatim}

\begin{verbatim}
<IPython.core.display.HTML object>
\end{verbatim}

\begin{verbatim}
<IPython.core.display.HTML object>
\end{verbatim}

\begin{verbatim}
<IPython.core.display.HTML object>
\end{verbatim}

\begin{verbatim}
<IPython.core.display.HTML object>
\end{verbatim}

\begin{verbatim}
<IPython.core.display.HTML object>
\end{verbatim}

$\displaystyle \text{True}$

\begin{verbatim}
<IPython.core.display.HTML object>
\end{verbatim}

\section{Example: Convergence with Mesh
Refinement}\label{sec-convergence-example}

The following plots show the FEM solution for increasing numbers of
elements, demonstrating how the numerical solution converges to the
exact solution as the mesh is refined.

\begin{Shaded}
\begin{Highlighting}[]
\ImportTok{import}\NormalTok{ matplotlib.pyplot }\ImportTok{as}\NormalTok{ plt}

\NormalTok{fig, axes }\OperatorTok{=}\NormalTok{ plt.subplots(}\DecValTok{2}\NormalTok{, }\DecValTok{2}\NormalTok{, figsize}\OperatorTok{=}\NormalTok{(}\DecValTok{10}\NormalTok{, }\DecValTok{8}\NormalTok{))}
\NormalTok{x\_exact }\OperatorTok{=}\NormalTok{ np.linspace(}\DecValTok{0}\NormalTok{, }\DecValTok{1}\NormalTok{, }\DecValTok{100}\NormalTok{)}
\NormalTok{y\_exact }\OperatorTok{=} \FloatTok{0.5} \OperatorTok{*}\NormalTok{ x\_exact }\OperatorTok{*}\NormalTok{ (}\DecValTok{1} \OperatorTok{{-}}\NormalTok{ x\_exact)}

\ControlFlowTok{for}\NormalTok{ idx, n\_elem }\KeywordTok{in} \BuiltInTok{enumerate}\NormalTok{([}\DecValTok{2}\NormalTok{, }\DecValTok{5}\NormalTok{, }\DecValTok{10}\NormalTok{, }\DecValTok{20}\NormalTok{]):}
\NormalTok{    ax }\OperatorTok{=}\NormalTok{ axes[idx }\OperatorTok{//} \DecValTok{2}\NormalTok{, idx }\OperatorTok{\%} \DecValTok{2}\NormalTok{]}

    \CommentTok{\# {-}{-}{-} Compute FEM solution {-}{-}{-}}
\NormalTok{    n\_nodes }\OperatorTok{=}\NormalTok{ n\_elem }\OperatorTok{+} \DecValTok{1}
\NormalTok{    h }\OperatorTok{=} \FloatTok{1.0} \OperatorTok{/}\NormalTok{ n\_elem}
\NormalTok{    K }\OperatorTok{=}\NormalTok{ np.zeros((n\_nodes, n\_nodes))}
\NormalTok{    F\_vec }\OperatorTok{=}\NormalTok{ np.zeros(n\_nodes)}
    \ControlFlowTok{for}\NormalTok{ e }\KeywordTok{in} \BuiltInTok{range}\NormalTok{(n\_elem):}
\NormalTok{        k\_e }\OperatorTok{=}\NormalTok{ np.array([[}\DecValTok{1}\NormalTok{, }\OperatorTok{{-}}\DecValTok{1}\NormalTok{], [}\OperatorTok{{-}}\DecValTok{1}\NormalTok{, }\DecValTok{1}\NormalTok{]]) }\OperatorTok{/}\NormalTok{ h}
\NormalTok{        f\_e }\OperatorTok{=}\NormalTok{ np.array([h }\OperatorTok{/} \DecValTok{2}\NormalTok{, h }\OperatorTok{/} \DecValTok{2}\NormalTok{])}
\NormalTok{        K[e:e}\OperatorTok{+}\DecValTok{2}\NormalTok{, e:e}\OperatorTok{+}\DecValTok{2}\NormalTok{] }\OperatorTok{+=}\NormalTok{ k\_e}
\NormalTok{        F\_vec[e:e}\OperatorTok{+}\DecValTok{2}\NormalTok{] }\OperatorTok{+=}\NormalTok{ f\_e}
\NormalTok{    a }\OperatorTok{=}\NormalTok{ np.linalg.solve(K[}\DecValTok{1}\NormalTok{:}\OperatorTok{{-}}\DecValTok{1}\NormalTok{, }\DecValTok{1}\NormalTok{:}\OperatorTok{{-}}\DecValTok{1}\NormalTok{], F\_vec[}\DecValTok{1}\NormalTok{:}\OperatorTok{{-}}\DecValTok{1}\NormalTok{])}
\NormalTok{    y\_fem }\OperatorTok{=}\NormalTok{ np.zeros(n\_nodes)}
\NormalTok{    y\_fem[}\DecValTok{1}\NormalTok{:}\OperatorTok{{-}}\DecValTok{1}\NormalTok{] }\OperatorTok{=}\NormalTok{ a}
\NormalTok{    x\_fem }\OperatorTok{=}\NormalTok{ np.linspace(}\DecValTok{0}\NormalTok{, }\DecValTok{1}\NormalTok{, n\_nodes)}

    \CommentTok{\# {-}{-}{-} Plot on the subplot {-}{-}{-}}
\NormalTok{    ax.plot(x\_exact, y\_exact, }\StringTok{\textquotesingle{}b{-}\textquotesingle{}}\NormalTok{, label}\OperatorTok{=}\StringTok{\textquotesingle{}Exact: y = 0.5x(1{-}x)\textquotesingle{}}\NormalTok{)}
\NormalTok{    ax.plot(x\_fem, y\_fem, }\StringTok{\textquotesingle{}ro{-}\textquotesingle{}}\NormalTok{, label}\OperatorTok{=}\SpecialStringTok{f\textquotesingle{}FEM (}\SpecialCharTok{\{}\NormalTok{n\_elem}\SpecialCharTok{\}}\SpecialStringTok{ elements)\textquotesingle{}}\NormalTok{)}
\NormalTok{    ax.set\_title(}\SpecialStringTok{f\textquotesingle{}}\SpecialCharTok{\{}\NormalTok{n\_elem}\SpecialCharTok{\}}\SpecialStringTok{ elements\textquotesingle{}}\NormalTok{)}
\NormalTok{    ax.set\_xlabel(}\StringTok{\textquotesingle{}x\textquotesingle{}}\NormalTok{)}
\NormalTok{    ax.set\_ylabel(}\StringTok{\textquotesingle{}y\textquotesingle{}}\NormalTok{)}
\NormalTok{    ax.grid(}\VariableTok{True}\NormalTok{)}
\NormalTok{    ax.legend(fontsize}\OperatorTok{=}\DecValTok{8}\NormalTok{)}

\NormalTok{plt.tight\_layout()}
\NormalTok{plt.show()}
\end{Highlighting}
\end{Shaded}

\begin{figure}[htbp]

{\centering \pandocbounded{\includegraphics[keepaspectratio]{ch01-introduction_files/figure-pdf/cell-8-output-1.png}}

}

\caption{FEM solution with varying mesh refinement}

\end{figure}%

\chapter{Deriving the Algebraic System and Galerkin's
Method}\label{deriving-the-algebraic-system-and-galerkins-method}

\section{Setup}\label{setup-1}

\section{Review: The Weak Form}\label{review-the-weak-form}

\begin{Shaded}
\begin{Highlighting}[]
\InformationTok{\textasciigrave{}\textasciigrave{}\textasciigrave{}\{python\}}
\InformationTok{\#| code{-}fold: true}
\InformationTok{print\_styled("Assuming constant A, B, C and Dirichlet BCs v(0)=v(L)=0:")}
\InformationTok{weak\_term1 = {-}Integral(A * diff(v, x) * diff(y, x), (x, 0, L))}
\InformationTok{weak\_term2 = Integral(v * B * diff(y, x), (x, 0, L))}
\InformationTok{weak\_term3 = Integral(v * C * y, (x, 0, L))}
\InformationTok{rhs\_term = Integral(v * F, (x, 0, L))}
\InformationTok{weak\_form\_eq = Eq(weak\_term1 + weak\_term2 + weak\_term3, rhs\_term)}
\InformationTok{print\_styled("The weak form is:")}
\InformationTok{Lprint(weak\_form\_eq)}
\InformationTok{\textasciigrave{}\textasciigrave{}\textasciigrave{}}
\end{Highlighting}
\end{Shaded}

Assuming constant A, B, C and Dirichlet BCs v(0)=v(L)=0:

The weak form is:

$\displaystyle - \int\limits_{0}^{L} A \frac{d}{d x} v{\left(x \right)} \frac{d}{d x} y{\left(x \right)}\, dx + \int\limits_{0}^{L} B v{\left(x \right)} \frac{d}{d x} y{\left(x \right)}\, dx + \int\limits_{0}^{L} C v{\left(x \right)} y{\left(x \right)}\, dx = \int\limits_{0}^{L} F{\left(x \right)} v{\left(x \right)}\, dx$

\section{Discretization Using Shape
Functions}\label{discretization-using-shape-functions}

\begin{Shaded}
\begin{Highlighting}[]
\InformationTok{\textasciigrave{}\textasciigrave{}\textasciigrave{}\{python\}}
\InformationTok{\#| code{-}fold: true}
\InformationTok{print\_styled("The weak form is:")}
\InformationTok{Lprint(weak\_form\_eq)}
\InformationTok{print\_styled("Substitute Discretized Approximations")}
\InformationTok{print\_styled("Approximate solution $y(x)$ and test function $v(x)$ using shape functions $\textbackslash{}phi\_k(x)$:")}

\InformationTok{y\_approx = Sum(a[j] * phi(j,x), (j, 1, N))}
\InformationTok{v\_approx = Sum(b[i] * phi(i,x), (i, 1, N))}

\InformationTok{print\_styled("Approximate solution $y(x)$:")}
\InformationTok{Lprint(Eq(y, y\_approx, evaluate=False))}

\InformationTok{print\_styled("Test function $v(x)$:")}
\InformationTok{Lprint(Eq(v, v\_approx, evaluate=False))}

\InformationTok{print\_styled("Derivatives:")}
\InformationTok{dy\_dx\_approx = Derivative(y\_approx, x)}
\InformationTok{dv\_dx\_approx = Derivative(v\_approx, x)}
\InformationTok{Lprint(Eq(Derivative(y,x), dy\_dx\_approx, evaluate=False))}
\InformationTok{Lprint(Eq(Derivative(v,x), dv\_dx\_approx, evaluate=False))}

\InformationTok{print\_styled("Substitute sums into the weak form terms:")}
\InformationTok{term1\_subst\_sym = weak\_term1.subs(\{v: v\_approx, y: y\_approx\})}
\InformationTok{term2\_subst\_sym = weak\_term2.subs(\{v: v\_approx, y: y\_approx\})}
\InformationTok{term3\_subst\_sym = weak\_term3.subs(\{v: v\_approx, y: y\_approx\})}
\InformationTok{rhs\_subst\_sym = rhs\_term.subs(\{v: v\_approx\})}

\InformationTok{print\_styled("Term 1 (${-}\textbackslash{}int A v\textquotesingle{} y\textquotesingle{} dx$):")}
\InformationTok{Lprint(term1\_subst\_sym)}
\InformationTok{print\_styled("Term 2 ($\textbackslash{}int B v y\textquotesingle{} dx$):")}
\InformationTok{Lprint(term2\_subst\_sym)}
\InformationTok{print\_styled("Term 3 ($\textbackslash{}int C v y dx$):")}
\InformationTok{Lprint(term3\_subst\_sym)}
\InformationTok{print\_styled("RHS Term ($\textbackslash{}int v F dx$):")}
\InformationTok{Lprint(rhs\_subst\_sym)}
\InformationTok{\textasciigrave{}\textasciigrave{}\textasciigrave{}}
\end{Highlighting}
\end{Shaded}

The weak form is:

$\displaystyle - \int\limits_{0}^{L} A \frac{d}{d x} v{\left(x \right)} \frac{d}{d x} y{\left(x \right)}\, dx + \int\limits_{0}^{L} B v{\left(x \right)} \frac{d}{d x} y{\left(x \right)}\, dx + \int\limits_{0}^{L} C v{\left(x \right)} y{\left(x \right)}\, dx = \int\limits_{0}^{L} F{\left(x \right)} v{\left(x \right)}\, dx$

Substitute Discretized Approximations

Approximate solution \(y(x)\) and test function \(v(x)\) using shape
functions \(\phi_k(x)\):

Approximate solution \(y(x)\):

$\displaystyle y{\left(x \right)} = \sum_{j=1}^{N} {a}_{j} \phi_{j}(x)$

Test function \(v(x)\):

$\displaystyle v{\left(x \right)} = \sum_{i=1}^{N} {b}_{i} \phi_{i}(x)$

Derivatives:

$\displaystyle \frac{d}{d x} y{\left(x \right)} = \frac{\partial}{\partial x} \sum_{j=1}^{N} {a}_{j} \phi_{j}(x)$

$\displaystyle \frac{d}{d x} v{\left(x \right)} = \frac{\partial}{\partial x} \sum_{i=1}^{N} {b}_{i} \phi_{i}(x)$

Substitute sums into the weak form terms:

Term 1 (\(-\int A v' y' dx\)):

$\displaystyle - \int\limits_{0}^{L} A \left(\frac{\partial}{\partial x} \sum_{j=1}^{N} {a}_{j} \phi_{j}(x)\right) \frac{\partial}{\partial x} \sum_{i=1}^{N} {b}_{i} \phi_{i}(x)\, dx$

Term 2 (\(\int B v y' dx\)):

$\displaystyle \int\limits_{0}^{L} B \left(\frac{\partial}{\partial x} \sum_{j=1}^{N} {a}_{j} \phi_{j}(x)\right) \sum_{i=1}^{N} {b}_{i} \phi_{i}(x)\, dx$

Term 3 (\(\int C v y dx\)):

$\displaystyle \int\limits_{0}^{L} C \left(\sum_{j=1}^{N} {a}_{j} \phi_{j}(x)\right) \sum_{i=1}^{N} {b}_{i} \phi_{i}(x)\, dx$

RHS Term (\(\int v F dx\)):

$\displaystyle \int\limits_{0}^{L} F{\left(x \right)} \sum_{i=1}^{N} {b}_{i} \phi_{i}(x)\, dx$

\section{Rearranging the Summation and
Integrals}\label{rearranging-the-summation-and-integrals}

\begin{Shaded}
\begin{Highlighting}[]
\InformationTok{\textasciigrave{}\textasciigrave{}\textasciigrave{}\{python\}}
\InformationTok{\#| code{-}fold: true}
\InformationTok{print\_styled("Assuming continuity, swap integration and summation, then factor out $b\_i$.")}
\InformationTok{print\_styled(r"The weak form $\textbackslash{}sum\_\{\textbackslash{}text\{term\}\} (\textbackslash{}text\{Term\}) = \textbackslash{}text\{RHS\}$ becomes:")}
\InformationTok{Lprint(r"\textbackslash{}sum\_\{i=1\}\^{}N b\_i \textbackslash{}left( [\textbackslash{}text\{Inner part for Term 1\}]\_i + [\textbackslash{}text\{Inner part for Term 2\}]\_i + [\textbackslash{}text\{Inner part for Term 3\}]\_i \textbackslash{}right) = \textbackslash{}sum\_\{i=1\}\^{}N b\_i [\textbackslash{}text\{Inner part for RHS\}]\_i")}

\InformationTok{print\_styled(r"Where the inner parts (coefficients of $b\_i$) are:")}

\InformationTok{term1\_inner = Sum( {-}A * Integral(diff(phi(i,x), x) * diff(phi(j,x), x), (x, 0, L)) * a[j], (j, 1, N))}
\InformationTok{term2\_inner = Sum(  B * Integral(phi(i,x) * diff(phi(j,x), x), (x, 0, L)) * a[j], (j, 1, N))}
\InformationTok{term3\_inner = Sum(  C * Integral(phi(i,x) * phi(j,x), (x, 0, L)) * a[j], (j, 1, N))}
\InformationTok{rhs\_inner = Integral(phi(i,x) * F, (x, 0, L))}

\InformationTok{print\_styled("Inner part for Term 1 (coefficient of $b\_i$):")}
\InformationTok{Lprint(term1\_inner)}
\InformationTok{print\_styled("Inner part for Term 2 (coefficient of $b\_i$):")}
\InformationTok{Lprint(term2\_inner)}
\InformationTok{print\_styled("Inner part for Term 3 (coefficient of $b\_i$):")}
\InformationTok{Lprint(term3\_inner)}
\InformationTok{print\_styled("Inner part for RHS (coefficient of $b\_i$):")}
\InformationTok{Lprint(rhs\_inner)}
\InformationTok{\textasciigrave{}\textasciigrave{}\textasciigrave{}}
\end{Highlighting}
\end{Shaded}

Assuming continuity, swap integration and summation, then factor out
\(b_i\).

The weak form \(\sum_{\text{term}} (\text{Term}) = \text{RHS}\) becomes:

$\displaystyle \sum_{i=1}^N b_i \left( [\text{Inner part for Term 1}]_i + [\text{Inner part for Term 2}]_i + [\text{Inner part for Term 3}]_i \right) = \sum_{i=1}^N b_i [\text{Inner part for RHS}]_i$

Where the inner parts (coefficients of \(b_i\)) are:

Inner part for Term 1 (coefficient of \(b_i\)):

$\displaystyle \sum_{j=1}^{N} - A {a}_{j} \int\limits_{0}^{L} \frac{d}{d x} \phi_{i}(x) \frac{d}{d x} \phi_{j}(x)\, dx$

Inner part for Term 2 (coefficient of \(b_i\)):

$\displaystyle \sum_{j=1}^{N} B {a}_{j} \int\limits_{0}^{L} \phi_{i}(x) \frac{d}{d x} \phi_{j}(x)\, dx$

Inner part for Term 3 (coefficient of \(b_i\)):

$\displaystyle \sum_{j=1}^{N} C {a}_{j} \int\limits_{0}^{L} \phi_{i}(x) \phi_{j}(x)\, dx$

Inner part for RHS (coefficient of \(b_i\)):

$\displaystyle \int\limits_{0}^{L} F{\left(x \right)} \phi_{i}(x)\, dx$

\section{Obtaining N Equations}\label{obtaining-n-equations}

\begin{Shaded}
\begin{Highlighting}[]
\InformationTok{\textasciigrave{}\textasciigrave{}\textasciigrave{}\{python\}}
\InformationTok{\#| code{-}fold: true}
\InformationTok{print\_styled("The rearranged form is:")}
\InformationTok{print\_styled("$\textbackslash{}\textbackslash{}sum\_\{i=1\}\^{}N b\_i ( \textbackslash{}\textbackslash{}text\{LHS\}\_i {-} \textbackslash{}\textbackslash{}text\{RHS\}\_i ) = 0$.")}
\InformationTok{print\_styled("Since this must hold for arbitrary $b\_i$, the term in parentheses must be zero for each $i$:")}
\InformationTok{print\_styled("  $\textbackslash{}\textbackslash{}text\{LHS\}\_i {-} \textbackslash{}\textbackslash{}text\{RHS\}\_i = 0$.")}

\InformationTok{equation\_for\_i = Eq(term1\_inner + term2\_inner + term3\_inner, rhs\_inner)}
\InformationTok{print\_styled("This gives N equations (for $i = 1, ..., N$):")}
\InformationTok{Lprint(equation\_for\_i)}
\InformationTok{\textasciigrave{}\textasciigrave{}\textasciigrave{}}
\end{Highlighting}
\end{Shaded}

The rearranged form is:

\(\sum_{i=1}^N b_i ( \text{LHS}_i - \text{RHS}_i ) = 0\).

Since this must hold for arbitrary \(b_i\), the term in parentheses must
be zero for each \(i\):

\(\text{LHS}_i - \text{RHS}_i = 0\).

This gives N equations (for \(i = 1, ..., N\)):

$\displaystyle \sum_{j=1}^{N} - A {a}_{j} \int\limits_{0}^{L} \frac{d}{d x} \phi_{i}(x) \frac{d}{d x} \phi_{j}(x)\, dx + \sum_{j=1}^{N} B {a}_{j} \int\limits_{0}^{L} \phi_{i}(x) \frac{d}{d x} \phi_{j}(x)\, dx + \sum_{j=1}^{N} C {a}_{j} \int\limits_{0}^{L} \phi_{i}(x) \phi_{j}(x)\, dx = \int\limits_{0}^{L} F{\left(x \right)} \phi_{i}(x)\, dx$

\section{Assembling the Matrix
System}\label{assembling-the-matrix-system}

\begin{Shaded}
\begin{Highlighting}[]
\InformationTok{\textasciigrave{}\textasciigrave{}\textasciigrave{}\{python\}}
\InformationTok{\#| code{-}fold: true}
\InformationTok{print\_styled("The $i${-}th equation:")}
\InformationTok{print\_styled("$\textbackslash{}\textbackslash{}sum\_\{j=1\}\^{}N [ ... ]\_\{ij\} a\_j = F\_i$")}
\InformationTok{print\_styled("leads to the matrix form:")}
\InformationTok{print\_styled("$K a = F$")}

\InformationTok{K\_ij = {-}A * Integral(diff(phi(i,x), x) * diff(phi(j,x), x), (x, 0, L)) + \textbackslash{}}
\InformationTok{        B * Integral(phi(i,x) * diff(phi(j,x), x), (x, 0, L)) + \textbackslash{}}
\InformationTok{        C * Integral(phi(i,x) * phi(j,x), (x, 0, L))}
\InformationTok{F\_i = Integral(phi(i,x) * F, (x, 0, L))}

\InformationTok{print\_styled("Where the matrix/vector components are:")}
\InformationTok{Lprint(Eq(Symbol(\textquotesingle{}K\_\{ij\}\textquotesingle{}), K\_ij))}
\InformationTok{Lprint(Eq(Symbol(\textquotesingle{}F\_i\textquotesingle{}), F\_i))}

\InformationTok{K\_sym = MatrixSymbol(\textquotesingle{}K\textquotesingle{}, N, N)}
\InformationTok{a\_sym = MatrixSymbol(\textquotesingle{}a\textquotesingle{}, N, 1)}
\InformationTok{F\_sym = MatrixSymbol(\textquotesingle{}F\textquotesingle{}, N, 1)}
\InformationTok{matrix\_eq = Eq(K\_sym * a\_sym, F\_sym, evaluate=False)}
\InformationTok{print\_styled("The matrix system:")}
\InformationTok{Lprint(matrix\_eq)}

\InformationTok{print\_styled("Writing the full system for $i=1, ..., N$ gives the matrix form $Ka=F$:")}

\InformationTok{latex\_matrix\_form = r"""}
\InformationTok{\textbackslash{}begin\{bmatrix\}}
\InformationTok{K\_\{11\} \& K\_\{12\} \& \textbackslash{}dots \& K\_\{1N\} \textbackslash{}\textbackslash{}}
\InformationTok{K\_\{21\} \& K\_\{22\} \& \textbackslash{}dots \& K\_\{2N\} \textbackslash{}\textbackslash{}}
\InformationTok{\textbackslash{}vdots \& \textbackslash{}vdots \& \textbackslash{}ddots \& \textbackslash{}vdots \textbackslash{}\textbackslash{}}
\InformationTok{K\_\{N1\} \& K\_\{N2\} \& \textbackslash{}dots \& K\_\{NN\}}
\InformationTok{\textbackslash{}end\{bmatrix\}}
\InformationTok{\textbackslash{}begin\{bmatrix\}}
\InformationTok{a\_\{1\} \textbackslash{}\textbackslash{}}
\InformationTok{a\_\{2\} \textbackslash{}\textbackslash{}}
\InformationTok{\textbackslash{}vdots \textbackslash{}\textbackslash{}}
\InformationTok{a\_\{N\}}
\InformationTok{\textbackslash{}end\{bmatrix\}}
\InformationTok{=}
\InformationTok{\textbackslash{}begin\{bmatrix\}}
\InformationTok{F\_\{1\} \textbackslash{}\textbackslash{}}
\InformationTok{F\_\{2\} \textbackslash{}\textbackslash{}}
\InformationTok{\textbackslash{}vdots \textbackslash{}\textbackslash{}}
\InformationTok{F\_\{N\}}
\InformationTok{\textbackslash{}end\{bmatrix\}}
\InformationTok{"""}
\InformationTok{Lprint(latex\_matrix\_form)}

\InformationTok{print\_styled("This linear system $Ka = F$ can be solved for the unknown coefficients vector $a = [a\_1, a\_2, ..., a\_N]\^{}T$.")}
\InformationTok{\textasciigrave{}\textasciigrave{}\textasciigrave{}}
\end{Highlighting}
\end{Shaded}

The \(i\)-th equation:

\(\sum_{j=1}^N [ ... ]_{ij} a_j = F_i\)

leads to the matrix form:

\(K a = F\)

Where the matrix/vector components are:

$\displaystyle K_{ij} = - A \int\limits_{0}^{L} \frac{d}{d x} \phi_{i}(x) \frac{d}{d x} \phi_{j}(x)\, dx + B \int\limits_{0}^{L} \phi_{i}(x) \frac{d}{d x} \phi_{j}(x)\, dx + C \int\limits_{0}^{L} \phi_{i}(x) \phi_{j}(x)\, dx$

$\displaystyle F_{i} = \int\limits_{0}^{L} F{\left(x \right)} \phi_{i}(x)\, dx$

The matrix system:

$\displaystyle K a = F$

Writing the full system for \(i=1, ..., N\) gives the matrix form
\(Ka=F\):

$\displaystyle 
\begin{bmatrix}
K_{11} & K_{12} & \dots & K_{1N} \\
K_{21} & K_{22} & \dots & K_{2N} \\
\vdots & \vdots & \ddots & \vdots \\
K_{N1} & K_{N2} & \dots & K_{NN}
\end{bmatrix}
\begin{bmatrix}
a_{1} \\
a_{2} \\
\vdots \\
a_{N}
\end{bmatrix}
=
\begin{bmatrix}
F_{1} \\
F_{2} \\
\vdots \\
F_{N}
\end{bmatrix}
$

This linear system \(Ka = F\) can be solved for the unknown coefficients
vector \(a = [a_1, a_2, ..., a_N]^T\).

\section{Note on Boundary Conditions}\label{note-on-boundary-conditions}

\begin{Shaded}
\begin{Highlighting}[]
\InformationTok{\textasciigrave{}\textasciigrave{}\textasciigrave{}\{python\}}
\InformationTok{\#| code{-}fold: true}
\InformationTok{print\_styled("We assumed $y(0)=0$ and $y(L)=0$. How are these applied to the matrix system?")}
\InformationTok{print\_styled("If using nodal basis functions where:")}
\InformationTok{print\_styled("$\textbackslash{}phi\_k(x\_\{node\}) = \textbackslash{}delta\_\{knode\}$")}
\InformationTok{print\_styled("then:")}
\InformationTok{print\_styled("$y(x\_1) = y(0) = \textbackslash{}sum a\_j \textbackslash{}phi\_j(x\_1) = a\_1$.")}
\InformationTok{print\_styled("Similarly")}
\InformationTok{print\_styled("$y(x\_N) = y(L) = a\_N$.")}
\InformationTok{print\_styled("So, for $y(0)=y(L)=0$, we know $a\_1 = 0$ and $a\_N = 0$.")}
\InformationTok{print\_styled("These known values must be incorporated when solving $Ka=F$. Common methods include:")}
\InformationTok{print\_styled("  {-} Modifying the K matrix and F vector (e.g., setting rows/columns to enforce the condition).")}
\InformationTok{print\_styled("  {-} Partitioning the system into known and unknown degrees of freedom.")}
\InformationTok{print\_styled("Handling boundary conditions precisely is a key part of FEM implementation and will be detailed later.")}
\InformationTok{\textasciigrave{}\textasciigrave{}\textasciigrave{}}
\end{Highlighting}
\end{Shaded}

We assumed \(y(0)=0\) and \(y(L)=0\). How are these applied to the
matrix system?

If using nodal basis functions where:

\(\phi_k(x_{node}) = \delta_{knode}\)

then:

\(y(x_1) = y(0) = \sum a_j \phi_j(x_1) = a_1\).

Similarly

\(y(x_N) = y(L) = a_N\).

So, for \(y(0)=y(L)=0\), we know \(a_1 = 0\) and \(a_N = 0\).

These known values must be incorporated when solving \(Ka=F\). Common
methods include:

\begin{itemize}
\tightlist
\item
  Modifying the K matrix and F vector (e.g., setting rows/columns to
  enforce the condition).
\end{itemize}

\begin{itemize}
\tightlist
\item
  Partitioning the system into known and unknown degrees of freedom.
\end{itemize}

Handling boundary conditions precisely is a key part of FEM
implementation and will be detailed later.

\section{Weighted Residuals and Galerkin's
Method}\label{weighted-residuals-and-galerkins-method}

\begin{Shaded}
\begin{Highlighting}[]
\InformationTok{\textasciigrave{}\textasciigrave{}\textasciigrave{}\{python\}}
\InformationTok{\#| code{-}fold: true}
\InformationTok{print\_styled("Let\textquotesingle{}s revisit the core idea from a different perspective.")}
\InformationTok{print\_styled("Our original differential equation is $A(y) = F$, where $A$ is the differential operator.")}
\InformationTok{def differential\_operator(y\_var):}
\InformationTok{    return A * diff(y\_var, x, 2) + B * diff(y\_var, x) + C * y\_var}
\InformationTok{eq = sp.Eq(differential\_operator(y), F)}
\InformationTok{Lprint(eq)}
\InformationTok{print\_styled("Next, we represent $y(x)$ as a linear combination of basis functions $\textbackslash{}phi\_j(x)$, resulting in an approximated function $y\^{}N(x) \textbackslash{}\textbackslash{}approx y(x)$:")}

\InformationTok{yn = Function(\textquotesingle{}y\^{}N\textquotesingle{})(x)}
\InformationTok{Lprint(Eq(yn, y\_approx, evaluate=False))}

\InformationTok{print\_styled("Now, we can define the residual $r\^{}N(x)$ for the approximate solution $y\^{}N(x)$ given by $A(y\^{}N){-}F=r\^{}N$:")}
\InformationTok{A\_op = lambda func: A * diff(func, x, 2) + B * diff(func, x) + C * func}
\InformationTok{r\_N = Symbol(\textquotesingle{}r\^{}N\textquotesingle{})}
\InformationTok{residual\_def = Eq(r\_N, A\_op(y\_approx) {-} F, evaluate=False)}
\InformationTok{Lprint(residual\_def)}
\InformationTok{print\_styled("The residual $r\^{}N(x)$ is the error in the approximation. The goal of FEM is to make this residual \textquotesingle{}small\textquotesingle{} in some sense.")}
\InformationTok{n = sp.Symbol(\textquotesingle{}n\textquotesingle{}, integer=True, positive=True)}
\InformationTok{PI = sp.Symbol(\textquotesingle{}PI\textquotesingle{}, positive=True)}
\InformationTok{print\_styled("Since $r\^{}N$ is a function on $[0,L]$, we can measure its \textquotesingle{}size\textquotesingle{} using:")}
\InformationTok{Lprint(Eq(PI, Integral(r\_N**2, (x, 0, L)), evaluate=False))}
\InformationTok{print\_styled("A possible approach is to minimize $\textbackslash{}\textbackslash{}Pi$ by taking the derivative with respect to the coefficients $a\_j$ and setting it to zero.")}
\InformationTok{\textasciigrave{}\textasciigrave{}\textasciigrave{}}
\end{Highlighting}
\end{Shaded}

Let's revisit the core idea from a different perspective.

Our original differential equation is \(A(y) = F\), where \(A\) is the
differential operator.

$\displaystyle A \frac{d^{2}}{d x^{2}} y{\left(x \right)} + B \frac{d}{d x} y{\left(x \right)} + C y{\left(x \right)} = F{\left(x \right)}$

Next, we represent \(y(x)\) as a linear combination of basis functions
\(\phi_j(x)\), resulting in an approximated function
\(y^N(x) \approx y(x)\):

$\displaystyle y^{N}{\left(x \right)} = \sum_{j=1}^{N} {a}_{j} \phi_{j}(x)$

Now, we can define the residual \(r^N(x)\) for the approximate solution
\(y^N(x)\) given by \(A(y^N)-F=r^N\):

$\displaystyle r^{N} = A \sum_{j=1}^{N} {a}_{j} \frac{d^{2}}{d x^{2}} \phi_{j}(x) + B \sum_{j=1}^{N} {a}_{j} \frac{d}{d x} \phi_{j}(x) + C \sum_{j=1}^{N} {a}_{j} \phi_{j}(x) - F{\left(x \right)}$

The residual \(r^N(x)\) is the error in the approximation. The goal of
FEM is to make this residual `small' in some sense.

Since \(r^N\) is a function on \([0,L]\), we can measure its `size'
using:

$\displaystyle \pi = \int\limits_{0}^{L} \left(r^{N}\right)^{2}\, dx$

A possible approach is to minimize \(\Pi\) by taking the derivative with
respect to the coefficients \(a_j\) and setting it to zero.

\subsection{The Galerkin Choice}\label{the-galerkin-choice}

\begin{Shaded}
\begin{Highlighting}[]
\InformationTok{\textasciigrave{}\textasciigrave{}\textasciigrave{}\{python\}}
\InformationTok{\#| code{-}fold: true}
\InformationTok{print\_styled("In the Galerkin approach, we define weighting functions $w\_i(x)$.")}
\InformationTok{w = Function(\textquotesingle{}w\textquotesingle{})}
\InformationTok{weighted\_residual\_eq = Eq(Integral(r\_N * w(i, x), (x, 0, L)), 0)}
\InformationTok{Lprint(weighted\_residual\_eq)}
\InformationTok{print\_styled("The **Method of Weighted Residuals** requires the residual to be orthogonal to the weighting functions $w\_i(x)$.")}
\InformationTok{print\_styled("The **Galerkin Method** makes a specific choice: use the *same* functions for weighting as for the basis/shape functions.")}
\InformationTok{print\_styled("This means we set $w\_i(x) = \textbackslash{}\textbackslash{}phi\_i(x)$:")}
\InformationTok{galerkin\_eq\_from\_WR = weighted\_residual\_eq.subs(\{w(i,x): phi(i,x)\})}
\InformationTok{Lprint(galerkin\_eq\_from\_WR)}
\InformationTok{print\_styled(r"This must hold for $i = 1, 2, ..., N$.")}

\InformationTok{print\_styled("Substituting the definition of $r\^{}N$ into the Galerkin equation gives:")}
\InformationTok{galerkin\_expanded = galerkin\_eq\_from\_WR.subs(\{r\_N: A\_op(y\_approx) {-} F\})}
\InformationTok{Lprint(galerkin\_expanded)}

\InformationTok{print\_styled("Applying integration by parts to the second derivative term and requiring that $\textbackslash{}\textbackslash{}phi\_i$ satisfy the boundary conditions, we obtain:")}
\InformationTok{print\_styled(r"$\textbackslash{}int\_0\^{}L \textbackslash{}left( {-}A \textbackslash{}frac\{d\textbackslash{}phi\_i\}\{dx\} \textbackslash{}sum\_\{j=1\}\^{}N\textbackslash{}frac\{d\textbackslash{}phi\_j\}\{dx\} + B \textbackslash{}phi\_i \textbackslash{}sum\_\{j=1\}\^{}N\textbackslash{}frac\{d\textbackslash{}phi\_j\}\{dx\} + C \textbackslash{}phi\_i \textbackslash{}sum\_\{j=1\}\^{}N\textbackslash{}phi\_j {-} \textbackslash{}phi\_i F \textbackslash{}right) dx = 0$")}

\InformationTok{print\_styled("This is the *exact same* $i${-}th equation we derived earlier from the weak form.")}
\InformationTok{print\_styled("Therefore, the Galerkin method provides a powerful theoretical justification: it systematically minimizes the residual by making it orthogonal to the basis functions used for the approximation.")}
\InformationTok{\textasciigrave{}\textasciigrave{}\textasciigrave{}}
\end{Highlighting}
\end{Shaded}

In the Galerkin approach, we define weighting functions \(w_i(x)\).

$\displaystyle \int\limits_{0}^{L} r^{N} w{\left(i,x \right)}\, dx = 0$

The \textbf{Method of Weighted Residuals} requires the residual to be
orthogonal to the weighting functions \(w_i(x)\).

The \textbf{Galerkin Method} makes a specific choice: use the
\emph{same} functions for weighting as for the basis/shape functions.

This means we set \(w_i(x) = \phi_i(x)\):

$\displaystyle \int\limits_{0}^{L} r^{N} \phi_{i}(x)\, dx = 0$

This must hold for \(i = 1, 2, ..., N\).

Substituting the definition of \(r^N\) into the Galerkin equation gives:

$\displaystyle \int\limits_{0}^{L} \left(A \sum_{j=1}^{N} {a}_{j} \frac{d^{2}}{d x^{2}} \phi_{j}(x) + B \sum_{j=1}^{N} {a}_{j} \frac{d}{d x} \phi_{j}(x) + C \sum_{j=1}^{N} {a}_{j} \phi_{j}(x) - F{\left(x \right)}\right) \phi_{i}(x)\, dx = 0$

Applying integration by parts to the second derivative term and
requiring that \(\phi_i\) satisfy the boundary conditions, we obtain:

\(\int_0^L \left( -A \frac{d\phi_i}{dx} \sum_{j=1}^N\frac{d\phi_j}{dx} + B \phi_i \sum_{j=1}^N\frac{d\phi_j}{dx} + C \phi_i \sum_{j=1}^N\phi_j - \phi_i F \right) dx = 0\)

This is the \emph{exact same} \(i\)-th equation we derived earlier from
the weak form.

Therefore, the Galerkin method provides a powerful theoretical
justification: it systematically minimizes the residual by making it
orthogonal to the basis functions used for the approximation.

\chapter{Function Spaces}\label{function-spaces}

Hilbert, L², and Sobolev Spaces

\hfill\break

\section{Setup}\label{setup-2}

\section{Motivation: Why Special Function
Spaces?}\label{motivation-why-special-function-spaces}

In FEM, we transform strong forms of PDEs (like
\(\nabla\cdot(k\nabla u)=f\)) into \textbf{weak forms} using integration
by parts. Weak forms look like: Find \(u\) such that

\[
\int_{\Omega}(k\nabla u \cdot \nabla v\,d\Omega) = \int_{\Omega}fv \,d\Omega
\]

for all suitable ``test functions'' \(v\).

This involves integrals of functions and their derivatives. \textbf{What
kind of functions \(u\) and \(v\) can we ``legally'' plug into these
integrals?} Do they need to be smooth? Can they have kinks?

\section{From Vectors to Functions}\label{from-vectors-to-functions}

You know \textbf{vector spaces} like \(\mathbb{R}^3\): vectors,
addition, scalar multiplication. We measure vector ``size'' (norm) using
the dot product:

\[
\|x\|^2 = x_1^2 + x_2^2 + x_3^2 = x \cdot x, \qquad \|x\| = \sqrt{x \cdot x}
\]

Functions can also form vector spaces! We can add functions (\(f+g\))
and scale them (\(cf\)). The challenge is defining a ``dot product'' and
``size'' (norm) for functions, and ensuring our space is ``complete''
(no missing limits).

\section{Hilbert Spaces: The General
Framework}\label{hilbert-spaces-the-general-framework}

A \textbf{Hilbert Space} is a vector space (often of functions) with an
\textbf{inner product} \(\langle f,g \rangle\). The inner product lets
us define:

\begin{itemize}
\tightlist
\item
  \textbf{Norm (Size):} \(\|f\| = \sqrt{\langle f,f \rangle}\)
\item
  \textbf{Orthogonality:} \(f\) is orthogonal to \(g\) if
  \(\langle f,g \rangle = 0\)
\end{itemize}

Crucially, Hilbert spaces are \textbf{complete}: every Cauchy sequence
converges to a point \emph{within} the space. This is vital for proving
solution existence and convergence of approximations like FEM.

\section{\texorpdfstring{\(L^2(\Omega)\): Square-Integrable
Functions}{L\^{}2(\textbackslash Omega): Square-Integrable Functions}}\label{l2omega-square-integrable-functions}

The most common Hilbert space in introductory FEM.

\(L^2(\Omega)\) is the space of functions \(f\) defined on a domain
\(\Omega\) whose square has a finite integral:

\[\int_{\Omega} |f(x)|^2 \,dx < \infty\]

Think: functions with finite ``energy'' or finite mean-square value.

\textbf{Inner Product:}
\[\langle f, g\rangle_{L^2} = \int_{\Omega} f(x) g(x) \,dx\]

\textbf{Norm (\(L^2\) Norm):}
\[\|f\|_{L^2} = \sqrt{\int_{\Omega} |f(x)|^2 \,dx}\]

\(L^2(\Omega)\) is a Hilbert space, also sometimes called
\(H^0(\Omega)\).

\subsection{Examples}\label{examples}

\begin{Shaded}
\begin{Highlighting}[]
\InformationTok{\textasciigrave{}\textasciigrave{}\textasciigrave{}\{python\}}
\InformationTok{\#| code{-}fold: true}
\InformationTok{\#| fig{-}cap: "Left: sin(x) is in L²[0, π]. Right: 1/x is NOT in L²[0, π] as its square diverges near x=0."}
\InformationTok{\#| output: true}
\InformationTok{x1 = np.linspace(0.01, np.pi, 400)}
\InformationTok{x2 = np.linspace(0, np.pi, 400)}

\InformationTok{f1 = np.sin(x2)}
\InformationTok{f2 = 1 / x1}

\InformationTok{fig, axs = plt.subplots(1, 2, figsize=(10, 4))}

\InformationTok{axs[0].plot(x2, f1, label=\textquotesingle{}f(x) = sin(x)\textquotesingle{})}
\InformationTok{axs[0].fill\_between(x2, f1**2, alpha=0.3, label=\textquotesingle{}f(x)²\textquotesingle{})}
\InformationTok{axs[0].set\_title(\textquotesingle{}Function in L²[0, π]\textquotesingle{})}
\InformationTok{axs[0].set\_xlabel(\textquotesingle{}x\textquotesingle{})}
\InformationTok{axs[0].set\_ylabel(\textquotesingle{}f(x)\textquotesingle{})}
\InformationTok{axs[0].legend()}
\InformationTok{axs[0].grid(True)}
\InformationTok{axs[0].set\_ylim(bottom=0)}

\InformationTok{axs[1].plot(x1, f2, label=\textquotesingle{}g(x) = 1/x\textquotesingle{})}
\InformationTok{axs[1].fill\_between(x1, f2**2, alpha=0.3, label=\textquotesingle{}g(x)²\textquotesingle{})}
\InformationTok{axs[1].set\_title(\textquotesingle{}Function NOT in L²[0, π]\textquotesingle{})}
\InformationTok{axs[1].set\_xlabel(\textquotesingle{}x\textquotesingle{})}
\InformationTok{axs[1].set\_ylabel(\textquotesingle{}g(x)\textquotesingle{})}
\InformationTok{axs[1].legend()}
\InformationTok{axs[1].grid(True)}
\InformationTok{axs[1].set\_ylim(0, 50)}

\InformationTok{plt.tight\_layout()}
\InformationTok{plt.show()}
\InformationTok{\textasciigrave{}\textasciigrave{}\textasciigrave{}}
\end{Highlighting}
\end{Shaded}

\begin{figure}[htbp]

{\centering \pandocbounded{\includegraphics[keepaspectratio]{ch03-function-spaces_files/figure-pdf/cell-3-output-1.png}}

}

\caption{Left: sin(x) is in L²{[}0, π{]}. Right: 1/x is NOT in L²{[}0,
π{]} as its square diverges near x=0.}

\end{figure}%

\begin{tcolorbox}[enhanced jigsaw, opacitybacktitle=0.6, leftrule=.75mm, colframe=quarto-callout-note-color-frame, toptitle=1mm, breakable, bottomtitle=1mm, opacityback=0, colback=white, rightrule=.15mm, colbacktitle=quarto-callout-note-color!10!white, arc=.35mm, left=2mm, toprule=.15mm, bottomrule=.15mm, titlerule=0mm, coltitle=black, title=\textcolor{quarto-callout-note-color}{\faInfo}\hspace{0.5em}{Note}]

Functions can be unbounded (like \(1/\sqrt{x}\) near 0 on \([0,1]\)) and
still be in \(L^2\). Functions can also be discontinuous and be in
\(L^2\).

\end{tcolorbox}

\section{Need for Derivatives: Sobolev
Spaces}\label{need-for-derivatives-sobolev-spaces}

Our weak form involves derivatives \(\nabla u \cdot \nabla v\).
Requiring functions \(u, v\) to be only in \(L^2(\Omega)\) is not enough
--- we need their derivatives to also be square-integrable.

\begin{tcolorbox}[enhanced jigsaw, opacitybacktitle=0.6, leftrule=.75mm, colframe=quarto-callout-note-color-frame, toptitle=1mm, breakable, bottomtitle=1mm, opacityback=0, colback=white, rightrule=.15mm, colbacktitle=quarto-callout-note-color!10!white, arc=.35mm, left=2mm, toprule=.15mm, bottomrule=.15mm, titlerule=0mm, coltitle=black, title=\textcolor{quarto-callout-note-color}{\faInfo}\hspace{0.5em}{Problem}]

What if \(u\) is continuous but has a kink (like \(|x|\))? Its classical
derivative isn't defined everywhere or might not be in \(L^2\).

\end{tcolorbox}

\section{\texorpdfstring{Sobolev Space
\(H^1(\Omega)\)}{Sobolev Space H\^{}1(\textbackslash Omega)}}\label{sobolev-space-h1omega}

\(H^1(\Omega)\) is the simplest Sobolev space relevant to many FEM
problems.

\begin{tcolorbox}[enhanced jigsaw, opacitybacktitle=0.6, leftrule=.75mm, colframe=quarto-callout-note-color-frame, toptitle=1mm, breakable, bottomtitle=1mm, opacityback=0, colback=white, rightrule=.15mm, colbacktitle=quarto-callout-note-color!10!white, arc=.35mm, left=2mm, toprule=.15mm, bottomrule=.15mm, titlerule=0mm, coltitle=black, title=\textcolor{quarto-callout-note-color}{\faInfo}\hspace{0.5em}{Definition}]

\(H^1(\Omega)\) is the space of functions \(f\) such that both \(f\) and
its first weak derivatives are in \(L^2(\Omega)\):

\[H^1(\Omega) = \{ f \in L^2(\Omega) \mid \nabla f \in L^2(\Omega) \}\]

A \textbf{weak derivative} is a generalization of the derivative using
integration by parts. It allows functions with ``corners'' (like
\(|x|\)) to still have a derivative within this framework, provided it's
an \(L^2\) function.

\(H^1(\Omega)\) is also a Hilbert space.

\end{tcolorbox}

\subsection{\texorpdfstring{\(H^1\) Inner Product and
Norm}{H\^{}1 Inner Product and Norm}}\label{h1-inner-product-and-norm}

\textbf{Inner Product:} Includes terms for both the function and its
derivative:

\[\langle f, g\rangle_{H^1} = \int_{\Omega} f(x) g(x) \,dx + \int_{\Omega} \nabla f(x) \cdot \nabla g(x) \,dx\]

\textbf{Norm (\(H^1\) Norm):}

\[\|f\|_{H^1}^2 = \|f\|_{L^2}^2 + \|\nabla f\|_{L^2}^2\]

If \(\|f\|_{H^1}\) is finite, the function is ``well-behaved enough''
for typical second-order PDE weak forms.

\subsection{\texorpdfstring{\(H^1\) vs \(L^2\):
Examples}{H\^{}1 vs L\^{}2: Examples}}\label{h1-vs-l2-examples}

\begin{Shaded}
\begin{Highlighting}[]
\InformationTok{\textasciigrave{}\textasciigrave{}\textasciigrave{}\{python\}}
\InformationTok{\#| code{-}fold: true}
\InformationTok{\#| fig{-}cap: "Left: x² is smooth, both function and derivative (2x) are in L², so x² ∈ H¹. Right: Step function is in L², but its derivative (a Dirac delta) is not in L², so Step ∉ H¹."}
\InformationTok{\#| output: true}
\InformationTok{x = np.linspace({-}1, 1, 400)}

\InformationTok{f1 = x**2}
\InformationTok{f1\_deriv = 2*x}
\InformationTok{f2 = np.sign(x)}

\InformationTok{fig, axs = plt.subplots(1, 2, figsize=(10, 4))}

\InformationTok{axs[0].plot(x, f1, label=\textquotesingle{}f(x) = x²\textquotesingle{})}
\InformationTok{axs[0].plot(x, f1\_deriv, label="f\textquotesingle{}(x) = 2x", linestyle=\textquotesingle{}{-}{-}\textquotesingle{})}
\InformationTok{axs[0].set\_title(\textquotesingle{}Function in H¹[{-}1, 1]\textquotesingle{})}
\InformationTok{axs[0].set\_xlabel(\textquotesingle{}x\textquotesingle{})}
\InformationTok{axs[0].set\_ylabel(\textquotesingle{}Value\textquotesingle{})}
\InformationTok{axs[0].legend()}
\InformationTok{axs[0].grid(True)}

\InformationTok{axs[1].plot(x, f2, label=\textquotesingle{}g(x) = sign(x) (Step)\textquotesingle{})}
\InformationTok{axs[1].set\_title(\textquotesingle{}Function in L²[{-}1, 1], but NOT H¹[{-}1, 1]\textquotesingle{})}
\InformationTok{axs[1].set\_xlabel(\textquotesingle{}x\textquotesingle{})}
\InformationTok{axs[1].set\_ylabel(\textquotesingle{}Value\textquotesingle{})}
\InformationTok{axs[1].legend()}
\InformationTok{axs[1].grid(True)}

\InformationTok{plt.tight\_layout()}
\InformationTok{plt.show()}
\InformationTok{\textasciigrave{}\textasciigrave{}\textasciigrave{}}
\end{Highlighting}
\end{Shaded}

\begin{figure}[htbp]

{\centering \pandocbounded{\includegraphics[keepaspectratio]{ch03-function-spaces_files/figure-pdf/cell-4-output-1.png}}

}

\caption{Left: x² is smooth, both function and derivative (2x) are in
L², so x² ∈ H¹. Right: Step function is in L², but its derivative (a
Dirac delta) is not in L², so Step ∉ H¹.}

\end{figure}%

Key observations:

\begin{itemize}
\tightlist
\item
  Smooth functions (like polynomials) are typically in \(H^1\)
\item
  Functions with ``jumps'' (like the step function) are often in \(L^2\)
  but not \(H^1\)
\item
  Functions with ``kinks'' (like \(|x|\)) can be in \(H^1\) because
  their weak derivative \(\text{sign}(x)\) exists and is in \(L^2\)
\end{itemize}

\section{Summary}\label{summary}

\begin{itemize}
\tightlist
\item
  \textbf{Weak Formulations}: FEM relies on integral (weak) forms of
  PDEs
\item
  \textbf{\(L^2\) (\(H^0\))}: Guarantees that integrals of functions
  themselves (\(\int fv\,dx\)) make sense. Useful for source terms, mass
  matrices
\item
  \textbf{\(H^1\)}: Guarantees that integrals involving first
  derivatives (\(\int \nabla f \cdot \nabla v\,dx\)) also make sense.
  The natural solution space for many second-order problems (heat,
  elasticity, potential flow). Ensures stiffness matrices are
  well-defined
\item
  Using these spaces provides a rigorous mathematical foundation for
  FEM, ensuring solutions exist and approximations converge
\end{itemize}

\section{\texorpdfstring{Beyond
\(H^1\)}{Beyond H\^{}1}}\label{beyond-h1}

Higher-order Sobolev spaces exist, like \(H^k(\Omega)\), which include
functions with weak derivatives up to order \(k\). Specific spaces such
as \(H(\text{div})\) and \(H(\text{curl})\) are useful for vector fields
in fluid dynamics and electromagnetics.

\part{Part II: 1D Finite Elements}

\chapter{Galerkin's Method and 1D
FEM}\label{galerkins-method-and-1d-fem}

\section{Setup}\label{setup-3}

\section{From Last Week}\label{from-last-week}

\textbf{Key insight}: The weak form from Lecture 1 and Galerkin's method
from Lecture 2 are fundamentally the same.

\textbf{Galerkin's approach} (review):

\begin{enumerate}
\def\labelenumi{\arabic{enumi}.}
\tightlist
\item
  Compute the residual \(Ay_N - f = r_N(x)\)
\item
  Force the residual to be orthogonal to each approximation function:
  \(\int_0^L r_N(x) \phi_i(x)\,dx = 0\), for \(i = 1, 2, \ldots, N\)
\item
  Solve the resulting set of coupled equations
\end{enumerate}

However, the primary challenge remains: \textbf{no systematic way for
choosing the approximation functions (yet)}. The Finite Element Method
(FEM) is designed specifically to address this problem --- it is based
on Galerkin's method, provides a computationally systematic and
efficient approach, removes restrictions on differentiability, and
creates a framework for handling complex geometries.

\section{A 1D Model Problem}\label{a-1d-model-problem}

Today we start solving our first PDE using the finite element method.

As is often the case in real life problems, materials are not
necessarily homogeneous. Moreover, the physical solution we are looking
for is not necessarily smooth.

Consider the 1D bar shown below, composed of two different elastic
solids:

\begin{Shaded}
\begin{Highlighting}[]
\InformationTok{\textasciigrave{}\textasciigrave{}\textasciigrave{}\{python\}}
\InformationTok{\#| code{-}fold: true}
\InformationTok{\#| fig{-}cap: "1D bar composed of two different materials"}
\InformationTok{create\_fem\_diagram(}
\InformationTok{    width=8,}
\InformationTok{    height=0.8,}
\InformationTok{    num\_elements=12,}
\InformationTok{    use\_gradient=False,}
\InformationTok{    element\_color=\textquotesingle{}\#c0c0c0\textquotesingle{},}
\InformationTok{    show\_element\_labels=True,}
\InformationTok{)}
\InformationTok{\textasciigrave{}\textasciigrave{}\textasciigrave{}}
\end{Highlighting}
\end{Shaded}

\begin{figure}[htbp]

{\centering \pandocbounded{\includegraphics[keepaspectratio]{ch04-1d-fem_files/figure-pdf/cell-3-output-1.png}}

}

\caption{1D bar composed of two different materials}

\end{figure}%

The bar is subjected to tractions \(t^*\) on its left side
(\(\Gamma_t\)), while the right side is fixed (\(\Gamma_u\)). The bar is
composed of two different materials with Young's moduli \(E_1\) and
\(E_2\).

\section{Strong Form}\label{strong-form}

\begin{Shaded}
\begin{Highlighting}[]
\InformationTok{\textasciigrave{}\textasciigrave{}\textasciigrave{}\{python\}}
\InformationTok{\#| code{-}fold: true}
\InformationTok{print\_styled("The strong form of the PDE:")}
\InformationTok{Strong\_form = diff(sigma,x)+f}
\InformationTok{Lprint(Eq(Strong\_form,0))}
\InformationTok{print\_styled(\textquotesingle{}This PDE is a balance of forces, with $f$ being the body force and $\textbackslash{}\textbackslash{}sigma$ the stress.\textquotesingle{})}
\InformationTok{print\_styled("The boundary conditions:")}
\InformationTok{print\_styled("$u|\_\{\textbackslash{}\textbackslash{}Gamma\_u\} = u\^{}*$")}
\InformationTok{print\_styled("$t|\_\{\textbackslash{}\textbackslash{}Gamma\_t\} = t\^{}*$")}
\InformationTok{\textasciigrave{}\textasciigrave{}\textasciigrave{}}
\end{Highlighting}
\end{Shaded}

The strong form of the PDE:

$\displaystyle f + \frac{d}{d x} \sigma{\left(x \right)} = 0$

This PDE is a balance of forces, with \(f\) being the body force and
\(\sigma\) the stress.

The boundary conditions:

\(u|_{\Gamma_u} = u^*\)

\(t|_{\Gamma_t} = t^*\)

Since we assume the body to be linearly elastic:

\begin{Shaded}
\begin{Highlighting}[]
\InformationTok{\textasciigrave{}\textasciigrave{}\textasciigrave{}\{python\}}
\InformationTok{\#| code{-}fold: true}
\InformationTok{print\_styled("The stress $\textbackslash{}\textbackslash{}sigma$ is given by:")}
\InformationTok{Elastic\_1d = E*epsilon}
\InformationTok{Lprint(Eq(sigma,Elastic\_1d))}
\InformationTok{print\_styled(\textquotesingle{}where $E$ is the Young\textbackslash{}\textquotesingle{}s modulus and $\textbackslash{}\textbackslash{}epsilon$ is the strain.\textquotesingle{})}
\InformationTok{\textasciigrave{}\textasciigrave{}\textasciigrave{}}
\end{Highlighting}
\end{Shaded}

The stress \(\sigma\) is given by:

$\displaystyle \sigma{\left(x \right)} = E \epsilon{\left(x \right)}$

where \(E\) is the Young's modulus and \(\epsilon\) is the strain.

For 1D small displacements:

\begin{Shaded}
\begin{Highlighting}[]
\InformationTok{\textasciigrave{}\textasciigrave{}\textasciigrave{}\{python\}}
\InformationTok{\#| code{-}fold: true}
\InformationTok{print\_styled("The strain is:")}
\InformationTok{Strain\_1d = diff(u,x)}
\InformationTok{Lprint(Eq(epsilon,Strain\_1d))}
\InformationTok{print\_styled("We can write the stress as:")}
\InformationTok{Lprint(Eq(sigma,E*diff(u,x)))}
\InformationTok{\textasciigrave{}\textasciigrave{}\textasciigrave{}}
\end{Highlighting}
\end{Shaded}

The strain is:

$\displaystyle \epsilon{\left(x \right)} = \frac{d}{d x} u{\left(x \right)}$

We can write the stress as:

$\displaystyle \sigma{\left(x \right)} = E \frac{d}{d x} u{\left(x \right)}$

\section{Getting the Weak Form}\label{getting-the-weak-form}

To get the weak form we multiply by a test function \(v\) and integrate
over the domain \(\Omega\):

\begin{Shaded}
\begin{Highlighting}[]
\InformationTok{\textasciigrave{}\textasciigrave{}\textasciigrave{}\{python\}}
\InformationTok{\#| code{-}fold: true}
\InformationTok{residual = Integral(r*v, (x,0,L))}
\InformationTok{weak\_term = Integral((diff(sigma,x)+f)*v, (x,0,L))}
\InformationTok{Lprint(Eq(weak\_term,0))}
\InformationTok{\textasciigrave{}\textasciigrave{}\textasciigrave{}}
\end{Highlighting}
\end{Shaded}

$\displaystyle \int\limits_{0}^{L} \left(f + \frac{d}{d x} \sigma{\left(x \right)}\right) v{\left(x \right)}\, dx = 0$

Using integration by parts:

\begin{Shaded}
\begin{Highlighting}[]
\InformationTok{\textasciigrave{}\textasciigrave{}\textasciigrave{}\{python\}}
\InformationTok{\#| code{-}fold: true}
\InformationTok{print\_styled("Using integration by parts:")}
\InformationTok{chain\_rule = Derivative(sigma*v,x)}
\InformationTok{Lprint(Eq(chain\_rule,chain\_rule.doit(),evaluate=False))}
\InformationTok{print\_styled("And thus:")}
\InformationTok{Lprint(Eq(v*Derivative(sigma,x), Derivative(sigma*v,x){-}sigma*Derivative(v,x)))}
\InformationTok{\textasciigrave{}\textasciigrave{}\textasciigrave{}}
\end{Highlighting}
\end{Shaded}

Using integration by parts:

$\displaystyle \frac{d}{d x} \sigma{\left(x \right)} v{\left(x \right)} = \sigma{\left(x \right)} \frac{d}{d x} v{\left(x \right)} + v{\left(x \right)} \frac{d}{d x} \sigma{\left(x \right)}$

And thus:

$\displaystyle v{\left(x \right)} \frac{d}{d x} \sigma{\left(x \right)} = - \sigma{\left(x \right)} \frac{d}{d x} v{\left(x \right)} + \frac{d}{d x} \sigma{\left(x \right)} v{\left(x \right)}$

\begin{Shaded}
\begin{Highlighting}[]
\InformationTok{\textasciigrave{}\textasciigrave{}\textasciigrave{}\{python\}}
\InformationTok{\#| code{-}fold: true}
\InformationTok{print\_styled("We can write the weak form as:")}
\InformationTok{weak\_form\_before\_bc = Integral(chain\_rule,(x,0,L)){-}Integral(sigma*Derivative(v,x),(x,0,L))+Integral(f*v,(x,0,L))}
\InformationTok{Lprint(Eq(weak\_form\_before\_bc,0))}
\InformationTok{\textasciigrave{}\textasciigrave{}\textasciigrave{}}
\end{Highlighting}
\end{Shaded}

We can write the weak form as:

$\displaystyle \int\limits_{0}^{L} f v{\left(x \right)}\, dx - \int\limits_{0}^{L} \sigma{\left(x \right)} \frac{d}{d x} v{\left(x \right)}\, dx + \int\limits_{0}^{L} \frac{d}{d x} \sigma{\left(x \right)} v{\left(x \right)}\, dx = 0$

Requiring \(v|_{\Gamma_u}=0\) and \(\sigma|_{\Gamma_t}=t^*\):

\begin{Shaded}
\begin{Highlighting}[]
\InformationTok{\textasciigrave{}\textasciigrave{}\textasciigrave{}\{python\}}
\InformationTok{\#| code{-}fold: true}
\InformationTok{print\_styled("Starting from the boundary term:")}
\InformationTok{weak\_form\_3 = Integral(chain\_rule,(x,0,L))}
\InformationTok{sv = sigma*v}
\InformationTok{Lprint(Eq(weak\_form\_3,sv.subs(\{x:L\}){-}sv.subs(\{x:0\})))}
\InformationTok{print\_styled("Recalling that $v$ is 0 on $\textbackslash{}\textbackslash{}Gamma\_u$:")}
\InformationTok{Lprint(Eq(weak\_form\_3,t\_star*v.subs(\{x:L\})))}
\InformationTok{print\_styled("The weak form becomes:")}
\InformationTok{weak\_LHS = Integral(sigma*Derivative(v,x),(x,0,L))}
\InformationTok{weak\_RHS = Integral(f*v,(x,0,L))+t\_star*v.subs(\{x:L\})}
\InformationTok{Lprint(Eq(weak\_LHS,weak\_RHS))}
\InformationTok{\textasciigrave{}\textasciigrave{}\textasciigrave{}}
\end{Highlighting}
\end{Shaded}

Starting from the boundary term:

$\displaystyle \int\limits_{0}^{L} \frac{d}{d x} \sigma{\left(x \right)} v{\left(x \right)}\, dx = - \sigma{\left(0 \right)} v{\left(0 \right)} + \sigma{\left(L \right)} v{\left(L \right)}$

Recalling that \(v\) is 0 on \(\Gamma_u\):

$\displaystyle \int\limits_{0}^{L} \frac{d}{d x} \sigma{\left(x \right)} v{\left(x \right)}\, dx = t^{*} v{\left(L \right)}$

The weak form becomes:

$\displaystyle \int\limits_{0}^{L} \sigma{\left(x \right)} \frac{d}{d x} v{\left(x \right)}\, dx = t^{*} v{\left(L \right)} + \int\limits_{0}^{L} f v{\left(x \right)}\, dx$

\section{The Weak Form}\label{the-weak-form}

To get the final weak form in terms of the displacement \(u\), we
introduce the constitutive and kinematic relations:

\begin{Shaded}
\begin{Highlighting}[]
\InformationTok{\textasciigrave{}\textasciigrave{}\textasciigrave{}\{python\}}
\InformationTok{\#| code{-}fold: true}
\InformationTok{print\_styled("Substituting the constitutive relation:")}
\InformationTok{weak\_LHS\_sub = weak\_LHS.subs(\{sigma:E*diff(u,x)\})}
\InformationTok{Lprint(Eq(weak\_LHS\_sub,weak\_RHS))}
\InformationTok{\textasciigrave{}\textasciigrave{}\textasciigrave{}}
\end{Highlighting}
\end{Shaded}

Substituting the constitutive relation:

$\displaystyle \int\limits_{0}^{L} E \frac{d}{d x} u{\left(x \right)} \frac{d}{d x} v{\left(x \right)}\, dx = t^{*} v{\left(L \right)} + \int\limits_{0}^{L} f v{\left(x \right)}\, dx$

The final equation is the \textbf{Weak Form} or \textbf{Variational
Formulation}. We can write it as:

\[\mathbb{B}(u,v) = \mathbb{F}(v)\]

where:

\begin{itemize}
\tightlist
\item
  \(\mathbb{B}(u,v) = \int_0^L \frac{dv}{dx} E \frac{du}{dx} \, dx\)
  represents internal virtual work
\item
  \(\mathbb{F}(v) = \int_0^L f v \, dx + t^* v|_{\Gamma_t}\) represents
  external virtual work
\end{itemize}

This form requires less smoothness on the solution \(u\) compared to the
strong form and is the starting point for the Finite Element Method.

\section{Function Spaces for the Weak
Form}\label{function-spaces-for-the-weak-form}

For well-defined integrals, we need:

\[\forall v, \quad \int_0^L \frac{dv}{dx} \sigma \, dx < \infty, \quad \int_0^L f v \, dx < \infty\]

This leads to using Sobolev spaces:

\[u \in H^1(\Omega) \iff \int_\Omega \frac{du}{dx}\frac{du}{dx} \, dx + \int_\Omega u\, u \, dx < \infty\]

Our problem becomes: Find \(u \in H^1(\Omega)\),
\(u|_{\Gamma_u} = u^*\), such that \(\forall v \in H^1(\Omega)\),
\(v|_{\Gamma_u} = 0\):

\[\int_0^L \frac{dv}{dx} E \frac{du}{dx} \, dx = \int_0^L f v \, dx + t^* v|_{\Gamma_t}\]

\section{Approximating the Solution}\label{approximating-the-solution}

We discretize the problem using:

\[u_N \approx u(x) = \sum_{j=1}^{N} a_j \phi_j(x), \quad v = \phi_i(x)\]

This gives us a linear system \([K]\{a\} = \{F\}\) where:

\[K_{ij} = \int_0^L \frac{d\phi_i}{dx} E \frac{d\phi_j}{dx} \, dx, \qquad F_i = \int_0^L f \phi_i \, dx + t^* \phi_i|_{\Gamma_t}\]

\section{Basis Functions in FEM}\label{basis-functions-in-fem}

The key insight of FEM is a systematic way to construct basis functions:

\begin{itemize}
\tightlist
\item
  Functions are built piecewise over subdomains (``elements'')
\item
  Usually low-degree polynomials
\item
  Smooth enough to be in \(H^1(\Omega)\)
\item
  Form a nodal basis where \(\phi_i(x_j) = \delta_{ij}\) (Kronecker
  delta)
\item
  The sum over the basis functions satisfies:
  \(\sum_{j\in \text{element}} \phi_j(x) = 1\) (partition of unity)
\item
  Each basis function is non-zero only on elements sharing node \(i\)
\end{itemize}

\subsection{Linear Elements}\label{linear-elements}

For linear elements, we define:

\[\phi_i(x) =
\begin{cases}
\frac{x-x_{i-1}}{h_i}, & x_{i-1} \leq x \leq x_i \\
1 - \frac{x-x_i}{h_{i+1}}, & x_i \leq x \leq x_{i+1} \\
0, & \text{otherwise}
\end{cases}\]

where \(h_i = x_i - x_{i-1}\). The derivatives are:

\[\frac{d\phi_i}{dx} =
\begin{cases}
\frac{1}{h_i}, & x_{i-1} \leq x \leq x_i \\
-\frac{1}{h_{i+1}}, & x_i \leq x \leq x_{i+1} \\
0, & \text{otherwise}
\end{cases}\]

\subsection{Visualization of Basis
Functions}\label{visualization-of-basis-functions}

\begin{Shaded}
\begin{Highlighting}[]
\InformationTok{\textasciigrave{}\textasciigrave{}\textasciigrave{}\{python\}}
\InformationTok{\#| code{-}fold: true}
\InformationTok{\#| fig{-}cap: "Lagrange shape functions of degree 1, 2, and 3"}
\InformationTok{fem\_viz = FEMShapeFunctions()}
\InformationTok{fig1 = fem\_viz.plot\_multielement\_shape\_functions(}
\InformationTok{    num\_elements=3,}
\InformationTok{    element\_type=\textquotesingle{}lagrange\textquotesingle{},}
\InformationTok{    degrees=[1, 2, 3])}
\InformationTok{\textasciigrave{}\textasciigrave{}\textasciigrave{}}
\end{Highlighting}
\end{Shaded}

\begin{figure}[htbp]

{\centering \pandocbounded{\includegraphics[keepaspectratio]{ch04-1d-fem_files/figure-pdf/cell-12-output-1.png}}

}

\caption{Lagrange shape functions of degree 1, 2, and 3}

\end{figure}%

\begin{Shaded}
\begin{Highlighting}[]
\InformationTok{\textasciigrave{}\textasciigrave{}\textasciigrave{}\{python\}}
\InformationTok{\#| code{-}fold: true}
\InformationTok{\#| fig{-}cap: "Derivatives of Lagrange shape functions"}
\InformationTok{fem\_viz = FEMShapeFunctions()}
\InformationTok{fig2 = fem\_viz.plot\_multielement\_shape\_function\_derivatives(}
\InformationTok{    num\_elements=3,}
\InformationTok{    element\_type=\textquotesingle{}lagrange\textquotesingle{},}
\InformationTok{    degrees=[1, 2]}
\InformationTok{)}
\InformationTok{\textasciigrave{}\textasciigrave{}\textasciigrave{}}
\end{Highlighting}
\end{Shaded}

\begin{figure}[htbp]

{\centering \pandocbounded{\includegraphics[keepaspectratio]{ch04-1d-fem_files/figure-pdf/cell-13-output-1.png}}

}

\caption{Derivatives of Lagrange shape functions}

\end{figure}%

\begin{Shaded}
\begin{Highlighting}[]
\InformationTok{\textasciigrave{}\textasciigrave{}\textasciigrave{}\{python\}}
\InformationTok{\#| code{-}fold: true}
\InformationTok{\#| fig{-}cap: "Partition of unity for linear elements"}
\InformationTok{fem\_viz = FEMShapeFunctions()}
\InformationTok{fig3 = fem\_viz.plot\_multielement\_partition\_of\_unity\_per\_element(}
\InformationTok{    num\_elements=3,}
\InformationTok{    degree=1}
\InformationTok{)}
\InformationTok{\textasciigrave{}\textasciigrave{}\textasciigrave{}}
\end{Highlighting}
\end{Shaded}

\begin{figure}[htbp]

{\centering \pandocbounded{\includegraphics[keepaspectratio]{ch04-1d-fem_files/figure-pdf/cell-14-output-1.png}}

}

\caption{Partition of unity for linear elements}

\end{figure}%

\section{Building the Stiffness
Matrix}\label{building-the-stiffness-matrix}

For example, to compute \(K_{11}\):

\[K_{11} = \int_0^L \frac{d\phi_1}{dx} E \frac{d\phi_1}{dx} \, dx = \int_0^{x_1} E \frac{1}{h_1^2} \, dx = \frac{1}{h_1^2} \int_0^{x_1} E(x) \, dx\]

Similarly for \(K_{12}\):

\[K_{12} = \int_0^L \frac{d\phi_1}{dx} E \frac{d\phi_2}{dx} \, dx = -\frac{1}{h_1^2} \int_0^{x_1} E(x) \, dx\]

And computing \(K_{13}\):

\[K_{13} = \int_0^L \frac{d\phi_1}{dx} E \frac{d\phi_3}{dx} \, dx = 0\]

since \(\phi_1\) and \(\phi_3\) have no overlapping support. This
pattern shows why FEM matrices have a \textbf{banded structure} ---
there are only local interactions between adjacent elements.

\section{Element-wise Assembly}\label{element-wise-assembly}

We recognize that:

\[K_{ij}^g = \sum_{\text{elem}} K_{ij}^e, \quad \text{where } K_{ij}^e = \int_{\Omega_e} \frac{d\phi_i}{dx} E \frac{d\phi_j}{dx} \, dx\]

\[F_i^g = \sum_{\text{elem}} F_i^e, \quad F_i^e = \int_{\Omega_e} \phi_i f \, dx + t^* \phi_i|_{\Gamma_t \cap \Omega_e}\]

This element-wise assembly is key to FEM efficiency.

\section{Reference Element
Transformation}\label{reference-element-transformation}

To make calculations systematic, we define a master (reference) element
in a local coordinate system \(\zeta \in [-1, 1]\).

We need a mapping from reference to physical coordinates:

\[x = \sum_{i=1}^2 \chi_i \hat{\phi}_i(\zeta) = M_x(\zeta)\]

where \(\chi_i\) are the physical coordinates of the nodes and
\(\hat{\phi}_i(\zeta)\) are the shape functions in reference
coordinates.

\textbf{Shape Functions in Reference Space:} For linear elements in 1D:

\[\hat{\phi}_1 = \frac{1}{2}(1-\zeta) \quad \text{and} \quad \hat{\phi}_2 = \frac{1}{2}(1+\zeta)\]

\section{Differential Relations}\label{differential-relations}

We need the mapping between derivatives:

\[\frac{d}{d\zeta} = \frac{dx}{d\zeta}\frac{d}{dx} = J \frac{d}{dx}\]

And the inverse relation:

\[\frac{d}{dx} = \frac{d\zeta}{dx}\frac{d}{d\zeta} = \frac{1}{J}\frac{d}{d\zeta}\]

where \(J = \frac{dx}{d\zeta}\) is the Jacobian of the transformation.

\section{Numerical Integration: Gaussian
Quadrature}\label{numerical-integration-gaussian-quadrature}

For numerical integration, we use Gaussian quadrature:

\[\int_{x_i}^{x_{i+1}} F(x) \, dx = \int_{-1}^1 F(x(\zeta)) J(\zeta) \, d\zeta \approx \sum_{q=1}^G w_q F(x(\zeta_q)) J(\zeta_q)\]

Gauss quadrature can integrate a polynomial of order \(2G-1\) exactly
with \(G\) points.

\begin{longtable}[]{@{}ccc@{}}
\toprule\noalign{}
Gauss rule & \(\zeta_i\) & \(w_i\) \\
\midrule\noalign{}
\endhead
\bottomrule\noalign{}
\endlastfoot
2 & \(\pm 0.5773502692\) & 1.0 \\
3 & 0.0 & 0.8888888889 \\
& \(\pm 0.7745966692\) & 0.5555555556 \\
\end{longtable}

Two points are sufficient for linear elements since the integrands in
\(K_{ij}\) have at most cubic terms.

\section{Computing Element Matrices with
Quadrature}\label{computing-element-matrices-with-quadrature}

The element stiffness matrix is computed as:

\[K_{ij}^e = \sum_{q=1}^G w_q \left[\frac{d\hat{\phi}_i}{d\zeta}\frac{d\zeta}{dx}\right] E \left[\frac{d\hat{\phi}_j}{d\zeta}\frac{d\zeta}{dx}\right] J \bigg|_{\zeta=\zeta_q}\]

Similarly for the force vector:

\[F_i^e = \sum_{q=1}^G w_q \hat{\phi}_i(\zeta_q) f(x(\zeta_q)) J(\zeta_q)\]

plus any boundary terms.

\section{Post-Processing}\label{post-processing}

After obtaining the nodal values \(a_i\), we can compute:

\begin{itemize}
\tightlist
\item
  \textbf{Displacements} at any point:
  \(u(x) = \sum_{i=1}^N a_i \phi_i(x)\)
\item
  \textbf{Strains} at any point:
  \(\varepsilon(x) = \frac{du}{dx} = \sum_{i=1}^N a_i \frac{d\phi_i}{dx}\)
\item
  \textbf{Stresses} at any point:
  \(\sigma(x) = E(x) \sum_{i=1}^N a_i \frac{d\phi_i}{dx}\)
\end{itemize}

\section{Implementing Boundary
Conditions}\label{implementing-boundary-conditions}

Essential (Dirichlet) boundary conditions (e.g., \(u(0) = u^*\)):

\[\begin{bmatrix}
1 & 0 & \ldots & \ldots & 0 \\
K_{21} & K_{22} & \ldots & \ldots & K_{2N} \\
\vdots & \vdots & \ddots & \vdots & \vdots \\
K_{N1} & K_{N2} & \ldots & \ldots & K_{NN}
\end{bmatrix}
\begin{bmatrix}
a_1 \\ a_2 \\ \vdots \\ a_N
\end{bmatrix}
=
\begin{bmatrix}
u^* \\ F_2 \\ \vdots \\ F_N + \ldots
\end{bmatrix}\]

\chapter{Time-Dependent 1D Problems}\label{time-dependent-1d-problems}

\section{Problem Introduction}\label{problem-introduction}

In the following, we use the FEM methods for spatial discretization and
the finite difference method for time discretization.

Consider the Balance of Linear momentum equation in 1D:

\[
\begin{aligned}
\text{grad} \cdot \sigma+f(x,t)&=\rho_0\frac{\partial^2u}{\partial t^2} \\
&=\rho_0\frac{\partial v}{\partial t} \\
&=\rho_0 a(x,t)
\end{aligned}
\]

where:

\begin{itemize}
\tightlist
\item
  \(\sigma(x,t)\) is the stress
\item
  \(f(x,t)\) is the body force
\item
  \(\rho_0\) is the density
\item
  \(u\) is the displacement
\item
  \(v\) is the velocity
\item
  \(a\) is the acceleration
\end{itemize}

\section{Single Point in Space}\label{single-point-in-space}

For a specific point in space, we can write the equation as:

\[
m\dot{v} = F = ma \quad (\text{I})
\]

where all the forces acting on the point are lumped into a single force
\(F\) and the mass is \(m=\rho_0\Delta x\).

In the following, we will use the Taylor expansion around the time
\(t+\theta \Delta t\) to expand the quantities of interest at times
\(t\) and \(t+\Delta t\) where \(\theta\) is a parameter that can be
chosen to optimize the time integration scheme.

\subsection{Velocity Expansions}\label{velocity-expansions}

For the velocities:

\[
\begin{aligned}
v(t + \Delta t) &= v(t + \theta \Delta t) + \frac{dv}{dt}\bigg|_{t+\theta \Delta t} (1 - \theta) \Delta t + \frac{1}{2} \frac{d^2 v}{dt^2}\bigg|_{t+\theta \Delta t} (1 - \theta)^2 (\Delta t)^2 + O(\Delta t)^3 \quad (2.1)\\
v(t) &= v(t + \theta \Delta t) - \frac{dv}{dt}\bigg|_{t+\theta \Delta t} \theta \Delta t + \frac{1}{2} \frac{d^2 v}{dt^2}\bigg|_{t+\theta \Delta t} \theta^2 (\Delta t)^2 + O(\Delta t)^3 \quad (2.2)
\end{aligned}
\]

\subsection{Position Expansions}\label{position-expansions}

For the position:

\[
\begin{aligned}
u(t + \Delta t) &= u(t + \theta \Delta t) + \frac{du}{dt}\bigg|_{t+\theta \Delta t} (1 - \theta) \Delta t + \frac{1}{2} \frac{d^2 u}{dt^2}\bigg|_{t+\theta \Delta t} (1 - \theta)^2 (\Delta t)^2 + O(\Delta t)^3 \quad (2.3) \\
u(t) &= u(t + \theta \Delta t) - \frac{du}{dt}\bigg|_{t+\theta \Delta t} \theta \Delta t + \frac{1}{2} \frac{d^2 u}{dt^2}\bigg|_{t+\theta \Delta t} \theta^2 (\Delta t)^2 + O(\Delta t)^3 \quad (2.4)
\end{aligned}
\]

\subsection{Force Averaging}\label{force-averaging}

For the Force we obtain a weighted average of the forces at the two time
steps:

\[
F(t+\theta \Delta t) = \theta F(t+\Delta t) + (1-\theta) F(t) \quad (2.5)
\]

\subsection{Velocity Time Derivative}\label{velocity-time-derivative}

Subtracting Equations (2.1) and (2.2) we obtain:

\[
\frac{dv}{dt}\bigg|_{t+\theta \Delta t} = \frac{v(t + \Delta t) - v(t)}{\Delta t} + \hat{O}(\Delta t) \quad (2.6)
\]

where \(\hat{O}(\Delta t)\) is the error term that depends on the choice
of \(\theta\). For \(\theta=1/2\) we obtain \(\hat{O}=O(\Delta t)^2\)
and otherwise we get \(\hat{O}=O(\Delta t)\).

By using equation (2.6) and inserting it into equation (I) we obtain:

\[
v(t + \Delta t) = v(t) + \frac{\Delta t}{m} F(t + \theta \Delta t) + \hat{O}(\Delta t)^2 \quad (2.7)
\]

\subsection{Velocity at Intermediate
Time}\label{velocity-at-intermediate-time}

If we take the weighted sum of (2.1) and (2.2) we obtain:

\[
v(t+\theta \Delta t) = \theta v(t+\Delta t) + (1-\theta) v(t) + O(\Delta t)^2 \quad (2.8)
\]

\subsection{Position Time Derivative}\label{position-time-derivative}

We can do the same for the positions and obtain:

\[
\frac{u(t+\Delta t)-u(t)}{\Delta t} = v(t+\theta \Delta t) + \hat{O}(\Delta t) \quad (2.9)
\]

Which combined with equation (2.8) gives us:

\[
u(t + \Delta t) = u(t) + (\theta v(t + \Delta t) + (1 - \theta)v(t)) \Delta t + \hat{O}(\Delta t)^2 \quad (2.10)
\]

\subsection{Final Position Update}\label{final-position-update}

Now, we can use equations (2.10) and (2.7) to obtain:

\[
u(t + \Delta t) = u(t) + v(t) \Delta t + \frac{\theta (\Delta t)^2}{m} F(t + \theta \Delta t) + \hat{O}(\Delta t)^2 \quad(2.11)
\]

\subsection{Time Integration Schemes}\label{time-integration-schemes}

When \(\theta = 1\) we obtain an implicit integration scheme (backward
Euler) and when \(\theta = 0\) we obtain an explicit integration scheme
(forward Euler).

The backward Euler scheme is unconditionally stable, while the forward
Euler scheme is conditionally stable and requires a small time step to
ensure stability.

For both we obtain \(\hat{O}(\Delta t)^2 = O(\Delta t)^2\).

For \(\theta = 1/2\) we obtain the ``midpoint rule'' which is
unconditionally stable and \(\hat{O}(\Delta t)^2 = O(\Delta t)^3\).

\section{FEM Formulation}\label{fem-formulation}

The continuum Equation we try to solve is:

\[
\rho_0\frac{\partial^2 u}{\partial t^2} = \rho_0\frac{\partial v}{\partial t} = \nabla \cdot \sigma + f(x,t) \quad (3.0)
\]

We will make our life easier by defining
\(F=\nabla \cdot \sigma + f(x,t)\).

Using the equations we derived above we can reformulate the problem as:

\[
\rho_0v(t+\Delta t) = \rho_0v(t) + \Delta t(\theta F(t+\Delta t) + (1-\theta) F(t)) \quad (3.1)
\]

\subsection{Weak Form Development}\label{weak-form-development}

Multiplying by the test function and integrating by parts we obtain:

\[
\int_\Omega v \cdot \rho_0 v(t + \Delta t) \, d\Omega = \int_\Omega v \cdot \rho_0 v(t) \, d\Omega + \Delta t \int_\Omega v \cdot (\theta \Psi(t + \Delta t) + (1 - \theta) \Psi(t)) \, d\Omega \quad (3.2)
\]

And after applying the divergence theorem and enforcing \(v=0\) on the
boundary we obtain:

\[
\begin{aligned}
\int_\Omega v \cdot \rho_0 v^{t+\Delta t} \, d\Omega &= \int_\Omega v \cdot \rho_0 v^{t} \, d\Omega \\
&+ \Delta t \theta \left( -\int_\Omega \nabla v : \sigma \, d\Omega + \int_{\Gamma_t} v \cdot (\sigma \cdot n) \, dA + \int_\Omega v \cdot f \, d\Omega \right)^{t+\Delta t} \\
&+ \Delta t (1 - \theta) \left( -\int_\Omega \nabla v : \sigma \, d\Omega + \int_{\Gamma_t} v \cdot t^* \, dA + \int_\Omega v \cdot f \, d\Omega \right)^{t} \qquad (3.3)
\end{aligned}
\]

\subsection{Discretized System}\label{discretized-system}

Plugging in our shape function, and defining the standard matrices with
\(\{a\}\) being the coefficients vector, we obtain:

\[
[M]\{a\}^{t+\Delta t} = [M]\{a\}^t + (\Delta t \theta) \left( -[K]\{a\}^{t+\Delta t} + \{R_f\}^{t+\Delta t} + \{R_t\}^{t+\Delta t} \right) \\
+ \Delta t (1 - \theta) \left( -[K]\{a\}^t + \{R_f\}^t + \{R_t\}^t \right) \qquad (3.4)
\]

where we use:

\[
\{a\}^{t+\Delta t} = \{a\}^{t}+\Delta t (\theta \{\dot{a}\})^{t+\Delta t} + (1 - \theta) \{\dot{a}\}^{t} \qquad (3.5)
\]

\section{Implicit Backward Euler
Scheme}\label{implicit-backward-euler-scheme}

For \(\theta=1\) we get the implicit backward Euler scheme:

\[
\left( [M]\{\ddot{a}\}^{t+\Delta t} + \Delta t [K]\{a\}^{t+\Delta t} \right) = [M]\{\ddot{a}\}^t + \Delta t \left( \{R_t\}^{t+\Delta t} + \{R_f\}^{t+\Delta t} \right) \qquad (3.6)
\]

where we use:

\[
\{a\}^{t+\Delta t} = \{a\}^{t}+\Delta t (\{\dot{a}\})^{t+\Delta t}  \qquad (3.7)
\]

\section{Explicit Forward Euler
Scheme}\label{explicit-forward-euler-scheme}

For \(\theta=0\) we get the explicit forward Euler scheme:

\[
\{\dot{a}\}^{t+\Delta t} = \{\dot{a}\}^t + \Delta t [M]^{-1} \left( -[K]\{a\}^t + \{R_f\}^t + \{R_t\}^t \right) \qquad (3.8)
\]

where we use:

\[
\{a\}^{t+\Delta t} = \{a\}^t + \Delta t \{\dot{a}\}^t \qquad (3.9)
\]

\part{Part III: Multi-Dimensional FEM}

\chapter{From 1D to Multi-Dimensional
FEM}\label{from-1d-to-multi-dimensional-fem}

\section{Mathematical Preliminaries}\label{mathematical-preliminaries}

Looking at the following equation:

\[
\text{div}(\mathbf{K}\cdot \text{grad}(T))
\]

if \(\mathbf{K}\) is a constant multiplied by the identity matrix, we
can write it as:

\[
\text{div}(\mathbf{K}\cdot \text{grad}(T))  = \text{div}(\mathbf{K} \left [\frac{\partial T}{\partial x} \mathbf{e}_1 +\frac{\partial T}{\partial y} \mathbf{e}_2+\frac{\partial T}{\partial z} \mathbf{e}_3 \right ]) = \mathbf{K} \left ( \frac{\partial^2 T}{\partial x^2}+\frac{\partial^2 T}{\partial y^2}+\frac{\partial^2 T}{\partial z^2} \right) = \mathbf{K}\nabla^2T
\]

\subsection{Temperature as a Scalar
Field}\label{temperature-as-a-scalar-field}

\[
T \text{ - Scalar} \Rightarrow \text{grad}(T) = \left(\frac{\partial T}{\partial x}\mathbf{e_1} + \frac{\partial T}{\partial y}\mathbf{e_2} + \frac{\partial T}{\partial z}\mathbf{e_3}\right) \text{ - vector}
\]

\subsection{Conductivity Matrix
Operation}\label{conductivity-matrix-operation}

\[
\mathbf{K} \cdot \text{grad}(T) = \begin{bmatrix} K_{11} & K_{12} & K_{13} \\ K_{21} & K_{22} & K_{23} \\ K_{31} & K_{32} & K_{33} \end{bmatrix} \text{transpose}\left(\left(\frac{\partial T}{\partial x}\mathbf{e_1} + \frac{\partial T}{\partial y}\mathbf{e_2} + \frac{\partial T}{\partial z}\mathbf{e_3}\right)\right)
\]

\subsection{Heat Flux Vector}\label{heat-flux-vector}

Let's define \(\mathbf{q} = -\mathbf{K} \cdot \text{grad}(T)\) as the
heat flux vector.

Divergence of heat flux:

\[
\text{div}(\mathbf{q}) = \frac{\partial q_x}{\partial x} + \frac{\partial q_y}{\partial y} + \frac{\partial q_z}{\partial z}
\]

\section{The Product Rule for
Divergence}\label{the-product-rule-for-divergence}

To answer the question ``What happens when we multiply the heat
conduction equation by a test function V?'', we need to use the
\textbf{product rule for divergence}:

\[
\text{div}(V \mathbf{K} \text{ grad}(T)) = \text{grad}(V) \cdot \mathbf{K} \cdot \text{grad}(T) + V \text{ div}(\mathbf{K} \cdot \text{grad}(T))
\]

However, if \(\mathbf{K}=k\) is constant, we can simplify this to:

\[
\text{div}(V \mathbf{K} \text{ grad}(T))=k\text{div}(V \text{ grad}(T))=k(\text{grad}(V) \cdot \text{grad}(T) + V \nabla^2(T))
\]

\subsection{Detailed Expansion}\label{detailed-expansion}

Let's expand \(\text{div}(V \mathbf{K} \text{ grad}(T))\) step by step:

\[\text{div}(V \mathbf{K} \text{ grad}(T)) = \frac{\partial}{\partial x}\left(V K \frac{\partial T}{\partial x}\right) + \frac{\partial}{\partial y}\left(V K \frac{\partial T}{\partial y}\right) + \frac{\partial}{\partial z}\left(V K \frac{\partial T}{\partial z}\right)\]

Applying the product rule to each term:

\[
= \frac{\partial V}{\partial x} K \frac{\partial T}{\partial x} + \frac{\partial V}{\partial y} K \frac{\partial T}{\partial y} + \frac{\partial V}{\partial z} K \frac{\partial T}{\partial z} + V K \frac{\partial^2 T}{\partial x^2} + V K \frac{\partial^2 T}{\partial y^2} + V K \frac{\partial^2 T}{\partial z^2}
\]

We can group these terms into two parts:

\textbf{Gradient Terms:}

\[
\text{grad}(V) \cdot \mathbf{K} \cdot \text{grad}(T) = \left(\frac{\partial V}{\partial x}, \frac{\partial V}{\partial y}, \frac{\partial V}{\partial z}\right) \cdot \mathbf{K} \cdot \left(\frac{\partial T}{\partial x}, \frac{\partial T}{\partial y}, \frac{\partial T}{\partial z}\right)
\]

For constant K, this becomes:

\[
= K\left(\frac{\partial V}{\partial x} \frac{\partial T}{\partial x} + \frac{\partial V}{\partial y} \frac{\partial T}{\partial y} + \frac{\partial V}{\partial z} \frac{\partial T}{\partial z}\right) = A
\]

\textbf{Second Derivative Terms:}

\[
V \text{ div}(\mathbf{K} \cdot \text{grad}(T)) = V K \left(\frac{\partial^2 T}{\partial x^2} + \frac{\partial^2 T}{\partial y^2} + \frac{\partial^2 T}{\partial z^2}\right) = B
\]

\subsection{The Key Relationship}\label{the-key-relationship}

From the product rule, we have:

\[
\text{div}(V \mathbf{K} \text{ grad}(T)) = A + B
\]

Therefore:

\[
B = \text{div}(V \mathbf{K} \text{ grad}(T)) - A
\]

Which gives us:

\[
\text{div}(\mathbf{K} \text{ grad}(T)) V = \text{div}(V \mathbf{K} \text{ grad}(T)) - \text{grad}(V) \cdot \mathbf{K} \cdot \text{grad}(T)
\]

\subsection{Physical Interpretation}\label{physical-interpretation}

This mathematical identity is the foundation for:

\begin{enumerate}
\def\labelenumi{\arabic{enumi}.}
\tightlist
\item
  \textbf{Integration by parts} in the weak formulation
\item
  \textbf{Converting strong form to weak form} in finite element methods
\item
  \textbf{Moving derivatives from the solution to the test function}
\end{enumerate}

The term \(\text{div}(V \mathbf{K} \text{ grad}(T))\) will become a
boundary integral when integrated over the domain, while
\(\text{grad}(V) \cdot \mathbf{K} \cdot \text{grad}(T)\) becomes the
bilinear form in the finite element formulation.

\subsection{Weak Form Foundation}\label{weak-form-foundation}

When integrated over a domain \(\Omega\), this relationship becomes:

\[\int_\Omega V \text{ div}(\mathbf{K} \text{ grad}(T)) \, d\Omega = \int_\Omega \text{div}(V \mathbf{K} \text{ grad}(T)) \, d\Omega - \int_\Omega \text{grad}(V) \cdot \mathbf{K} \cdot \text{grad}(T) \, d\Omega\]

Using the divergence theorem, the first integral on the right becomes a
boundary integral, leading to the standard weak form used in finite
element analysis.

\section{FEM in 2D/3D: Heat Equation}\label{fem-in-2d3d-heat-equation}

\subsection{Problem Setup}\label{problem-setup}

\textbf{Governing equation:} \(\text{div}(\mathbf{q}) = 0\) where we
assumed that there is no source term for simplicity.

\textbf{Constitutive equation (Fourier's law):}
\(\mathbf{q} = -\mathbf{K} \cdot \text{grad}(T)\)

\textbf{Boundary conditions:}

\begin{itemize}
\tightlist
\item
  \(\Gamma_{T_0}: T = T_0\) (Essential BC)
\item
  \(\Gamma_{T_1}: T = T_1\) (Essential BC)
\item
  \(\Gamma_q: \mathbf{q} = \mathbf{q}^*_n\) (Natural BC)
\end{itemize}

where \(\mathbf{q}^*_n\) is the prescribed heat flux on the boundary
\(\Gamma_q\) in the direction of the (outwards) normal to the surface.

\textbf{Combined form:}
\(\text{div}(\mathbf{q}) = -\text{div}(\mathbf{K} \cdot \text{grad}(T)) = 0\)

\section{Weak Form Derivation}\label{weak-form-derivation}

\textbf{Step 1:} Multiply by test function V and integrate over domain

\[\int_\Omega \text{div}(\mathbf{K} \cdot \text{grad}(T)) V \, d\Omega = 0\]

\textbf{Step 2:} Apply the product rule result

\[\int_\Omega \text{div}(V \mathbf{K} \cdot \text{grad}(T)) \, d\Omega - \int_\Omega \text{grad}(V) \cdot \mathbf{K} \cdot \text{grad}(T) \, d\Omega = 0\]

\textbf{Step 3:} Apply divergence theorem to first integral

\[\int_\Omega \text{div}(V \mathbf{K} \cdot \text{grad}(T)) \, d\Omega = \int_\Omega \text{div}(V \mathbf{q}) \, d\Omega = \int_{\partial\Omega} V \mathbf{q} \cdot \mathbf{n} \, d\Gamma\]

\textbf{Final weak form:}

\[\int_{\partial\Omega} V \mathbf{q} \cdot \mathbf{n} \, d\Gamma - \int_\Omega \text{grad}(V) \cdot \mathbf{K} \cdot \text{grad}(T) \, d\Omega = 0\]

\textbf{Finite element approximation:}

\[V_h = \sum_{i=1}^N b_i \phi_i(x,y), \quad T_h = \sum_{j=1}^N a_j \phi_j(x,y) \Rightarrow \{a_j\}\]

\section{2D/3D Finite Elements: Geometry and Shape
Functions}\label{d3d-finite-elements-geometry-and-shape-functions}

\subsection{Transition from 1D to
2D/3D}\label{transition-from-1d-to-2d3d}

\textbf{Key differences:}

\begin{itemize}
\tightlist
\item
  \textbf{1D:} Geometric description is simpler - element domain is
  simply a line
\item
  \textbf{2D/3D:} Geometry plays a crucial role and complexity increases
  significantly
\end{itemize}

\textbf{Most common finite element geometries:}

\begin{itemize}
\tightlist
\item
  \textbf{2D:} Triangular elements and Rectangular elements
\item
  \textbf{3D:} Tetrahedron elements and Hexahedron elements
\end{itemize}

\subsection{2D Rectangular Elements: Linear Shape
Functions}\label{d-rectangular-elements-linear-shape-functions}

\textbf{Design requirements for shape functions:}

\begin{itemize}
\tightlist
\item
  \textbf{Simplicity:} Describe change in the element's domain with
  minimal complexity
\item
  \textbf{Local support:} Each shape function corresponds to one node,
  value of ``1'' at one node, ``0'' at all other nodes
\end{itemize}

\textbf{1D analogy:} \(u^e = a_i \phi_i^g + a_{i+1} \phi_{i+1}^g\)
(where \(u^e \in [x_i, x_{i+1}]\))

\textbf{2D extension:} For a linear rectangular element, we need
\textbf{four shape functions} to describe \(u^e\)

\subsection{Rectangular Element with Linear Shape
Functions}\label{rectangular-element-with-linear-shape-functions}

Consider a rectangular element with nodes numbered as follows:

\begin{verbatim}
4 ────── 3
│        │  b
│   ue   │
│        │
1 ────── 2
    a
\end{verbatim}

\textbf{Shape functions for specific rectangular element:}

\[\phi_1^g = \frac{1}{ab}(a-x)(b-y)\]

\[\phi_2^g = \frac{1}{ab}x(b-y)\]

\[\phi_3^g = \frac{1}{ab}xy\]

\[\phi_4^g = \frac{1}{ab}(a-x)y\]

\textbf{Verification:}

\begin{itemize}
\tightlist
\item
  \(\phi_1^g(0,0) = 1\), \(\phi_1^g(a,0) = 0\), \(\phi_1^g(0,b) = 0\),
  \(\phi_1^g(a,b) = 0\) ✓
\item
  \(\phi_2^g(0,0) = 0\), \(\phi_2^g(a,0) = 1\), \(\phi_2^g(0,b) = 0\),
  \(\phi_2^g(a,b) = 0\) ✓
\end{itemize}

\textbf{Element solution:}
\(u^e = a_1 \phi_1^g + a_2 \phi_2^g + a_3 \phi_3^g + a_4 \phi_4^g = u^e(x,y)\)

\section{Master Element Concept in
2D}\label{master-element-concept-in-2d}

\subsection{Mapping to Master Element}\label{mapping-to-master-element}

\textbf{Purpose:} Map general elements to a standardized \textbf{master
element} for easier computation

\textbf{Coordinate transformation:}
\([x,y] \rightarrow [\zeta_1, \zeta_2] \in [-1,1] \times [-1,1]\)

\textbf{Master element node numbering convention:}

\begin{verbatim}
4 ────── 3   ζ₂
│        │    ↑
│        │    │
│        │    │
1 ────── 2    └─→ ζ₁
\end{verbatim}

\textbf{Master element shape functions:}

\[\hat{\phi}_1 = \phi_1^e(\zeta_1, \zeta_2) = \frac{1}{4}(1-\zeta_1)(1-\zeta_2)\]

\[\hat{\phi}_2 = \phi_2^e(\zeta_1, \zeta_2) = \frac{1}{4}(1+\zeta_1)(1-\zeta_2)\]

\[\hat{\phi}_3 = \phi_3^e(\zeta_1, \zeta_2) = \frac{1}{4}(1+\zeta_1)(1+\zeta_2)\]

\[\hat{\phi}_4 = \phi_4^e(\zeta_1, \zeta_2) = \frac{1}{4}(1-\zeta_1)(1+\zeta_2)\]

\textbf{Bi-linear nature:} Each \(\hat{\phi}_i\) equals ``1'' at the
\(i\)-th node and ``0'' at other nodes, creating \textbf{bi-linear}
shape functions.

\textbf{Example expansion:}
\(\hat{\phi}_1 = \frac{1}{4}(1-\zeta_1)(1-\zeta_2) = \frac{1}{4}(1-\zeta_1-\zeta_2+\zeta_1\zeta_2)\)

The term \(\zeta_1\zeta_2\) makes it bi-linear (linear in both
coordinates).

\pandocbounded{\includegraphics[keepaspectratio]{images/GLmap.png}}

\section{Jacobian and Coordinate
Transformation}\label{jacobian-and-coordinate-transformation}

\subsection{Jacobian Matrix}\label{jacobian-matrix}

\textbf{Definition:}

\[\mathbf{F} = \begin{bmatrix} \frac{\partial x}{\partial \zeta_1} & \frac{\partial x}{\partial \zeta_2} \\ \frac{\partial y}{\partial \zeta_1} & \frac{\partial y}{\partial \zeta_2} \end{bmatrix}, \quad J = \text{Det}(\mathbf{F}) = \frac{\partial x}{\partial \zeta_1}\frac{\partial y}{\partial \zeta_2} - \frac{\partial x}{\partial \zeta_2}\frac{\partial y}{\partial \zeta_1}\]

\textbf{Critical requirement:} \(J > 0\) throughout the element for a
``good'' mapping

\subsection{Isoparametric Mapping}\label{isoparametric-mapping}

We will use \textbf{isoparametric} mapping, where the same shape
functions are used for both geometry and field variables.

\textbf{Coordinate transformation:}

\[x (\chi_1,\chi_2)= \sum_{i=1}^4 \chi_{1i} \hat{\phi}_i = \chi_{11}\hat{\phi}_1 + \chi_{12}\hat{\phi}_2 + \chi_{13}\hat{\phi}_3 + \chi_{14}\hat{\phi}_4\]

\[y(\chi_1, \chi_2) = \sum_{i=1}^4 \chi_{2i} \hat{\phi}_i\]

Where \(\chi_{1i}\) are global x-coordinates and \(\chi_{2i}\) are
global y-coordinates of element nodes.

\section{Element Quality: Good vs Bad
Elements}\label{element-quality-good-vs-bad-elements}

\subsection{Acceptable Elements}\label{acceptable-elements}

\textbf{Case 1 - Rectangular element:}
\(0 < J(\zeta_1, \zeta_2) < \infty\) throughout element

\begin{itemize}
\tightlist
\item
  Jacobian is constant
\item
  Acceptable
\end{itemize}

\textbf{Case 3 - Distorted but convex:}
\(0 < J(\zeta_1, \zeta_2) < \infty\) throughout element

\begin{itemize}
\tightlist
\item
  Jacobian varies but remains positive and bounded
\item
  Acceptable
\end{itemize}

\subsection{Unacceptable Elements}\label{unacceptable-elements}

\textbf{Case 2 - Incorrect node numbering:} \(J(\zeta_1, \zeta_2) < 0\)
throughout element

\begin{itemize}
\tightlist
\item
  Nodes numbered incorrectly, turning element ``inside out''
\item
  Unacceptable
\end{itemize}

\textbf{Case 4 - Partially negative Jacobian:}
\(J(\zeta_1, \zeta_2) < 0\) in some regions

\begin{itemize}
\tightlist
\item
  Can cause singularities in stiffness matrix
\item
  Unacceptable
\end{itemize}

\begin{figure}[htbp]

{\centering \pandocbounded{\includegraphics[keepaspectratio]{images/elementscases.png}}

}

\caption{Element quality cases}

\end{figure}%

\textbf{Key insight for linear elements:} The primary indicator of
problematic elements is \textbf{non-convexity}, even with correct
numbering.

\section{Bookkeeping in 2D Finite
Elements}\label{bookkeeping-in-2d-finite-elements}

\subsection{Difference from 1D}\label{difference-from-1d}

\textbf{1D:} Node numbering is intuitive and element connectivity is
straightforward

\begin{verbatim}
1 ── 2 ── 3 ── 4 ── 5 ── 6 ── 7 ── 8 ── 9
\end{verbatim}

\textbf{2D:} More complex connectivity patterns require systematic
bookkeeping

\subsection{Element Connectivity
Tables}\label{element-connectivity-tables}

\textbf{Element-to-node connectivity:}

\begin{longtable}[]{@{}lllll@{}}
\toprule\noalign{}
Element & Node 1 & Node 2 & Node 3 & Node 4 \\
\midrule\noalign{}
\endhead
\bottomrule\noalign{}
\endlastfoot
1 & 20 & 21 & 2 & 1 \\
2 & 21 & 22 & 3 & 2 \\
\end{longtable}

\textbf{Node coordinate table:}

\begin{longtable}[]{@{}lll@{}}
\toprule\noalign{}
Node & X-coord & Y-coord \\
\midrule\noalign{}
\endhead
\bottomrule\noalign{}
\endlastfoot
1 & X₁ & Y₁ \\
2 & X₂ & Y₂ \\
\ldots{} & \ldots{} & \ldots{} \\
\end{longtable}

\textbf{Importance:} Essential for assembling the global stiffness
matrix \([K_{ij}^g]\) and load vector \(\{F_i^g\}\), and for enforcing
boundary conditions.

\section{Shape Functions: Differential
Properties}\label{shape-functions-differential-properties}

\subsection{Coordinate Transformation
Relations}\label{coordinate-transformation-relations}

\textbf{Forward transformation (ζ → x):}

\[\frac{\partial}{\partial \zeta_1} = \frac{\partial}{\partial x}\frac{\partial x}{\partial \zeta_1} + \frac{\partial}{\partial y}\frac{\partial y}{\partial \zeta_1}\]

\[\frac{\partial}{\partial \zeta_2} = \frac{\partial}{\partial x}\frac{\partial x}{\partial \zeta_2} + \frac{\partial}{\partial y}\frac{\partial y}{\partial \zeta_2}\]

\textbf{Inverse transformation (x → ζ):}

\[\frac{\partial}{\partial x} = \frac{\partial}{\partial \zeta_1}\frac{\partial \zeta_1}{\partial x} + \frac{\partial}{\partial \zeta_2}\frac{\partial \zeta_2}{\partial x}\]

\[\frac{\partial}{\partial y} = \frac{\partial}{\partial \zeta_1}\frac{\partial \zeta_1}{\partial y} + \frac{\partial}{\partial \zeta_2}\frac{\partial \zeta_2}{\partial y}\]

\textbf{Matrix form:}

\[\begin{bmatrix} dx \\ dy \end{bmatrix} = \mathbf{F} \begin{bmatrix} d\zeta_1 \\ d\zeta_2 \end{bmatrix}, \quad \begin{bmatrix} d\zeta_1 \\ d\zeta_2 \end{bmatrix} = \mathbf{F}^{-1} \begin{bmatrix} dx \\ dy \end{bmatrix}\]

\subsection{Gradient Calculation in Master
Element}\label{gradient-calculation-in-master-element}

\textbf{For test function V:}

\[\text{grad}(V) = \begin{bmatrix} \frac{\partial V}{\partial x} \\ \frac{\partial V}{\partial y} \end{bmatrix} = \begin{bmatrix} \frac{\partial \zeta_1}{\partial x} & \frac{\partial \zeta_2}{\partial x} \\ \frac{\partial \zeta_1}{\partial y} & \frac{\partial \zeta_2}{\partial y} \end{bmatrix}^T \begin{bmatrix} \frac{\partial V}{\partial \zeta_1} \\ \frac{\partial V}{\partial \zeta_2} \end{bmatrix}\]

\textbf{For temperature field T:}

\[\text{grad}(T) = \begin{bmatrix} \frac{\partial T}{\partial x} \\ \frac{\partial T}{\partial y} \end{bmatrix} = \begin{bmatrix} \frac{\partial \zeta_1}{\partial x} & \frac{\partial \zeta_2}{\partial x} \\ \frac{\partial \zeta_1}{\partial y} & \frac{\partial \zeta_2}{\partial y} \end{bmatrix}^T \begin{bmatrix} \frac{\partial T}{\partial \zeta_1} \\ \frac{\partial T}{\partial \zeta_2} \end{bmatrix}\]

\chapter{The B Matrix}\label{the-b-matrix}

\section{Extension to 3D Vector
Problems}\label{extension-to-3d-vector-problems}

\subsection{Linear Elasticity Example}\label{linear-elasticity-example}

\textbf{Governing equation:}
\(\text{div}(\boldsymbol{\sigma}) + \mathbf{f} = 0\)

\textbf{Constitutive relation:}
\(\boldsymbol{\sigma} = \mathbf{E} \cdot \boldsymbol{\varepsilon} = \mathbf{E} \cdot \text{grad}^s(\mathbf{u})\)

\textbf{Divergence theorem in 3D:}

\[\int_\Omega \text{div}(\boldsymbol{\sigma}) \, d\Omega = \int_{\partial\Omega} \boldsymbol{\sigma} \cdot \mathbf{n} \, da\]

\textbf{Key extensions for 3D:}

\begin{itemize}
\tightlist
\item
  \textbf{Vector unknowns:} Displacement field
  \(\mathbf{u} = [u_x, u_y, u_z]^T\)
\item
  \textbf{Matrix operations:} Stress and strain tensors
\item
  \textbf{Global/local transformations:} More complex connectivity
\item
  \textbf{Vector problem connectivity:} Multiple DOFs per node
\end{itemize}

\section{3D Element Example}\label{d-element-example}

\begin{figure}[htbp]

{\centering \pandocbounded{\includegraphics[keepaspectratio]{images/hex20.png}}

}

\caption{20-node hexahedron element}

\end{figure}%%
\begin{figure}[htbp]

{\centering \pandocbounded{\includegraphics[keepaspectratio]{images/hex20r.png}}

}

\caption{20-node hexahedron element (refined)}

\end{figure}%%
\begin{figure}[htbp]

{\centering \pandocbounded{\includegraphics[keepaspectratio]{images/hex20f.png}}

}

\caption{20-node hexahedron element (field view)}

\end{figure}%

\section{The B Matrix: Relating Nodal Values to Field
Gradients}\label{the-b-matrix-relating-nodal-values-to-field-gradients}

In the Finite Element Method, we approximate a continuous field (like
temperature T or displacement u) within an element using shape functions
\(N_j\) and nodal values. For example, the temperature T(x,y) at any
point within an element can be approximated as:

\[T(x,y) \approx \sum_{j=1}^{n_{\text{nodes}}} N_j(x,y) a_j = \mathbf{N} \mathbf{a}^e\]

where \(N_j(x,y)\) are the shape functions, \(a_j\) are the nodal
temperatures for the element, \(\mathbf{N}\) is the row vector of shape
functions, and \(\mathbf{a}^e\) is the column vector of nodal
temperatures for the element.

The physical behavior of the system (e.g., heat flux, stress/strain)
often depends on the gradient of this field. The B matrix is a
cornerstone of FEM as it directly relates the nodal values of an element
to the gradient of the field (or its derivatives) within that element.

\section{For Scalar Field Problems (Heat
Conduction)}\label{for-scalar-field-problems-heat-conduction}

In heat conduction, the heat flux \(\mathbf{q}\) is related to the
temperature gradient \(\text{grad}(T)\) by Fourier's Law:
\(\mathbf{q} = -\mathbf{K} \cdot \text{grad}(T)\).

The weak form integral that leads to the stiffness matrix,
\(\int_{\Omega^e} \text{grad}(V) \cdot \mathbf{K} \cdot \text{grad}(T) \, d\Omega\),
requires evaluating these gradients.

Let \(T = \mathbf{N} \mathbf{a}^e\). The temperature gradient in 2D is:

\[\text{grad}(T) = \begin{Bmatrix} \frac{\partial T}{\partial x} \\ \frac{\partial T}{\partial y} \end{Bmatrix} = \begin{Bmatrix} \frac{\partial(\sum N_j a_j)}{\partial x} \\ \frac{\partial(\sum N_j a_j)}{\partial y} \end{Bmatrix} = \sum_{j=1}^{n_{\text{nodes}}} \begin{bmatrix} \frac{\partial N_j}{\partial x} \\ \frac{\partial N_j}{\partial y} \end{bmatrix} a_j\]

This can be written in matrix form as:

\[\text{grad}(T) = \mathbf{B} \mathbf{a}^e\]

Where the B matrix for a 2D scalar problem (like heat conduction) for an
element with \(n_{\text{nodes}}\) is:

\[\mathbf{B} = \begin{bmatrix}
\frac{\partial N_1}{\partial x} & \frac{\partial N_2}{\partial x} & \ldots & \frac{\partial N_{n_{\text{nodes}}}}{\partial x} \\
\frac{\partial N_1}{\partial y} & \frac{\partial N_2}{\partial y} & \ldots & \frac{\partial N_{n_{\text{nodes}}}}{\partial y}
\end{bmatrix}\]

Each column j of the B matrix corresponds to node j and contains the
partial derivatives of its shape function \(N_j\) with respect to the
global coordinates. For a 3D scalar problem, another row for
\(\frac{\partial N_j}{\partial z}\) would be added.

\section{Role in Element Stiffness
Matrix}\label{role-in-element-stiffness-matrix}

The element stiffness matrix for heat conduction, derived from the weak
form, is given by:

\[[\mathbf{k}^e] = \int_{\Omega^e} \mathbf{B}^T \mathbf{K} \mathbf{B} \, d\Omega\]

Here, \(\mathbf{K}\) is the material property matrix (thermal
conductivity tensor). For an isotropic material in 2D with scalar
conductivity k,
\(\mathbf{K} = \begin{bmatrix} k & 0 \\ 0 & k \end{bmatrix}\). The
integration is performed over the volume (or area in 2D) of the element
\(\Omega^e\).

\section{Calculation of B Matrix using Master Element
Coordinates}\label{calculation-of-b-matrix-using-master-element-coordinates}

Shape functions \(N_j\) are typically defined in local (master element)
coordinates \((\zeta_1, \zeta_2)\) (often denoted \(\xi, \eta\)). Their
derivatives with respect to global coordinates (x,y) are found using the
chain rule and the Jacobian matrix \(\mathbf{F}\) of the coordinate
transformation:

\[\begin{Bmatrix} \frac{\partial N_j}{\partial x} \\ \frac{\partial N_j}{\partial y} \end{Bmatrix} = \mathbf{F}^{-T} \begin{Bmatrix} \frac{\partial N_j}{\partial \zeta_1} \\ \frac{\partial N_j}{\partial \zeta_2} \end{Bmatrix}\]

where \(\mathbf{F}^{-T} = (\mathbf{F}^{-1})^T\).

The Jacobian matrix \(\mathbf{F}\) is defined as
\(\mathbf{F} = \begin{bmatrix} \frac{\partial x}{\partial \zeta_1} & \frac{\partial x}{\partial \zeta_2} \\ \frac{\partial y}{\partial \zeta_1} & \frac{\partial y}{\partial \zeta_2} \end{bmatrix}\).

The derivatives \(\frac{\partial N_j}{\partial \zeta_1}\) and
\(\frac{\partial N_j}{\partial \zeta_2}\) are easily computed from the
master element shape function definitions. The components of
\(\mathbf{F}^{-T}\) depend on the specific element geometry.

\section{For Vector Field Problems
(Elasticity)}\label{for-vector-field-problems-elasticity}

The concept of the B matrix is fundamental in structural mechanics
(elasticity), where it relates nodal displacements to strains within an
element.

The displacement field \(\mathbf{u}\) is approximated as
\(\mathbf{u} = \mathbf{N} \mathbf{d}^e\), where \(\mathbf{d}^e\) is the
vector of all nodal displacements for the element, and \(\mathbf{N}\) is
the matrix of shape functions arranged appropriately for a vector field.

The strain \(\boldsymbol{\varepsilon}\) is obtained by differentiating
the displacement field:
\(\boldsymbol{\varepsilon} = \mathbf{L} \mathbf{u}\), where
\(\mathbf{L}\) is a differential operator matrix. Combining these gives:

\[\boldsymbol{\varepsilon} = \mathbf{L}(\mathbf{N} \mathbf{d}^e) = (\mathbf{L} \mathbf{N}) \mathbf{d}^e = \mathbf{B} \mathbf{d}^e\]

The B matrix in elasticity thus contains derivatives of shape functions.
The element stiffness matrix is then computed as:

\[[\mathbf{k}^e] = \int_{\Omega^e} \mathbf{B}^T \mathbf{D} \mathbf{B} \, d\Omega\]

where \(\mathbf{D}\) is the material elasticity matrix (constitutive
matrix relating stress to strain).

\section{B Matrix for 2D Elasticity (Plane Stress / Plane
Strain)}\label{b-matrix-for-2d-elasticity-plane-stress-plane-strain}

For a 2D elasticity problem, the displacement at any point within an
element is
\(\mathbf{u}(x,y) = \begin{Bmatrix} u(x,y) \\ v(x,y) \end{Bmatrix}\).

The nodal displacement vector for a node i is
\(\mathbf{d}_i = \begin{Bmatrix} u_i \\ v_i \end{Bmatrix}\).

The displacement field is interpolated as:

\[\begin{Bmatrix} u \\ v \end{Bmatrix} = \sum_{i=1}^{n_{\text{nodes}}} \begin{bmatrix} N_i(x,y) & 0 \\ 0 & N_i(x,y) \end{bmatrix} \begin{Bmatrix} u_i \\ v_i \end{Bmatrix}\]

The strain vector in 2D is:

\[\boldsymbol{\varepsilon} = \begin{Bmatrix} \varepsilon_x \\ \varepsilon_y \\ \gamma_{xy} \end{Bmatrix} = \begin{Bmatrix} \frac{\partial u}{\partial x} \\ \frac{\partial v}{\partial y} \\ \frac{\partial u}{\partial y} + \frac{\partial v}{\partial x} \end{Bmatrix}\]

The B matrix contribution for node i is:

\[\mathbf{B}_i = \begin{bmatrix}
\frac{\partial N_i}{\partial x} & 0 \\
0 & \frac{\partial N_i}{\partial y} \\
\frac{\partial N_i}{\partial y} & \frac{\partial N_i}{\partial x}
\end{bmatrix}\]

\[\mathbf{B} = [\mathbf{B}_1 \quad \mathbf{B}_2 \quad \ldots \quad \mathbf{B}_{n_{\text{nodes}}}]\]

So, if an element has \(n_{\text{nodes}}\) nodes, the B matrix will have
3 rows (for \(\varepsilon_x, \varepsilon_y, \gamma_{xy}\)) and
\(2 \times n_{\text{nodes}}\) columns.

\section{B Matrix for 3D Elasticity}\label{b-matrix-for-3d-elasticity}

For a 3D elasticity problem, the displacement at any point is
\(\mathbf{u}(x,y,z) = \begin{Bmatrix} u(x,y,z) \\ v(x,y,z) \\ w(x,y,z) \end{Bmatrix}\).

The nodal displacement vector for node i is
\(\mathbf{d}_i = \begin{Bmatrix} u_i \\ v_i \\ w_i \end{Bmatrix}\).

The displacement field is interpolated as:

\[\begin{Bmatrix} u \\ v \\ w \end{Bmatrix} = \sum_{i=1}^{n_{\text{nodes}}} \begin{bmatrix} N_i(x,y,z) & 0 & 0 \\ 0 & N_i(x,y,z) & 0 \\ 0 & 0 & N_i(x,y,z) \end{bmatrix} \begin{Bmatrix} u_i \\ v_i \\ w_i \end{Bmatrix}\]

The strain vector (engineering strains) in 3D is:

\[\boldsymbol{\varepsilon} = \begin{Bmatrix} \varepsilon_x \\ \varepsilon_y \\ \varepsilon_z \\ \gamma_{xy} \\ \gamma_{yz} \\ \gamma_{zx} \end{Bmatrix} = \begin{Bmatrix}
\frac{\partial u}{\partial x} \\
\frac{\partial v}{\partial y} \\
\frac{\partial w}{\partial z} \\
\frac{\partial u}{\partial y} + \frac{\partial v}{\partial x} \\
\frac{\partial v}{\partial z} + \frac{\partial w}{\partial y} \\
\frac{\partial w}{\partial x} + \frac{\partial u}{\partial z}
\end{Bmatrix}\]

For each node, its contribution to the strain is related to the nodal
displacements of that node:

\[\mathbf{B}_i = \begin{bmatrix}
\frac{\partial N_i}{\partial x} & 0 & 0 \\
0 & \frac{\partial N_i}{\partial y} & 0 \\
0 & 0 & \frac{\partial N_i}{\partial z} \\
\frac{\partial N_i}{\partial y} & \frac{\partial N_i}{\partial x} & 0 \\
0 & \frac{\partial N_i}{\partial z} & \frac{\partial N_i}{\partial y} \\
\frac{\partial N_i}{\partial z} & 0 & \frac{\partial N_i}{\partial x}
\end{bmatrix}\]

The full element B matrix is assembled as
\(\mathbf{B} = [\mathbf{B}_1 \quad \mathbf{B}_2 \quad \ldots \quad \mathbf{B}_{n_{\text{nodes}}}]\).

If an element has \(n_{\text{nodes}}\) nodes, the B matrix will have 6
rows and \(3 \times n_{\text{nodes}}\) columns. The derivatives
\(\frac{\partial N_i}{\partial x}\),
\(\frac{\partial N_i}{\partial y}\), and
\(\frac{\partial N_i}{\partial z}\) are again found using the Jacobian
transformation from the master element coordinates.

\section{Key Takeaway (B Matrix)}\label{key-takeaway-b-matrix}

The B matrix is a crucial component in finite element analysis.

It contains the spatial derivatives of the element shape functions,
effectively translating the discrete nodal degrees of freedom into
continuous gradient fields (like temperature gradients) or deformation
measures (like mechanical strains) within an element.

This matrix is essential for forming the element stiffness matrix and,
consequently, for solving the overall system of equations that describes
the physical problem.

The components of B are generally functions of position within the
element (unless using simple elements like constant strain triangles),
and are evaluated numerically during the integration process to form the
element stiffness matrix.

\section{Shape Functions: Differential
Properties}\label{shape-functions-differential-properties-1}

To evaluate terms like \(\text{grad}(V)\) and \(\text{grad}(T)\) in the
weak form, we need the derivatives of shape functions (and consequently,
field variables) with respect to global coordinates (x,y).

However, shape functions \(N(\zeta_1, \zeta_2)\) are defined in the
master element coordinates \((\zeta_1, \zeta_2)\).

We use the chain rule and the Jacobian of the coordinate mapping to
transform these derivatives.

\subsection{Coordinate Transformation and the Jacobian
Matrix}\label{coordinate-transformation-and-the-jacobian-matrix}

Recall the isoparametric mapping from master coordinates
\((\zeta_1, \zeta_2)\) to global coordinates (x,y):

\[x = \sum_{i=1}^{n_{\text{nodes}}} N_i(\zeta_1, \zeta_2) x_i\]

\[y = \sum_{i=1}^{n_{\text{nodes}}} N_i(\zeta_1, \zeta_2) y_i\]

The Jacobian matrix of this transformation is:

\[\mathbf{F} = \begin{bmatrix} \frac{\partial x}{\partial \zeta_1} & \frac{\partial x}{\partial \zeta_2} \\ \frac{\partial y}{\partial \zeta_1} & \frac{\partial y}{\partial \zeta_2} \end{bmatrix}\]

The determinant of this matrix, \(J = \det(\mathbf{F})\), is used in
changing variables for integration:
\(dx \, dy = J \, d\zeta_1 \, d\zeta_2\).

Differentials are related by:

\[\begin{Bmatrix} dx \\ dy \end{Bmatrix} = \mathbf{F} \begin{Bmatrix} d\zeta_1 \\ d\zeta_2 \end{Bmatrix}\]

\[\begin{Bmatrix} d\zeta_1 \\ d\zeta_2 \end{Bmatrix} = \mathbf{F}^{-1} \begin{Bmatrix} dx \\ dy \end{Bmatrix}\]

\[\mathbf{F}^{-1} = \begin{bmatrix} \frac{\partial \zeta_1}{\partial x} & \frac{\partial \zeta_1}{\partial y} \\ \frac{\partial \zeta_2}{\partial x} & \frac{\partial \zeta_2}{\partial y} \end{bmatrix}\]

\subsection{Transformation of
Derivatives}\label{transformation-of-derivatives}

We need to express derivatives with respect to global coordinates (x,y)
in terms of derivatives with respect to master coordinates
\((\zeta_1, \zeta_2)\), since our shape functions \(N_i\) are given as
\(N_i(\zeta_1, \zeta_2)\).

Consider a function \(V(\zeta_1, \zeta_2)\). Its derivatives with
respect to \(\zeta_1\) and \(\zeta_2\) can be related to its derivatives
with respect to x and y using the chain rule:

\[\frac{\partial V}{\partial \zeta_1} = \frac{\partial V}{\partial x} \frac{\partial x}{\partial \zeta_1} + \frac{\partial V}{\partial y} \frac{\partial y}{\partial \zeta_1}\]

\[\frac{\partial V}{\partial \zeta_2} = \frac{\partial V}{\partial x} \frac{\partial x}{\partial \zeta_2} + \frac{\partial V}{\partial y} \frac{\partial y}{\partial \zeta_2}\]

\[\begin{Bmatrix} \frac{\partial V}{\partial \zeta_1} \\ \frac{\partial V}{\partial \zeta_2} \end{Bmatrix} = \begin{bmatrix} \frac{\partial x}{\partial \zeta_1} & \frac{\partial y}{\partial \zeta_1} \\ \frac{\partial x}{\partial \zeta_2} & \frac{\partial y}{\partial \zeta_2} \end{bmatrix} \begin{Bmatrix} \frac{\partial V}{\partial x} \\ \frac{\partial V}{\partial y} \end{Bmatrix} = \mathbf{F}^T \begin{Bmatrix} \frac{\partial V}{\partial x} \\ \frac{\partial V}{\partial y} \end{Bmatrix}\]

\[\begin{Bmatrix} \frac{\partial V}{\partial x} \\ \frac{\partial V}{\partial y} \end{Bmatrix} = (\mathbf{F}^T)^{-1} \begin{Bmatrix} \frac{\partial V}{\partial \zeta_1} \\ \frac{\partial V}{\partial \zeta_2} \end{Bmatrix} = \mathbf{F}^{-T} \begin{Bmatrix} \frac{\partial V}{\partial \zeta_1} \\ \frac{\partial V}{\partial \zeta_2} \end{Bmatrix}\]

\[\mathbf{F}^{-T} = \left(\begin{bmatrix} \frac{\partial x}{\partial \zeta_1} & \frac{\partial x}{\partial \zeta_2} \\ \frac{\partial y}{\partial \zeta_1} & \frac{\partial y}{\partial \zeta_2} \end{bmatrix}\right)^T = \begin{bmatrix} \frac{\partial x}{\partial \zeta_1} & \frac{\partial y}{\partial \zeta_1} \\ \frac{\partial x}{\partial \zeta_2} & \frac{\partial y}{\partial \zeta_2} \end{bmatrix}\]

\subsection{Gradient Calculation in Master Element
Coordinates}\label{gradient-calculation-in-master-element-coordinates}

Using the relationship derived above, the gradient of a test function V
(or temperature field T) in global coordinates is:

\[\text{grad}(V) = \begin{Bmatrix} \frac{\partial V}{\partial x} \\ \frac{\partial V}{\partial y} \end{Bmatrix} = \mathbf{F}^{-T} \begin{Bmatrix} \frac{\partial V}{\partial \zeta_1} \\ \frac{\partial V}{\partial \zeta_2} \end{Bmatrix} = \begin{bmatrix} \frac{\partial x}{\partial \zeta_1} & \frac{\partial y}{\partial \zeta_1} \\ \frac{\partial x}{\partial \zeta_2} & \frac{\partial y}{\partial \zeta_2} \end{bmatrix} \begin{Bmatrix} \frac{\partial V}{\partial \zeta_1} \\ \frac{\partial V}{\partial \zeta_2} \end{Bmatrix}\]

\[\frac{\partial V}{\partial x} = \frac{\partial x}{\partial \zeta_1} \frac{\partial V}{\partial \zeta_1} + \frac{\partial x}{\partial \zeta_2} \frac{\partial V}{\partial \zeta_2}\]

\[\frac{\partial V}{\partial y} = \frac{\partial y}{\partial \zeta_1} \frac{\partial V}{\partial \zeta_1} + \frac{\partial y}{\partial \zeta_2} \frac{\partial V}{\partial \zeta_2}\]

\subsection{Implementation Steps}\label{implementation-steps}

To implement this:

\begin{enumerate}
\def\labelenumi{\arabic{enumi}.}
\item
  For each shape function \(N_i(\zeta_1, \zeta_2)\), calculate its
  derivatives \(\frac{\partial N_i}{\partial \zeta_1}\) and
  \(\frac{\partial N_i}{\partial \zeta_2}\) analytically in the master
  element.
\item
  At each integration point \((\zeta_1, \zeta_2)\) within the master
  element:

  \begin{itemize}
  \tightlist
  \item
    Calculate the Jacobian matrix \(\mathbf{F}\) using the derivatives
    of the mapping functions \(x(\zeta_1, \zeta_2)\) and
    \(y(\zeta_1, \zeta_2)\).
  \item
    Compute \(\mathbf{F}^{-1}\) and then \(\mathbf{F}^{-T}\).
  \item
    Use \(\mathbf{F}^{-T}\) to transform the local derivatives
    \(\begin{Bmatrix} \frac{\partial N_i}{\partial \zeta_1} \\ \frac{\partial N_i}{\partial \zeta_2} \end{Bmatrix}\)
    into global derivatives
    \(\begin{Bmatrix} \frac{\partial N_i}{\partial x} \\ \frac{\partial N_i}{\partial y} \end{Bmatrix}\).
  \item
    These global derivatives form the columns of the B matrix for scalar
    problems, or parts of the columns for vector problems.
  \end{itemize}
\end{enumerate}

\section{Summary: Key Concepts}\label{summary-key-concepts}

\subsection{Mathematical Foundation}\label{mathematical-foundation}

\begin{itemize}
\tightlist
\item
  Product rule for divergence enables weak form derivation
\item
  Divergence theorem converts domain integrals to boundary integrals
\item
  Integration by parts transfers derivatives to test functions
\end{itemize}

\subsection{Geometric Considerations}\label{geometric-considerations}

\begin{itemize}
\tightlist
\item
  Master element concept simplifies integration
\item
  Jacobian must be positive for valid elements
\item
  Element quality depends on convexity and node ordering
\end{itemize}

\subsection{Implementation Details}\label{implementation-details}

\begin{itemize}
\tightlist
\item
  Shape functions provide local approximation
\item
  Connectivity tables manage global assembly
\item
  Coordinate transformations enable numerical integration
\end{itemize}

\subsection{Extension to Higher
Dimensions}\label{extension-to-higher-dimensions}

\begin{itemize}
\tightlist
\item
  Same mathematical principles apply
\item
  Increased complexity in bookkeeping and assembly
\item
  Vector problems require multiple DOFs per node
\end{itemize}

\chapter{2D Heat Transfer}\label{d-heat-transfer}

\section{The Strong Form}\label{the-strong-form-1}

The governing equation for 2D heat transfer is:

\[\text{div}(\mathbf{K} \cdot \text{grad}(T)) + Q = 0\]

with boundary conditions:

\[T=T_A \text{ on } \partial\Omega_A, \quad q=q^* \text{ on } \partial\Omega_B\]

where: - \(\mathbf{K}\) is the thermal conductivity tensor - \(Q\) is
the heat source term - \(q^*\) is a heat flux boundary condition

\section{The Weak Form}\label{the-weak-form-1}

To develop the weak form, we multiply the strong form by a test function
\(V\) and integrate over the domain:

\[\int_{\Omega} \text{grad}(V) \cdot \mathbf{K} \cdot \text{grad}(T) \, d\Omega = \int_{\Omega} Q V \, d\Omega - \int_{\partial\Omega_B}Vq^* \, da\]

\subsection{Test and Trial Functions}\label{test-and-trial-functions}

We choose to take both test and trial functions from the same function
space, which is spanned by the shape functions \({\phi}_i\) defined via
the master element coordinates \((\zeta_1, \zeta_2)\) and master element
shape functions \(\hat{\phi}\).

\[\begin{aligned}
T=\sum_{j=1}^N a_j \phi_j \\
V=\sum_{i=1}^N b_i \phi_i
\end{aligned}\]

Replacing the integral over sums with sum over integrals:

\[\sum_{i=1}^Nb_i\left[ \sum_{j=1}^N \left(\int_{\Omega}\text{grad}{\phi_i} \cdot \mathbf{K} \cdot \text{grad}{\phi_j} \, d\Omega\right) a_j \right]   -\sum_{i=1}^Nb_i\left(\int_{\Omega}\phi_i Q d\Omega -\int_{\partial \Omega_B}\phi_iq^*da\right)=0\]

\section{Element Stiffness Matrix}\label{element-stiffness-matrix}

The above equation represents a summation of \textbf{element} integrals
over \textbf{all nodes}.

For element \(e\) whose nodes are \(I,J\):

\[K_{IJ}^e = \int_{\Omega^e}\text{grad}{\phi_I} \cdot \mathbf{K} \cdot \text{grad}{\phi_J} \, d\Omega\]

Assuming isotropic material with constant thermal conductivity
\(\mathbf{K} = k \mathbf{I}\), we can simplify:

\[K_{IJ}^e = k \int_{\Omega^e}\text{grad}{\phi_I} \cdot \text{grad}{\phi_J} \, d\Omega\]

This leads to the element stiffness matrix:

\[[\mathbf{k}^e] = k \int_{\Omega^e} \mathbf{B}^T \mathbf{B} \, d\Omega\]

where \(\mathbf{B}\) is the B matrix containing the derivatives of the
shape functions \(\phi_i\) with respect to global coordinates (x,y):

\[[\mathbf{B}\phi^e] = \begin{bmatrix} \frac{\partial \phi_1}{\partial x_1} & \frac{\partial \phi_1}{\partial x_2} \\ \frac{\partial \phi_2}{\partial x_1} & \frac{\partial \phi_2}{\partial x_2} \\ \frac{\partial \phi_3}{\partial x_1} & \frac{\partial \phi_3}{\partial x_2} \\ \frac{\partial \phi_4}{\partial x_1} & \frac{\partial \phi_4}{\partial x_2} \end{bmatrix}\]

Such that:

\[K^e=\int_{\Omega^e}k[B\phi^e][B\phi^e]^T \, d\Omega\]

\section{Isoparametric Mapping and
Jacobian}\label{isoparametric-mapping-and-jacobian}

Using isoparametric mapping, we can express the shape functions in terms
of master element coordinates \((\zeta_1, \zeta_2)\) which are used for
interpolating the space as well:

\[\mathbf{x}(\zeta_1, \zeta_2) = \sum_{I=1}^{n_{\text{nodes}}} x_I\hat{\phi}_I(\zeta_1, \zeta_2)\]

The derivatives transform according to:

\[\frac{\partial\phi_I}{\partial x_j}= \sum_{k=1}^{dim}\frac{\partial\hat{\phi}_I}{\partial \zeta_k}\frac{\partial \zeta_k}{\partial x_j}\]

The term \([\mathbf{B}\hat{\phi}^e]_{Ij}\) can be expressed as:

\[\frac{\partial\hat{\phi}_I}{\partial x_j}=\sum_{k=1}^{dim}\frac{\partial{\phi}_I}{\partial x_j}\frac{\partial x_j}{\partial \zeta_k}\]

Or equivalently:

\[\mathbf{B}\hat{\phi^e}=\mathbf{B}\phi^e\mathbf{F} \\
\mathbf{B}{\phi}^e=\mathbf{B}{\hat{\phi^e}}\mathbf{F}^{-1}\]

With the Jacobian matrix defined as:

\[F_{ij}=\frac{\partial x_i}{\partial \zeta_j}=\sum_{I=1}^{nodes}x_I^j\frac{\partial\hat{\phi}_I(\mathbf{\zeta})}{\partial\zeta_j}\]

The determinant of \(\mathbf{F}\) (denoted \(J\)) is used to transform
between the area in the master element and the area in the global
coordinates:

\[d\Omega = J d\hat{\Omega}\]

where \(J = \det(\mathbf{F})\) and \(d\hat{\Omega}\) is the area in the
master element coordinates.

\section{Weak Form Terms in Master Element
Coordinates}\label{weak-form-terms-in-master-element-coordinates}

We can now write all the terms in the weak form at the element level:

\[\begin{aligned}
\mathbf{K}^e =\int_{\Omega^e} \mathbf{K}(\mathbf{x}(\mathbf{\zeta}))(\mathbf{B}\hat{\phi^e})(\mathbf{F}^T\mathbf{F})^{-1}  (\mathbf{B}\hat{\phi^e})^T \, Jd\hat{\Omega} \\
\mathbf{f}^e = \int_{\Omega^e} \hat{\phi^e}(\mathbf{\zeta})f(\mathbf{x}(\mathbf{\zeta})) \, Jd\hat{\Omega}
\end{aligned}\]

\section{Flux Boundary Condition in Master Element
Coordinates}\label{flux-boundary-condition-in-master-element-coordinates}

The only remaining task is to define the flux boundary condition
\(\int_{\partial\Omega_B}\phi_iq^*da\) in terms of the master element
coordinates.

We start with relating \(da\) (global) and \(dA\) (local):

\[d\mathbf{\zeta} = \frac{d\mathbf{\zeta}}{dA} \, dA=\mathbf{M}\, dA \\
d\mathbf{x} = \frac{d\mathbf{x}}{da} \, da \, = \mathbf{m}\, da\]

The differential area relationship is:

\[dx_i = \sum_{j=1}^{dim}\frac{\partial x_i}{\partial \zeta_j} d\zeta_j = \sum_{j=1}^{dim} F_{ij} d\zeta_j\Rightarrow d\mathbf{x}=\mathbf{F}d\mathbf{\zeta}\]

Using the norm properties:

\[d\mathbf{\zeta} \cdot d\mathbf{\zeta}=dA^2 \quad ; \quad d\mathbf{x} \cdot d\mathbf{x}=da^2\]

\[da^2=\mathbf{F}d\mathbf{\zeta} \cdot \mathbf{F}d\mathbf{\zeta} = dA^2(\mathbf{F}\mathbf{M})\cdot(\mathbf{F}\mathbf{M})\]

Therefore:

\[\frac{da}{dA}=\sqrt{(\mathbf{F}\mathbf{M})^T(\mathbf{F}\mathbf{M})}\]

The boundary integral in terms of master element coordinates becomes:

\[Q_i^s = -\int_{\partial\Omega_B}\phi_iq^*da =  -\int_{\partial\hat{\Omega}_B}\hat{\phi}_i(\mathbf{\zeta})q^*(\mathbf{\zeta}(\mathbf{x}))\sqrt{(\mathbf{F}\mathbf{M})^T(\mathbf{F}\mathbf{M})}dA\]

The flux boundary condition is applied on some line of the 2D element,
which means either \(\zeta_1\) or \(\zeta_2\) will be constant and equal
to \(\pm 1\).

\section{2D Finite Element Assembly
Algorithm}\label{d-finite-element-assembly-algorithm}

\begin{tcolorbox}[enhanced jigsaw, opacitybacktitle=0.6, leftrule=.75mm, colframe=quarto-callout-note-color-frame, toptitle=1mm, breakable, bottomtitle=1mm, opacityback=0, colback=white, rightrule=.15mm, colbacktitle=quarto-callout-note-color!10!white, arc=.35mm, left=2mm, toprule=.15mm, bottomrule=.15mm, titlerule=0mm, coltitle=black, title={Algorithm: 2D Finite Element Assembly}]

\textbf{Input:} Number of elements \(N_e\), Gauss points \(n_g\)

\textbf{Output:} Global stiffness matrix \(\mathbf{K}\) and force vector
\(\mathbf{F}\)

\textbf{Steps:}

for \(\text{elem} = 1\) to \(N_e\) do

\(\quad\) Set \(\mathbf{K}^e = \mathbf{0}\),
\(\mathbf{f}^e = \mathbf{0}\); Form \(\mathbf{X}^e\)

\(\quad\) for \(m = 1\) to \(n_g\) do (Gauss points)

\(\quad\quad\) for \(n = 1\) to \(n_g\) do

\(\quad\quad\quad\)
\(\boldsymbol{\zeta} = [\zeta_1, \zeta_2] = [\text{Points}[m], \text{Points}[n]]\)

\(\quad\quad\quad\) Evaluate
\(\mathbf{x}(\boldsymbol{\zeta}) = (\mathbf{X}^e)^T \hat{\boldsymbol{\phi}}(\boldsymbol{\zeta})\),
\(\mathbf{F}\), \(\det(\mathbf{F})\)

\(\quad\quad\quad\)
\(\mathbf{K}^e \leftarrow \mathbf{K}^e + w_m w_n K(\mathbf{x}(\boldsymbol{\zeta})) \mathbf{D}{\hat{\boldsymbol{\phi}}^e} (\mathbf{F}^T \mathbf{F})^{-1} (\mathbf{D}{\hat{\boldsymbol{\phi}}^e})^T \det(\mathbf{F})\)

\(\quad\quad\quad\)
\(\mathbf{f}^e \leftarrow \mathbf{f}^e + w_m w_n f(\mathbf{x}(\boldsymbol{\zeta})) \hat{\boldsymbol{\phi}}(\boldsymbol{\zeta}) \det(\mathbf{F})\)

\(\quad\quad\) end for

\(\quad\) end for

\(\quad\) Assemble into global matrix \([\mathbf{K}]{N \times N}\) and
vector \([\mathbf{F}]{N \times 1}\)

end for

\end{tcolorbox}

\part{Part IV: Elasticity}

\chapter{Continuum Mechanics
Fundamentals}\label{continuum-mechanics-fundamentals}

\section{Displacement and
Configuration}\label{displacement-and-configuration}

In continuum mechanics, we deal with continuous media and their
deformation under external forces.

\textbf{Displacement} \(\mathbf{u}\): The change in position of a point
in the body, defined as \(\mathbf{u} = \mathbf{x} - \mathbf{x}_0\),
where \(\mathbf{x}_0\) is the original position, often written using
capital \(\mathbf{X}\).

\textbf{Configuration:}

\begin{itemize}
\item
  \(\mathbf{X}\): The original position of a point in the body before
  deformation, with \(E_i\) being the coordinate system in the reference
  configuration: \(\mathbf{X}=X_iE_i\)
\item
  \(\mathbf{x}\): The current position of the same point after
  deformation, with \(e_i\) being the coordinate system in the current
  configuration: \(\mathbf{x}=x_ie_i\)
\end{itemize}

\section{Deformation Gradient}\label{deformation-gradient}

The \textbf{Deformation Gradient} \(\mathbf{F}\) is a tensor that
describes the local deformation of the material. It relates the current
configuration to the reference configuration:

\[\mathbf{F} = \frac{\partial \mathbf{x}}{\partial \mathbf{X}} = \begin{bmatrix} \frac{\partial x_1}{\partial X_1} & \frac{\partial x_1}{\partial X_2} & \ldots \\ \frac{\partial x_2}{\partial X_1} & \frac{\partial x_2}{\partial X_2} & \ldots \\ \vdots & \vdots & \ddots \end{bmatrix}\]

\section{Strain}\label{strain}

\textbf{Strain} \(\boldsymbol{\varepsilon}\) is a measure of deformation
defined as the symmetric part of the deformation gradient:

\[\boldsymbol{\varepsilon} = \frac{1}{2}(\mathbf{F}^T \mathbf{F} - \mathbf{I})\]

where \(\mathbf{I}\) is the identity tensor.

Under the assumption of small deformations, the strain can be
approximated as:

\[\boldsymbol{\varepsilon} \approx \frac{1}{2}(\nabla \mathbf{u} + \nabla \mathbf{u}^T) = \nabla^s\mathbf{u}\]

where \(\nabla \mathbf{u}\) is the gradient of the displacement field
and \(\nabla^s\mathbf{u}\) is the symmetric gradient.

\section{Stress and Constitutive
Relation}\label{stress-and-constitutive-relation}

\textbf{Stress} \(\boldsymbol{\sigma}\) is a measure of internal forces
within the material, defined using a constitutive relation that relates
stress to strain. For linear elasticity:

\[\boldsymbol{\sigma} = \mathbf{C} : \boldsymbol{\varepsilon}\]

where \(\mathbf{C}\) is the elasticity tensor (material property matrix)
and \(\boldsymbol{\sigma}\) is the Cauchy stress tensor.

\section{Tractions}\label{tractions}

\textbf{Tractions} \(\mathbf{t}\) represent the force per unit area
acting on a surface, related to stress by:

\[\mathbf{t} = \boldsymbol{\sigma} \cdot \mathbf{n}\]

where \(\mathbf{n}\) is the outward normal to the surface.

\section{Mass Conservation}\label{mass-conservation}

The principle that mass is conserved in a closed system. In continuum
mechanics, this is expressed as:

\[\int_{\Omega} \rho \, dv = \text{constant}= \int_{\Omega} \rho_0 \, dV\Rightarrow \rho J = \rho_0\]

where \(\rho\) is the current density, \(\rho_0\) is the reference
density, and \(J = \det(\mathbf{F})\) is the Jacobian determinant of the
deformation gradient.

\section{Balance of Linear Momentum}\label{balance-of-linear-momentum}

The balance of linear momentum states:

\[\int_{\Omega} \rho \frac{\partial^2 \mathbf{x}}{\partial t^2} \, dv = \int_{\Omega} \mathbf{f} \,dv +\int_{\partial \Omega} \mathbf{t} \, da\]

This leads to the equation of motion:

\[\text{div}(\boldsymbol{\sigma}) + \mathbf{f} = \rho \ddot{\mathbf{u}}\]

\section{Balance of Angular Momentum}\label{balance-of-angular-momentum}

The balance of angular momentum is expressed as:

\[\frac{d}{dt}\int_{\Omega} \mathbf{x} \times (\rho v)\, dv = \int_{\Omega} \mathbf{x} \times \mathbf{f} \,dv +\int_{\partial \Omega} \mathbf{x} \times \mathbf{t} \, da\]

\section{The Weak Form Derivation}\label{the-weak-form-derivation}

For quasi-static problems, we assume the acceleration terms are zero, so
the balance of linear momentum becomes:

\[\nabla \cdot \boldsymbol{\sigma} + \mathbf{f} = 0\]

with

\[\nabla \cdot \boldsymbol{\sigma} = \left(\frac{\partial \sigma_{ij}}{\partial x_j}\right)\mathbf{e}_i\]

where \(\sigma_{ij}\) are the components of the stress tensor.

\subsection{Multiplication by Test Function and
Integration}\label{multiplication-by-test-function-and-integration}

We multiply by a test function \(\mathbf{v}\) and integrate over the
domain:

\[\int_{\Omega} (\nabla \cdot \boldsymbol{\sigma} + \mathbf{f}) \cdot \mathbf{v} \, d\Omega = 0\]

\subsection{Product Rule Application}\label{product-rule-application}

We use the product rule for the divergence of a tensor contracted with a
vector:

\[\nabla \cdot (\boldsymbol{\sigma}\cdot\mathbf{v}) = (\nabla \cdot \boldsymbol{\sigma})\cdot \mathbf{v} + \boldsymbol{\sigma}:\nabla(\mathbf{v})\]

where:

\[\boldsymbol{\sigma}:\nabla(\mathbf{v}) = \sigma_{ij}\frac{\partial v_j}{\partial x_i}\]

Rearranging:

\[(\nabla \cdot \boldsymbol{\sigma}) \cdot \mathbf{v} = \nabla \cdot (\boldsymbol{\sigma}\mathbf{v}) - \boldsymbol{\sigma} : \nabla\mathbf{v}\]

\subsection{Divergence Theorem
Application}\label{divergence-theorem-application}

Applying the divergence theorem to the integrated equation:

\[\int_{\Omega} \nabla \cdot (\boldsymbol{\sigma}\mathbf{v}) d\Omega - \int_{\Omega} \boldsymbol{\sigma} : \nabla\mathbf{v} d\Omega + \int_{\Omega} \mathbf{f} \cdot \mathbf{v} d\Omega = 0\]

\[\implies \int_{\partial\Omega} (\boldsymbol{\sigma}\mathbf{v}) \cdot \mathbf{n} d\Gamma - \int_{\Omega} \boldsymbol{\sigma} : \nabla\mathbf{v} d\Omega + \int_{\Omega} \mathbf{f} \cdot \mathbf{v} d\Omega = 0\]

The term (\(\boldsymbol{\sigma}\boldsymbol{n}\)) is the traction vector
\(\boldsymbol{t}\). The term \(\nabla \boldsymbol{v}\) is the gradient
of the test function, and its symmetric part is the strain tensor for
the test function, \(\boldsymbol{\varepsilon}(\boldsymbol{v})\).

\subsection{Final Weak Form}\label{final-weak-form}

This derivation yields the weak form of the momentum equation, which
forms the foundation for finite element elasticity analysis.

\chapter{Quasi-Static Elasticity in
3D}\label{quasi-static-elasticity-in-3d}

\section{Strong Form and Method of Weighted
Residuals}\label{strong-form-and-method-of-weighted-residuals}

Starting from the equation of motion in the quasi-static case:

\[\nabla \cdot \boldsymbol{\sigma} + \mathbf{f} =  0\]

We multiply by a test function \(\mathbf{v}\) and integrate over the
domain \(\Omega\). This is the method of weighted residuals, where
\(\boldsymbol{r}\) is the residual:

\[\int_{\Omega} \left(\nabla\cdot \boldsymbol{\sigma}+\boldsymbol{f}\right)\cdot\mathbf{v} \, d\Omega = 0=\int_{\Omega}\boldsymbol{r}\cdot\boldsymbol{v}d\Omega\]

\pandocbounded{\includegraphics[keepaspectratio]{images/3dElast.png}}

\section{Product Rule and Integration by
Parts}\label{product-rule-and-integration-by-parts}

Recall the product rule for the divergence of a tensor contracted with a
vector:

\[\nabla \cdot (\boldsymbol{\sigma}\mathbf{v}) = (\nabla \cdot \boldsymbol{\sigma}) \cdot \mathbf{v} + \boldsymbol{\sigma} : \nabla\mathbf{v}\]

We can write:

\[(\nabla \cdot \boldsymbol{\sigma}) \cdot \mathbf{v} = \nabla \cdot (\boldsymbol{\sigma}\mathbf{v}) - \boldsymbol{\sigma} : \nabla\mathbf{v}\]

Substituting this into our integrated equation and applying the
divergence theorem:

\[\int_{\Omega} \nabla \cdot (\boldsymbol{\sigma}\mathbf{v}) d\Omega - \int_{\Omega} \boldsymbol{\sigma} : \nabla\mathbf{v} d\Omega + \int_{\Omega} \mathbf{f} \cdot \mathbf{v} d\Omega = 0\]

\[\implies \int_{\partial\Omega} (\boldsymbol{\sigma}\mathbf{v}) \cdot \mathbf{n} d\Gamma - \int_{\Omega} \boldsymbol{\sigma} : \nabla\mathbf{v} d\Omega + \int_{\Omega} \mathbf{f} \cdot \mathbf{v} d\Omega = 0\]

The term (\(\boldsymbol{\sigma}\boldsymbol{n}\)) is the traction vector
\(\boldsymbol{t}\). The term \(\nabla \boldsymbol{v}\) is the gradient
of the test function, and its symmetric part is the strain tensor for
the test function, \(\boldsymbol{\varepsilon}(\boldsymbol{v})\).

\section{Boundary Conditions and Natural Boundary
Terms}\label{boundary-conditions-and-natural-boundary-terms}

The boundary \(\partial \Omega\) is composed of two parts:

\begin{itemize}
\tightlist
\item
  \(\Gamma_u\) (where displacements are prescribed)
\item
  \(\Gamma_t\) (where tractions are prescribed)
\end{itemize}

On \(\Gamma_u\) we enforce Dirichlet boundary conditions, meaning the
displacement field \(\mathbf{u}\) is prescribed as \(\mathbf{u}^*\).

On \(\Gamma_t\) we enforce Neumann boundary conditions, meaning the
traction vector \(\boldsymbol{\sigma n}=\boldsymbol{t}\) is prescribed
as \(\boldsymbol{t}^*\).

This is the natural boundary condition that will be incorporated into
the weak form. The boundary integral becomes:

\[\int_{\partial\Omega} \mathbf{t} \cdot \mathbf{v} d\Gamma = \int_{\Gamma_t} \mathbf{t}^* \cdot \mathbf{v} d\Gamma + \int_{\Gamma_u} \mathbf{t} \cdot \underbrace{\mathbf{v}}_{=0} d\Gamma = \int_{\Gamma_t} \mathbf{t}^* \cdot \mathbf{v} d\Gamma\]

\section{Final Weak Form}\label{final-weak-form-1}

We can now state our problem in its final weak form:

Find the displacement field \(\boldsymbol{u} \in H^1(\Omega)\) such that
\(\boldsymbol{u} = \boldsymbol{u}^*\) on the Dirichlet boundary
\(\Gamma_u\) for all test functions \(\boldsymbol{v}\in H^1\) for which
\(\boldsymbol{v}|_{\Gamma_u}=0\) satisfying the weak form:

\[\int_{\Omega} \boldsymbol{\sigma}(\mathbf{u}) : \boldsymbol{\varepsilon}(\mathbf{v}) d\Omega = \int_{\Omega} \mathbf{f} \cdot \mathbf{v} d\Omega + \int_{\Gamma_t} \mathbf{t}^* \cdot \mathbf{v} d\Gamma\]

\section{Constitutive Law and Elasticity
Tensor}\label{constitutive-law-and-elasticity-tensor}

In the weak form, we have the stress tensor \(\boldsymbol{\sigma}\)
expressed in terms of the strain tensor \(\boldsymbol{\varepsilon}\) and
the displacement field \(\mathbf{u}\).

We use the constitutive law for linear elasticity
\(\boldsymbol{\sigma}=\boldsymbol{L}:\boldsymbol{\varepsilon}\), where
\(\boldsymbol{L}\) is the fourth-order elasticity tensor. For an
isotropic material, the matrix form of \(\boldsymbol{L}\) (in Voigt
notation) is given by:

\[\mathbf{L} = \frac{E}{(1+\nu)(1-2\nu)}
\begin{pmatrix}
1-\nu & \nu & \nu & 0 & 0 & 0 \\
\nu & 1-\nu & \nu & 0 & 0 & 0 \\
\nu & \nu & 1-\nu & 0 & 0 & 0 \\
0 & 0 & 0 & \frac{1-2\nu}{2} & 0 & 0 \\
0 & 0 & 0 & 0 & \frac{1-2\nu}{2} & 0 \\
0 & 0 & 0 & 0 & 0 & \frac{1-2\nu}{2}
\end{pmatrix}\]

where E is Young's modulus and ν is Poisson's ratio.

The term \(\boldsymbol{\sigma}=\boldsymbol{L}:\boldsymbol{\varepsilon}\)
refers to the double contraction of the elasticity tensor with the
strain tensor, which can be expressed in index notation as:

\[\sigma_{ij} = L_{ijkl} \varepsilon_{kl}\]

where \(\sigma_{ij}\) are the components of the stress tensor,
\(L_{ijkl}\) are the components of the elasticity tensor, and
\(\varepsilon_{kl}\) are the components of the strain tensor.

\section{Voigt Notation for Strain and
Stress}\label{voigt-notation-for-strain-and-stress}

In Voigt notation, the strain tensor is represented as a vector:

\[\boldsymbol{\varepsilon} = \begin{bmatrix}
\varepsilon_{11} \\ \varepsilon_{22} \\ \varepsilon_{33} \\ 2\varepsilon_{23} \\ 2\varepsilon_{31} \\ 2\varepsilon_{12} \end{bmatrix}\]

and the stress tensor is represented as:

\[\boldsymbol{\sigma} = \begin{bmatrix}
\sigma_{11} \\ \sigma_{22} \\ \sigma_{33} \\ \sigma_{23} \\ \sigma_{31} \\ \sigma_{12} \end{bmatrix}\]

\section{Bilinear and Linear Forms}\label{bilinear-and-linear-forms}

We continue from our weak form:

\[\int_{\Omega} \boldsymbol{\sigma}(\mathbf{u}) : \boldsymbol{\varepsilon}(\mathbf{v}) \, d\Omega = \int_{\Omega} \mathbf{f} \cdot \mathbf{v} \, d\Omega + \int_{\Gamma_t} \mathbf{t}^* \cdot \mathbf{v} \, d\Gamma\]

Now, we substitute the constitutive law into the left-hand side:

\[\int_{\Omega} (\boldsymbol{L}\cdot\boldsymbol{\varepsilon}(\mathbf{u})) \cdot \boldsymbol{\varepsilon}(\mathbf{v}) \, d\Omega = \int_{\Omega} \mathbf{f} \cdot \mathbf{v} \, d\Omega + \int_{\Gamma_t} \mathbf{t}^* \cdot \mathbf{v} \, d\Gamma\]

Where we dropped the double contraction (\(:\)) and replaced it with
\(\cdot\) since we use Voigt notation for the strain and stress tensors.

This equation has a standard structure. The left side is linear with
respect to both the solution u and the test function v. The right side
is linear only with respect to the test function v.

We can define a \textbf{bilinear form} \(B(u,v)\) from the left-hand
side, which represents the internal virtual work:

\[B(\mathbf{u}, \mathbf{v}) := \int_{\Omega} \boldsymbol{\varepsilon}(\mathbf{v}) \cdot \boldsymbol{L} \cdot \boldsymbol{\varepsilon}(\mathbf{u}) \, d\Omega\]

And we define a \textbf{linear form} \(F(v)\) from the right-hand side,
which represents the external virtual work done by body forces and
applied tractions:

\[F(\mathbf{v}) := \int_{\Omega} \mathbf{f} \cdot \mathbf{v} \, d\Omega + \int_{\Gamma_t} \mathbf{t}^* \cdot \mathbf{v} \, d\Gamma\]

The problem is now stated in its abstract form: Find the displacement u
(satisfying Dirichlet BCs) such that for all valid test functions v:

\[B(\mathbf{u}, \mathbf{v}) = F(\mathbf{v})\]

This abstract form is the direct starting point for the Finite Element
discretization, where we will replace the infinite-dimensional function
spaces with finite-dimensional approximations.

Recalling the definition of the strain tensor:

\[\boldsymbol{\varepsilon}(\mathbf{v}) = \frac{1}{2}(\nabla \mathbf{v} + \nabla \mathbf{v}^T)\]

We can express the bilinear form in terms of the gradient of the test
function:

\[B(\mathbf{u}, \mathbf{v}) = \int_{\Omega} \boldsymbol{\varepsilon}(\mathbf{v}) \cdot \boldsymbol{L} \cdot \boldsymbol{\varepsilon}(\mathbf{u}) \, d\Omega = \int_{\Omega} \left(\frac{1}{2}(\nabla \mathbf{v} + \nabla \mathbf{v}^T)\right) : \boldsymbol{L} : \left(\frac{1}{2}(\nabla \mathbf{u} + \nabla \mathbf{u}^T)\right) \, d\Omega\]

\section{Function Spaces for Test and Trial
Functions}\label{function-spaces-for-test-and-trial-functions}

To ensure our weak form is mathematically well-posed, the solution
(trial) function \(\boldsymbol{u}\) and the test function
\(\boldsymbol{v}\) can't be just any function. They must belong to
specific function spaces.

Why do we need restrictions? The bilinear form \textbf{B(u,v)} involves
integrals of derivatives of the functions. For these integrals to be
well-defined and finite, the functions must have a certain degree of
smoothness.

\subsection{Hilbert-Sobolev Spaces}\label{hilbert-sobolev-spaces}

A function u belongs to \(H^1\) if both the function and its first
derivatives are square-integrable over the domain Ω:

\[H^1(\Omega) = \left( \mathbf{u} \in L^2(\Omega) \; \middle| \; \nabla\mathbf{u} \in L^2(\Omega) \right)\]

In other words:

\[\int_{\Omega}\nabla \boldsymbol{u} \cdot\nabla \boldsymbol{u} \, d\Omega + \int_{\Omega} \boldsymbol{u} \cdot \boldsymbol{u} \, d\Omega < \infty\]

This ensures that the energy norm is finite.

\subsection{Handling Dirichlet
Conditions}\label{handling-dirichlet-conditions}

The test functions \(\boldsymbol{v}\) must also respect the homogeneous
Dirichlet boundary conditions. This leads us to the space
\(H^1_0(\Omega)\), which is a subspace of \(H^1(\Omega)\).

It contains all functions that are zero on the Dirichlet boundary
\(\Gamma_u\):

\[H^1_0(\Omega) = \left[ \mathbf{v} \in H^1(\Omega) \; \middle| \; \mathbf{v} = \mathbf{0} \text{ on } \Gamma_u \right]\]

\section{The Lifting Technique for Non-Zero
Displacements}\label{the-lifting-technique-for-non-zero-displacements}

If our prescribed displacement \(\boldsymbol{u}^*\) on the boundary
\(\Gamma_u\) is not zero, we must handle this carefully using the
\textbf{lifting} technique.

Mathematically, we handle this by splitting the solution u into two
parts:

\[\mathbf{u} = \mathbf{u}_0 + \mathbf{u}_D\]

where \(\boldsymbol{u}_D\) is a known function (aka ``lifting
function'') that satisfies the Dirichlet boundary condition on
\(\Gamma_u\), and \(\boldsymbol{u}_0\) is the unknown part of the
solution that we seek.

\subsection{Modified Weak Form}\label{modified-weak-form}

We substitute this decomposition into our weak form:

\[B(\mathbf{u}_0 + \mathbf{u}_D, \mathbf{v}) = F(\mathbf{v})\]

Using the linearity of B(⋅,v):

\[B(\mathbf{u}_0, \mathbf{v}) + B(\mathbf{u}_D, \mathbf{v}) = F(\mathbf{v})\]

We then move all the known quantities to the right-hand side:

\[B(\mathbf{u}_0, \mathbf{v}) = F(\mathbf{v}) - B(\mathbf{u}_D, \mathbf{v})\]

Our new problem is to find the homogeneous component
\(\boldsymbol{u}_0 \in H^1_0(\Omega)\) that satisfies this modified
equation for all test functions \(\boldsymbol{v} \in H^1_0(\Omega)\):

\[B(\mathbf{u}_0, \mathbf{v}) = F(\mathbf{v}) - B(\mathbf{u}_D, \mathbf{v})\]

This means we are looking for a solution \(\boldsymbol{u}_0\) that
satisfies the weak form with respect to the homogeneous space
\(H^1_0(\Omega)\), while \(\boldsymbol{u}_D\) is already known and
satisfies the Dirichlet boundary condition.

Once \(\boldsymbol{u}_0\) is found, the final solution is simply
\(\boldsymbol{u} = \boldsymbol{u}_0 + \boldsymbol{u}_D\). In practice,
this is handled during the assembly of the finite element system.

\section{Principle of Minimum Potential
Energy}\label{principle-of-minimum-potential-energy}

An alternative but equivalent way to derive the weak form for elasticity
problems is by minimizing a potential energy functional. This principle
states that the exact solution to an elastic problem is the one that
minimizes the total potential energy.

\subsection{Total Potential Energy}\label{total-potential-energy}

The total potential energy of the system is the sum of the internal
strain energy and the potential energy of the external loads:

\[\Pi(\mathbf{u}) = \underbrace{\frac{1}{2} \int_{\Omega} \boldsymbol{\sigma}(\mathbf{u}) : \boldsymbol{\varepsilon}(\mathbf{u}) \, d\Omega}_{\text{Strain Energy}} - \underbrace{\left( \int{\Omega} \mathbf{f} \cdot \mathbf{u} \, d\Omega + \int_{\Gamma_t} \bar{\mathbf{t}} \cdot \mathbf{u} \, d\Gamma \right)}_{\text{Potential of External Loads}}\]

\subsection{Finding the Minimum}\label{finding-the-minimum}

The solution u that minimizes the functional \(\Pi\) is found by setting
its directional derivative to zero for any arbitrary perturbation (test
function) v. This is analogous to setting the first derivative to zero
in standard calculus:

\[\lim_{\epsilon \to 0} \frac{\Pi(\mathbf{u} + \epsilon\mathbf{v}) - \Pi(\mathbf{u})}{\epsilon} = 0\]

Calculating this derivative (which is the first variation of the
functional) yields:

\[\delta\Pi = \int_{\Omega} \boldsymbol{\sigma}(\mathbf{u}) : \boldsymbol{\varepsilon}(\mathbf{v}) \, d\Omega - \int_{\Omega} \mathbf{f} \cdot \mathbf{v} \, d\Omega - \int_{\Gamma_t} \bar{\mathbf{t}} \cdot \mathbf{v} \, d\Gamma = 0\]

Rearranging this gives us our familiar weak form:

\[\int_{\Omega} \boldsymbol{\sigma}(\mathbf{u}) : \boldsymbol{\varepsilon}(\mathbf{v}) \, d\Omega = \int_{\Omega} \mathbf{f} \cdot \mathbf{v} \, d\Omega + \int_{\Gamma_t} \bar{\mathbf{t}} \cdot \mathbf{v} \, d\Gamma\]

This shows that the solution satisfying the weak form is also the one
that minimizes the total potential energy of the system.

\section{Proof that Weak Form Solution Minimizes
Energy}\label{proof-that-weak-form-solution-minimizes-energy}

We previously showed that minimizing the potential energy functional
leads to the weak form. Now we will prove the converse: the unique
solution u to the weak form is the unique minimizer of the potential
energy functional \(L(\mathbf{w})\).

Define the Energy Functional:

\[L(\mathbf{w}) \overset{\text{def}}{=} \frac{1}{2} B(\mathbf{w}, \mathbf{w}) - F(\mathbf{w})\]

The Goal: We want to show that \(J(u)\leq J(w)\) for any other
admissible function \(\boldsymbol{w}\) (a function that satisfies the
same Dirichlet boundary conditions as \(\boldsymbol{u}\)).

Let's evaluate the energy of any admissible function w: We can write any
such w as the sum of the true solution u and an error function e, so
w=u+e.

Since u and w both satisfy the same Dirichlet boundary conditions, the
error function e must be zero on the Dirichlet boundary \(\Gamma_u\) and
thus belongs to the space \(H^1_0(\Omega)\), which is the space of
functions that are square-integrable and have square-integrable first
derivatives, and are zero on the Dirichlet boundary.

Let's substitute w=u+e into the energy functional:

\[L(\mathbf{w}) = L(\mathbf{u}+\mathbf{e}) = \frac{1}{2}B(\mathbf{u}+\mathbf{e}, \mathbf{u}+\mathbf{e}) - F(\mathbf{u}+\mathbf{e})\]

Using the linearity of the forms, this expands to:

\[L(\mathbf{w}) = \frac{1}{2}\left( B(\mathbf{u},\mathbf{u}) + 2B(\mathbf{u},\mathbf{e}) + B(\mathbf{e},\mathbf{e}) \right) - \left( F(\mathbf{u}) + F(\mathbf{e}) \right)\]

Now, we rearrange the expression:

\[L(\mathbf{w}) = \underbrace{\left( \frac{1}{2}B(\mathbf{u},\mathbf{u}) - F(\mathbf{u}) \right)}_{L(\mathbf{u})} + \underbrace{\left( B(\mathbf{u},\mathbf{e}) - F(\mathbf{e}) \right)}_{=0} + \underbrace{\frac{1}{2}B(\mathbf{e},\mathbf{e})}_{\frac{1}{2}|\mathbf{e}|_E^2}\]

The middle term is zero because u is the true solution, so it satisfies
the weak form \(B(u,v)=F(v)\) for any test function v, including our
error function e.

The Final Result: This simplifies to a fundamental relationship:

\[L(\mathbf{w}) = L(\mathbf{u}) + \frac{1}{2} | \mathbf{w} - \mathbf{u} |_E^2\]

Since the energy norm \(||w-u||_E^2\) is always greater than or equal to
zero, this proves that \(L(w)\geq L(u)\).

The energy of any admissible function is greater than or equal to the
energy of the true solution. The minimum is achieved only when
\(||w-u||_E^2=0\), which means \(w=u\).

\section{FEM Approximation and Error
Estimates}\label{fem-approximation-and-error-estimates}

The Finite Element Method does not find the exact solution, but rather
an approximate solution \(u_h\). \textbf{The accuracy of this
approximation is a central concern.}

Discretization: We divide our domain Ω into smaller elements of
characteristic size h. We then approximate the solution using simple
polynomial functions (e.g., linear, quadratic) on each element.

Error Measurement: The error is the difference between the exact
solution \(u\) and the finite element solution \(u_h\). We often measure
this error using the energy norm:

\[|| \mathbf{e} ||^2_E = || \mathbf{u} - \mathbf{u}_h ||^2_E = B(\mathbf{u} - \mathbf{u}_h, \mathbf{u} - \mathbf{u}_h)\]

This norm is physically meaningful as it represents the strain energy of
the error.

\subsection{A Priori Error Estimate}\label{a-priori-error-estimate}

For well-posed problems, we can estimate the error before solving. The
error in the energy norm is bounded by a constant \(C\) (independent of
the mesh) times terms involving the element size h and the polynomial
order p of the shape functions:

\[|| \mathbf{u} - \mathbf{u}_h ||^2_E \le C^2(u,p) h^{2\gamma}\]

Where \(\gamma=\text{min}(r-1,p)\), with r being the regularity of the
solution and p being the polynomial order of the finite element shape
functions.

Key Takeaway: The error decreases as the mesh is refined (as h gets
smaller) or as the polynomial order of the elements (p) is increased. We
can compare the solutions for two different meshes or polynomial orders
to estimate the convergence rate of the method.

\part{Part V: Implementation \& Error Analysis}

\chapter{FEM Implementation}\label{fem-implementation}

\section{Discretize the Domain}\label{discretize-the-domain}

The first step in any Finite Element analysis is to move from the
continuous domain Ω to a discrete one. We replace the single integral
over the whole domain with a sum of integrals over small, simple
subdomains called elements.

Our weak form is:

\[\underbrace{\int_{\Omega} \boldsymbol{\varepsilon}(\mathbf{v})^T \mathbf{L} \boldsymbol{\varepsilon}(\mathbf{u}) \, d\Omega}_{\text{Bilinear Form } B(\mathbf{u},\mathbf{v})} = \underbrace{\int{\Omega} \mathbf{v}^T \mathbf{f} \, d\Omega + \int_{\Gamma_t} \mathbf{v}^T \bar{\mathbf{t}} \, d\Gamma}_{\text{Linear Form } F(\mathbf{v})}\]

Note: We've switched to transpose notation for Voigt vectors, common in
implementation.

We approximate the domain Ω as a mesh of \(N_{el}\) finite elements,
each with its own subdomain \(\Omega_e\).

The integral over Ω becomes a sum of integrals over each element
\(\Omega_e\):

The bilinear and linear forms are now expressed as sums:

\[B(\mathbf{u},\mathbf{v}) = \sum_{e=1}^{N_{el}} \int_{\Omega_e} \boldsymbol{\varepsilon}(\mathbf{v})^T \mathbf{L} \boldsymbol{\varepsilon}(\mathbf{u}) \, d\Omega_e\]

\[F(\mathbf{v}) = \sum_{e=1}^{N_{el}} \left( \int_{\Omega_e} \mathbf{v}^T \mathbf{f} \, d\Omega_e + \int_{\Gamma_{t,e}} \mathbf{v}^T \bar{\mathbf{t}} \, d\Gamma_e \right)\]

where \(\Gamma_{t,e}\) is the part of the element's boundary that lies
on the traction boundary \(\Gamma_t\).

\section{Approximate Fields with Shape
Functions}\label{approximate-fields-with-shape-functions}

Within each element, we approximate the continuous displacement field
using simple polynomial functions (shape functions) and the displacement
values at the element's nodes.

The displacement u(x) at any point inside an element is interpolated
from the element's nodal displacements \(d_e\) using shape functions:

\[\mathbf{u}(\mathbf{x}) \approx \mathbf{u}_h(\mathbf{x}) = \mathbf{N}(\mathbf{x}) \mathbf{d}_e\]

N(x) is the shape function matrix. For a 3D problem, it relates the 3
displacement components to the nodal values. For an element with n
nodes, it's a 3×3n matrix:

\[\mathbf{N} = \begin{bmatrix} \phi_1 & 0 & 0 & \phi_2 & 0 & 0 & \dots & \phi_n & 0 & 0 \\ 0 & \phi_1 & 0 & 0 & \phi_2 & 0 & \dots & 0 & \phi_n & 0 \\ 0 & 0 & \phi_1 & 0 & 0 & \phi_2 & \dots & 0 & 0 & \phi_n \end{bmatrix}\]

where \(\phi_i(x)\) is the scalar shape function for node i.

We do exactly the same for the test function v, relating it to a vector
of arbitrary nodal values \(c_e\):

\[\mathbf{v}(\mathbf{x}) \approx \mathbf{v}_h(\mathbf{x}) = \mathbf{N}(\mathbf{x}) \mathbf{c}_e\]

\section{Compute Strains with the
B-Matrix}\label{compute-strains-with-the-b-matrix}

With the displacement field approximated, we can find the strain by
taking its derivatives. This process creates the crucial
strain-displacement matrix, B.

The strain vector \(\boldsymbol{\varepsilon}\) is found by applying a
differential operator matrix, ∂, to the displacement vector
\(\boldsymbol{u}\):

\[\boldsymbol{\varepsilon} = \partial \mathbf{u} \quad \text{where for 3D engineering strain, } \quad \partial = \begin{pmatrix} \partial_x & 0 & 0 \\ 0 & \partial_y & 0 \\ 0 & 0 & \partial_z \\ 0 & \partial_z & \partial_y \\ \partial_z & 0 & \partial_x \\ \partial_y & \partial_x & 0 \end{pmatrix}\]

For the approximated displacement field:

\[\boldsymbol{\varepsilon}_h = \partial \left( \mathbf{N}\mathbf{d}_e \right) = \left( \partial\mathbf{N} \right) \mathbf{d}_e\]

We define the strain-displacement matrix B as the derivatives of the
shape function matrix:

\[\mathbf{B} := \partial\mathbf{N}\]

This gives the direct algebraic relationship between strain and nodal
displacements for an element:

\[\boldsymbol{\varepsilon}_h = \mathbf{B} \mathbf{d}_e\]

The B matrix contains the spatial derivatives of the shape functions and
is a function of the position within the element.

\section{The B-Matrix from the Symmetric
Gradient}\label{the-b-matrix-from-the-symmetric-gradient}

The crucial B-matrix is not arbitrary; it arises directly from the
fundamental definition of the small strain tensor, ε, as the symmetric
part of the displacement gradient.

\subsection{Start with the Strain
Definition}\label{start-with-the-strain-definition}

\[\boldsymbol{\varepsilon} = \frac{1}{2} \left( \nabla\mathbf{u} + (\nabla\mathbf{u})^T \right)\]

\subsection{Form the Voigt Strain
Vector}\label{form-the-voigt-strain-vector}

For implementation, we use Voigt notation with engineering shear strains
\(\gamma_{ij}\), where
\(\gamma_{ij} = \varepsilon_{ij} + \varepsilon_{ji}\), so
\(\gamma_{ij} = 2\varepsilon_{ij}\) for symmetric strains:

\[\boldsymbol{\varepsilon}_{\text{Voigt}} = \begin{pmatrix} \varepsilon_{xx} \\ \varepsilon_{yy} \\ \varepsilon_{zz} \\ \gamma_{yz} \\ \gamma_{xz} \\ \gamma_{xy} \end{pmatrix} = \begin{pmatrix} \partial u / \partial x \\ \partial v / \partial y \\ \partial w / \partial z \\ (\partial v / \partial z) + (\partial w / \partial y) \\ (\partial u / \partial z) + (\partial w / \partial x) \\ (\partial u / \partial y) + (\partial v / \partial x) \end{pmatrix}\]

\subsection{Define the Differential
Operator}\label{define-the-differential-operator}

We can express this relationship using a differential operator matrix,
∂, that acts on the displacement vector:

\[\boldsymbol{\varepsilon}_{\text{Voigt}} = \underbrace{ \begin{pmatrix} \partial_x & 0 & 0 \\ 0 & \partial_y & 0 \\ 0 & 0 & \partial_z \\ 0 & \partial_z & \partial_y \\ \partial_z & 0 & \partial_x \\ \partial_y & \partial_x & 0 \end{pmatrix} }_{\partial} \underbrace{ \begin{pmatrix} u \\ v \\ w \end{pmatrix} }_{\mathbf{u}}\]

This explicitly shows how the structure of ∂ is determined by the
physics of strain.

\subsection{Introduce the FE
Approximation}\label{introduce-the-fe-approximation}

Finally, substitute the approximation for the displacement field,
\(u_h=Nd_e\), into the strain definition:

\[\boldsymbol{\varepsilon}_h = \partial (\mathbf{N}\mathbf{d}_e) = (\partial\mathbf{N})\mathbf{d}_e\]

This gives us our definition of the B-matrix, \(B:=\partial N\), and the
final element-level relationship:

\[\boldsymbol{\varepsilon}_h = \mathbf{B}\mathbf{d}_e\]

\section{\texorpdfstring{The Algebraic System:
\(k_e d_e = f_e\)}{The Algebraic System: k\_e d\_e = f\_e}}\label{the-algebraic-system-k_e-d_e-f_e}

By substituting the shape function approximations for u and v into the
element integral, we transform the differential equation into a system
of linear algebraic equations.

Consider the bilinear form for one element, and substitute
\(\boldsymbol{\varepsilon}(u)=\boldsymbol{B}\boldsymbol{d}_e\) and
\(\boldsymbol{\varepsilon}(v)=\boldsymbol{B}\boldsymbol{c}_e\):

\[ \int_{\Omega_e} (\mathbf{B}\mathbf{c}_e)^T \mathbf{L} (\mathbf{B}\mathbf{d}_e) \, d\Omega_e = \mathbf{c}_e^T \left( \int{\Omega_e} \mathbf{B}^T \mathbf{L} \mathbf{B} \, d\Omega_e \right) \mathbf{d}_e \]

We define the term in the parentheses as the \textbf{element stiffness
matrix} \(k_e\):

\[ \mathbf{k}_e := \int{\Omega_e} \mathbf{B}^T \mathbf{L} \mathbf{B} \, d\Omega_e \]

Similarly, for the linear form, we get the \textbf{element force vector}
\(f_e\):

\[ \mathbf{f}_e := \int{\Omega_e} \mathbf{N}^T \mathbf{f} \, d\Omega_e + \int_{\Gamma_{t,e}} \mathbf{N}^T \bar{\mathbf{t}} \, d\Gamma_e \]

The weak form for the element must hold for any arbitrary \(c_e\), which
can only be true if:

\[ \mathbf{k}_e \mathbf{d}_e = \mathbf{f}_e \]

\section{Numerical Integration}\label{numerical-integration}

The integrals for \(k_e\) and \(f_e\) involve the B and N matrices,
which can be complex polynomials. These are almost always computed
numerically using \textbf{Gaussian Quadrature}.

The idea is to approximate an integral by a weighted sum of the
integrand's value at specific sample points, known as Gauss points.

First, we map the element \(\Omega_e\) to a simple reference element.
The integral transforms as:

\[\int_{\Omega_e} g(\mathbf{x}) \, d\Omega_e = \int_{\hat{\Omega}} g(\mathbf{x}(\boldsymbol{\xi})) \, |\mathbf{J}(\boldsymbol{\xi})| \, d\hat{\Omega}\]

where ∣J∣ is the determinant of the Jacobian of the coordinate mapping.

The integral over the reference element is then approximated by the sum:

\[\int_{\hat{\Omega}} \tilde{g}(\boldsymbol{\xi}) \, d\hat{\Omega} \approx \sum_{i=1}^{N_{gp}} w_i \tilde{g}(\boldsymbol{\xi}_i)\]

where \(w_i\) are the Gauss weights and \(\boldsymbol{\xi}_i\) are the
Gauss point coordinates.

\subsection{Application to Element
Stiffness}\label{application-to-element-stiffness}

Applying this to our stiffness matrix:

\[\mathbf{k}_e = \int{\Omega_e} \mathbf{B}^T \mathbf{L} \mathbf{B} \, d\Omega \approx \sum_{i=1}^{N_{gp}} \left( \mathbf{B}(\boldsymbol{\xi}_i)^T \mathbf{L} \mathbf{B}(\boldsymbol{\xi}_i) \right) |\mathbf{J}(\boldsymbol{\xi}_i)| w_i\]

The force vector \(f_e\) is computed in the same way:

\[\mathbf{f}_e = \int{\Omega_e} \mathbf{N}^T \mathbf{f} \, d\Omega_e + \int_{\Gamma_{t,e}} \mathbf{N}^T \mathbf{t}^* \, d\Gamma_e \approx \sum_{i=1}^{N_{gp}} \left( \mathbf{N}(\boldsymbol{\xi}_i)^T \mathbf{f}(\boldsymbol{\xi}_i) + \mathbf{N}(\boldsymbol{\xi}_i)^T \mathbf{t}^*(\boldsymbol{\xi}_i) \right) |\mathbf{J}(\boldsymbol{\xi}_i)| w_i\]

\section{Algorithmic Overview}\label{algorithmic-overview}

\subsection{Step-by-step Procedure}\label{step-by-step-procedure}

\begin{enumerate}
\def\labelenumi{\arabic{enumi}.}
\item
  Initialize: Create empty global stiffness matrix K and force vector F.
\item
  Loop over Elements: For each element e in the mesh\ldots{}
\item
  Initialize Element Matrices: Create empty local stiffness matrix
  \(k_e\) and force vector \(f_e\).
\item
  Loop over Gauss Points: For each Gauss point i in the element\ldots{}
\item
  Calculate at Point: Evaluate the B matrix and the Jacobian determinant
  ∣J∣ at this point's coordinates.
\item
  Accumulate Integral: Compute the integrand for \(k_e\) and \(f_e\).
  Multiply by the Gauss weight and ∣J∣ and add to the element matrices.
\item
  Assemble: Once the loops over Gauss points are done, add the completed
  element matrices \(k_e\) and \(f_e\) into the global \(K\) and \(F\)
  based on the element's node numbering.
\item
  Apply Boundary Conditions: Modify the global system \(K d = F\) to
  enforce the prescribed Dirichlet displacements.
\item
  Solve: Solve the final system of linear equations for the vector of
  all nodal displacements d.
\item
  Post-Process: Use the computed displacements d to find element strains
  and stresses using the same matrices B and L as before at the Gauss
  points.
\end{enumerate}

\section{Full FEM Algorithm}\label{full-fem-algorithm}

\begin{tcolorbox}[enhanced jigsaw, opacitybacktitle=0.6, leftrule=.75mm, colframe=quarto-callout-note-color-frame, toptitle=1mm, breakable, bottomtitle=1mm, opacityback=0, colback=white, rightrule=.15mm, colbacktitle=quarto-callout-note-color!10!white, arc=.35mm, left=2mm, toprule=.15mm, bottomrule=.15mm, titlerule=0mm, coltitle=black, title={Algorithm: Finite Element Procedure for 3D Elasticity}]

\textbf{Input:} Mesh (\({\Omega_e}\)), Material Properties
(\(\mathbf{L}\)), Body Forces (\(\mathbf{f}\)), Boundary Conditions
(\(\bar{\mathbf{u}}, \bar{\mathbf{t}}\))

\textbf{Output:} Nodal Displacements (\(\mathbf{d}\)), Element Stresses
(\(\boldsymbol{\sigma}\)), Element Strains
(\(\boldsymbol{\varepsilon}\))

\textbf{Steps:}

Initialize Global Stiffness \(\mathbf{K} \leftarrow \mathbf{0}\) and
Force \(\mathbf{F} \leftarrow \mathbf{0}\).

For each element e in the mesh:

\begin{enumerate}
\def\labelenumi{\alph{enumi}.}
\item
  Initialize Element Stiffness \(\mathbf{k}_e \leftarrow \mathbf{0}\)
  and Force \(\mathbf{f}_e \leftarrow \mathbf{0}\).
\item
  For each Gauss point i in element e:

  \begin{itemize}
  \item
    Compute \(\mathbf{B}(\boldsymbol{\xi}_i)\), Jacobian determinant
    \(|\mathbf{J}(\boldsymbol{\xi}_i)|\), and weight \(w_i\).
  \item
    Accumulate stiffness:
    \(\mathbf{k}_e \mathrel{+}= \mathbf{B}^T \mathbf{L} \mathbf{B} \, |\mathbf{J}| w_i\)
  \item
    Accumulate body forces:
    \(\mathbf{f}_e \mathrel{+}= \mathbf{N}^T \mathbf{f} \, |\mathbf{J}| w_i\)
  \end{itemize}
\item
  Assemble traction boundary forces into \(\mathbf{f}_e\) if applicable.
\item
  Assemble local into global:
  Assemble\((\mathbf{K}, \mathbf{F}, \mathbf{k}_e, \mathbf{f}_e)\).
\end{enumerate}

Apply Dirichlet boundary conditions to the global system
\((\mathbf{K}, \mathbf{F})\).

Solve the linear system \(\mathbf{K}\mathbf{d} = \mathbf{F}\) for the
displacement vector \(\mathbf{d}\).

For each element e in the mesh (Post-processing):

\begin{enumerate}
\def\labelenumi{\alph{enumi}.}
\item
  Extract element displacements \(d_e\) from global d.
\item
  Compute strains \(\boldsymbol{\varepsilon}_e = \mathbf{B} d_e\) and
  stresses
  \(\boldsymbol{\sigma}_e = \mathbf{L} \boldsymbol{\varepsilon}_e\) at
  desired points.
\end{enumerate}

\end{tcolorbox}

\section{Surface Integrals and Nanson's
Formula}\label{surface-integrals-and-nansons-formula}

To compute integrals over an element's surface in the real domain, such
as for traction boundary conditions, we must map them back to the simple
reference domain. This requires a transformation for the surface area
element itself.

Consider the surface integral over the traction boundary \(\Gamma_t\) of
an element:

\[\int_{\Gamma_{t,e}} \bar{\mathbf{t}} \, d\Gamma_e\]

The integration domain \(\Gamma_{t,e}\) is a 2D surface of the 3D
element \(\Gamma_e\). The surface domain \(\Gamma_{t,e}\) is potentially
distorted in real space. Mapping it to the reference element will allow
us to perform the surface integral over a planar well-defined surface
(e.g.~a triangle or quadrilateral).

The mapping between the reference (\(\boldsymbol{\xi}\)) and real space
(\(\boldsymbol{x}\)) is done using the coordinate mapping and its
gradient (Jacobian):

\[|\boldsymbol{F}| = |\frac{\partial \boldsymbol{x}}{\partial \boldsymbol{\xi}}| = \boldsymbol{J}\]

The traction \(\bar{\mathbf{t}}\) is defined by the stress and the
vector normal to the surface \(\mathbf{n}\):

\[\bar{\mathbf{t}} = \boldsymbol{\sigma} \cdot \mathbf{n}\]

\subsection{Nanson's Formula}\label{nansons-formula}

\textbf{Nanson's Formula} provides the transformation for surface
integrals:

\[\boldsymbol{n}d\Gamma_e = |\boldsymbol{J}|\boldsymbol{J}^{-T}\cdot \boldsymbol{N} d\hat{\Gamma}\]

where \(\boldsymbol{N}\) is the normal vector in the reference element,
and \(d\hat{\Gamma}=d\xi_i d\xi_j\) is the differential area in the
reference element:

\[\mathbf{n} d\Gamma = \det(\mathbf{J}) (\mathbf{J}^{-T}) \hat{\mathbf{n}}  d\hat{\Gamma}\]

The surface Jacobian is defined as:

\[J_s =  \det(\mathbf{J}) (\mathbf{J}^{-T}) \hat{\mathbf{n}}\]

The surface integral becomes:

\[\int_{\Gamma_e} g(\mathbf{x})  d\Gamma = \int_{\hat{\Gamma}} g(\mathbf{x}(\boldsymbol{\xi})) J_s(\boldsymbol{\xi})  d\hat{\Gamma}\]

\chapter{Error Estimation and
Adaptivity}\label{error-estimation-and-adaptivity}

\section{Simple Error Estimates for Vector
Problems}\label{simple-error-estimates-for-vector-problems}

After computing a solution \(\mathbf{u}_h\), we need to assess its
quality.

\textbf{A posteriori error estimators} use the solution itself to
estimate the error in each element, guiding us on where the mesh needs
improvement.

The core idea is to check for violations of physical principles. For
elasticity, the exact stress tensor is continuous across element
boundaries. Our FE stress field, \(\boldsymbol{\sigma}_h\), is usually
discontinuous. The size of this \textbf{stress jump} can be used as an
indicator of local error.

For a face \(\gamma\) between elements \(K^+\) and \(K^-\), the jump is:

\[\mathbf{j}_{\gamma} = \boldsymbol{\sigma}_h^+ \mathbf{n} - \boldsymbol{\sigma}_h^- \mathbf{n}\]

The local error indicator \(\eta_K\) for element \(K\) sums the jumps
over its faces:

\[\eta_K^2 = \frac{1}{2} \sum_{\gamma \subset \partial K} \int_{\gamma} ||\mathbf{j}_{\gamma}||^2 \, d\Gamma\]

The total estimated error is then \(\eta^2 = \sum_K \eta_K^2\). This
gives us a computable error value for each element.

\section{Local Mesh Refinement
(h-Adaptivity)}\label{local-mesh-refinement-h-adaptivity}

The power of a posteriori error estimates lies in driving
\textbf{adaptive mesh refinement}.

Instead of uniformly refining the entire mesh (which is inefficient), we
add elements only where the error is large. This is a simple and
powerful feedback loop.

\subsection{The Adaptive Algorithm}\label{the-adaptive-algorithm}

\textbf{1. Solve:} Compute the FE solution \(\mathbf{u}_h\) on the
current mesh.

\textbf{2. Estimate:} For every element \(K\), compute the local error
indicator \(\eta_K\).

\textbf{3. Mark:} Identify elements for refinement. A common strategy is
to mark all elements \(K\) where the error is a significant fraction of
the maximum:

\[\text{Mark element } K \text{ if } \eta_K > \theta \cdot \eta_{\max} \quad (\text{e.g., } \theta = 0.5)\]

\textbf{4. Refine:} Subdivide each marked element into smaller ones
(\textbf{h-refinement}).

\textbf{5. Repeat:} Loop back to Step 1 and solve on the new, locally
denser mesh.

\pandocbounded{\includegraphics[keepaspectratio]{images/mref.png}}

\section{A Posteriori Recovery Methods: The Core
Idea}\label{a-posteriori-recovery-methods-the-core-idea}

Recovery methods work on a simple principle:

The raw FE stress field (\(\boldsymbol{\sigma}_h\)) is ``noisy'' and
inaccurate, but we can recover a better, smoother stress field
(\(\boldsymbol{\sigma}^*\)) from it. The difference between the two is
an excellent measure of the error.

\textbf{Why recover?} FE stresses are calculated from derivatives, which
amplifies error and makes the field discontinuous. The recovered field
\(\boldsymbol{\sigma}^*\) is continuous and much closer to the true
stress \(\boldsymbol{\sigma}\).

\subsection{Error Estimation Formula}\label{error-estimation-formula}

We estimate the error by assuming the recovered solution is the
``truth.'' The error in energy for an element \(K\) is the energy of the
difference:

\[\eta_K^2 = \int_{K} (\boldsymbol{\sigma}^* - \boldsymbol{\sigma}_h)^T \mathbf{L}^{-1} (\boldsymbol{\sigma}^* - \boldsymbol{\sigma}_h) \, dK\]

The main challenge is finding a good way to ``recover''
\(\boldsymbol{\sigma}^*\).

\section{The Zienkiewicz-Zhu (ZZ)
Estimator}\label{the-zienkiewicz-zhu-zz-estimator}

The \textbf{Zienkiewicz-Zhu (ZZ) estimator} is a classic recovery
technique that builds a smooth stress field using special, highly
accurate points within elements.

\subsection{Superconvergent Points}\label{superconvergent-points}

For many elements, there are specific locations---called
\textbf{superconvergent points}---where the calculated stresses are much
more accurate than anywhere else. For common elements, the \textbf{Gauss
points} used for integration (or some of them) are superconvergent.

\subsection{The ZZ Method}\label{the-zz-method}

The method constructs a global, continuous stress field
\(\boldsymbol{\sigma}^*\) that best fits the high-quality stress values
found at these superconvergent points.

\begin{enumerate}
\def\labelenumi{\arabic{enumi}.}
\tightlist
\item
  The recovered field is approximated using the \emph{same shape
  functions} as the displacement field:
\end{enumerate}

\[\boldsymbol{\sigma}^*(\mathbf{x}) = \mathbf{N}(\mathbf{x}) \boldsymbol{\sigma}_{\text{nodal}}^*\]

Here, \(\boldsymbol{\sigma}_{\text{nodal}}^*\) is a vector of unknown,
continuous stress values at the mesh nodes.

\begin{enumerate}
\def\labelenumi{\arabic{enumi}.}
\setcounter{enumi}{1}
\tightlist
\item
  These nodal stresses are found by a global \textbf{least-squares fit}
  to the superconvergent point data. We find
  \(\boldsymbol{\sigma}_{\text{nodal}}^*\) that minimizes:
\end{enumerate}

\[\sum_{\text{elements}} \sum_{\text{Gauss pts}} || \mathbf{N}(\boldsymbol{\xi}_i)\boldsymbol{\sigma}_{\text{nodal}}^* - \boldsymbol{\sigma}_h(\boldsymbol{\xi}_i) ||^2\]

\section{Superconvergent Patch Recovery
(SPR)}\label{superconvergent-patch-recovery-spr}

Solving a global least-squares problem can be expensive. A more
practical variant is \textbf{Superconvergent Patch Recovery (SPR)},
which works locally on patches of elements.

The SPR algorithm finds the recovered stress at a \textbf{single node}
p:

\subsection{5-Step SPR Algorithm}\label{step-spr-algorithm}

\begin{enumerate}
\def\labelenumi{\arabic{enumi}.}
\item
  \textbf{Define Patch:} Identify all elements connected to node p.
\item
  \textbf{Assume Polynomial:} Assume a simple polynomial for the
  recovered stress over the patch:
\end{enumerate}

\[\boldsymbol{\sigma}^*(\mathbf{x}) = \mathbf{P}(\mathbf{x}) \mathbf{a}\]

\begin{enumerate}
\def\labelenumi{\arabic{enumi}.}
\setcounter{enumi}{2}
\item
  \textbf{Sample Points:} Collect the raw FE stress values
  \(\boldsymbol{\sigma}_h\) from all superconvergent (Gauss) points
  within the patch.
\item
  \textbf{Local Fit:} Solve a small least-squares problem to find the
  polynomial coefficients \(\mathbf{a}\) that best fit the sampled data.
\item
  \textbf{Evaluate at Node:} Use the fitted polynomial to calculate the
  recovered stress value directly at the node's location:
  \(\boldsymbol{\sigma}^*_p = \mathbf{P}(\mathbf{x}_p)\mathbf{a}\).
\end{enumerate}

This is repeated for all nodes to define the continuous field
\(\boldsymbol{\sigma}^*\), which is then used to compute the error
indicators \(\eta_K\).

\section{Comparing Error Estimators}\label{comparing-error-estimators}

A direct comparison highlights the different strategies for evaluating
FE error. One method checks for violations of physical laws, while the
other tries to construct a better solution.

\subsection{Stress Jump (Residual)
Estimator}\label{stress-jump-residual-estimator}

\textbf{Core Philosophy:} ``How much does my solution violate the
physical laws?''

\textbf{What it Measures:} Violation of equilibrium at element
interfaces.

\textbf{Physical Intuition:} Unbalanced forces at element boundaries.

\textbf{Mathematical Basis:} Rigorous. Provides a proven upper bound on
the error.

\textbf{Pros:} - Mathematically sound - Conceptually linked to the PDE

\textbf{Cons:} - Can sometimes underestimate the error - Doesn't provide
a ``better'' solution to look at

\subsection{Zienkiewicz-Zhu (Recovery)
Estimator}\label{zienkiewicz-zhu-recovery-estimator}

\textbf{Core Philosophy:} ``Can I construct a better solution from the
one I have?''

\textbf{What it Measures:} The difference between the raw FE stress and
a ``cleaner'' field.

\textbf{Physical Intuition:} Filtering a noisy signal to get a smooth
one.

\textbf{Mathematical Basis:} Heuristic. Relies on the phenomenon of
superconvergence.

\textbf{Pros:} - Often very accurate - Provides a useful, smooth stress
field \(\sigma^*\) for visualization

\textbf{Cons:} - More complex to implement

\subsection{Summary}\label{summary-1}

The Jump Method is like a ``referee'' checking if the rules are
followed. The ZZ Method is like an ``artist'' creating a better picture
from a blurry photo.

\section{A Posteriori Residual Methods: The
Concept}\label{a-posteriori-residual-methods-the-concept}

Residual-based methods offer a more mathematically rigorous approach.
They don't need a ``recovered'' solution. Instead, they directly measure
how poorly the FE solution satisfies the original governing differential
equation.

\subsection{The Residual}\label{the-residual}

The residual is what's left over when you plug the approximate solution
\(\mathbf{u}_h\) back into the governing PDE. If \(\mathbf{u}_h\) were
the exact solution, the residual would be zero everywhere:

\[\mathbf{r} = \nabla \cdot \boldsymbol{\sigma}_h + \mathbf{f} \neq \mathbf{0}\]

\subsection{Key Idea}\label{key-idea}

These methods produce a proven \textbf{upper bound} on the solution
error. This means we can guarantee the true error is no larger than our
estimate (times a constant):

\[\| \mathbf{u} - \mathbf{u}_h \|_E \le \eta_{\text{residual}}\]

The final estimator is a sum of terms measuring the residual inside
elements, the jumps between elements, and the error in satisfying
boundary conditions.

\section{Breaking Down the Residual Error
Estimate}\label{breaking-down-the-residual-error-estimate}

The total residual error is a sum of local contributions that capture
all the ways our solution fails to be exact. The squared error is
bounded by a sum of these terms:

\[\| \mathbf{u} - \mathbf{u}_h \|_E^2 \le C_1 \sum_{K} h_K^2 \| \mathbf{r}_K \|^2_{L^2(K)} + C_2 \sum_{\gamma} h_{\gamma} \| \mathbf{j}_{\gamma} \|^2_{L^2(\gamma)} +C_3 \sum_{\gamma_t} h_{\gamma} \| \mathbf{t}_{\gamma_t} \|^2_{L^2(\gamma_t)}\]

\subsection{Three Components of the
Residual}\label{three-components-of-the-residual}

\textbf{1. Interior Residual (\(\mathbf{r}_K\)):} Measures how well the
equilibrium equation is satisfied \emph{inside} each element \(K\):

\[\mathbf{r}_K = \nabla \cdot \boldsymbol{\sigma}_h + \mathbf{f}\]

\textbf{2. Interface Residual (\(\mathbf{j}_{\gamma}\)):} This is the
familiar \textbf{stress jump} across interior element faces, \(\gamma\).
It measures the violation of force balance between elements:

\[\mathbf{j}_{\gamma} = \boldsymbol{\sigma}_h^+ \mathbf{n} - \boldsymbol{\sigma}_h^- \mathbf{n}\]

\textbf{3. Boundary Residual (\(\mathbf{t}_{\gamma_t}\)):} Measures how
well the solution satisfies the prescribed traction \(\bar{\mathbf{t}}\)
on Neumann boundaries \(\gamma_t\):

\[\mathbf{t}_{\gamma_t} = \bar{\mathbf{t}} - \boldsymbol{\sigma}_h \mathbf{n}\]

\subsection{Local Error Bound}\label{local-error-bound}

The local error at an element is given by:

\[C_1  h_K^2 \| \mathbf{r}_K \|^2_{L^2(K)} + C_2 h_{\gamma} \| \mathbf{j}_{\gamma} \|^2_{L^2(\gamma)} +C_3 h_{\gamma} \| \mathbf{t}_{\gamma_t} \|^2_{L^2(\gamma_t)}\]

This decomposition allows us to identify which elements have the largest
contributions to the total error, guiding adaptive refinement
strategies.


\backmatter


\end{document}
